% =========================================================================
% DRAFT: Uses \documentclass{article} for local
% compilation.  For POPL 2027 submission switch to:
%   \documentclass[acmsmall,review,anonymous]{acmart}
% with the ACM CCS/keywords block below uncommented.
%
% ACM metadata for final submission:
%   \acmJournal{PACMPL} \acmVolume{0} \acmNumber{POPL}
%   \acmYear{2027} \acmMonth{1} \setcopyright{none}
%
% CCS:
%   Theory of computation~Proof theory        [500]
%   Theory of computation~Logic and verification [300]
%   Theory of computation~Type theory         [300]
%
% Keywords: flag algebras, Lean 4, proof by reflection, semidefinite
%   programming, Turan density, interactive theorem proving,
%   tactic metaprogramming, extremal combinatorics
% =========================================================================
\documentclass{article}

\usepackage{hyperref}
\usepackage{amsmath}
\usepackage{mathtools}
\usepackage{amssymb}
\usepackage{amsthm}
\usepackage{dsfont}
\usepackage{stmaryrd}
\usepackage{xcolor}
\usepackage{listings}
\usepackage{booktabs}
\usepackage{microtype}
\usepackage{fullpage}
\usepackage{authblk}
\usepackage{aliascnt}
\usepackage{tikz}

% ---- Theorem environments -----------------------------------------------
\newtheorem{theorem}{Theorem}[section]

\newaliascnt{lemma}{theorem}
\newtheorem{lemma}[lemma]{Lemma}
\aliascntresetthe{lemma}

\newaliascnt{corollary}{theorem}
\newtheorem{corollary}[corollary]{Corollary}
\aliascntresetthe{corollary}

\newaliascnt{definition}{theorem}
\newtheorem{definition}[definition]{Definition}
\aliascntresetthe{definition}

\newaliascnt{example}{theorem}
\newtheorem{example}[example]{Example}
\aliascntresetthe{example}

\newaliascnt{remark}{theorem}
\newtheorem{remark}[remark]{Remark}
\aliascntresetthe{remark}

\usepackage{cleveref}
\crefname{theorem}{Theorem}{Theorems}
\Crefname{theorem}{Theorem}{Theorems}
\crefname{lemma}{Lemma}{Lemmas}
\Crefname{lemma}{Lemma}{Lemmas}
\crefname{corollary}{Corollary}{Corollaries}
\Crefname{corollary}{Corollary}{Corollaries}
\crefname{definition}{Definition}{Definitions}
\Crefname{definition}{Definition}{Definitions}
\crefname{example}{Example}{Examples}
\Crefname{example}{Example}{Examples}
\crefname{remark}{Remark}{Remarks}
\Crefname{remark}{Remark}{Remarks}

% ---- Macros -------------------------------------------------------------
\newcommand{\R}{\mathbb{R}}
\newcommand{\Q}{\mathbb{Q}}
\newcommand{\N}{\mathbb{N}}
\newcommand{\Fin}[1]{\mathsf{Fin}(#1)}
\newcommand{\Flag}[1]{\mathcal{F}^{#1}}
\newcommand{\FlagAlg}[1]{\mathcal{A}^{#1}}
\newcommand{\poshom}[1]{\mathrm{Hom}^+(\mathcal{A}^{#1}, \mathbb{R})}
\newcommand{\den}[2]{p(#1;\,#2)}
\newcommand{\tdensity}[2]{\pi(#1;\,#2)}
\newcommand{\lean}[1]{\texttt{#1}}
\newcommand{\emptyt}{\emptyset_{\mathrm{t}}}
\newcommand{\leansmul}{\mathbin{\vcenter{\hbox{$\scriptstyle\bullet$}}}}
\newcommand{\gw}[1]{\textcolor{red}{\textbf{[GW: #1]}}}
\newcommand{\sh}[1]{\textcolor{red}{\textbf{[SH: #1]}}}
\newcommand{\Ext}{\operatorname{Ext}}

\tikzset{
  graph vertex/.style={circle, draw, fill=white, inner sep=0pt, minimum size=4.5pt},
  graph root/.style={circle, draw, fill=black, inner sep=0pt, minimum size=4.5pt},
  graph edge/.style={line width=0.45pt},
  graph nonedge/.style={line width=0.45pt, dashed},
  inline/.style={scale=0.20, baseline=-0.5ex},
  fv/.style={circle, fill=white, draw, inner sep=0pt, minimum size=3pt},
  fr/.style={circle, fill=black, inner sep=0pt, minimum size=3pt}
}

% Inline flag macros: fv = unlabeled (open white), fr = labeled (filled black)
% 2-vertex flags
\newcommand{\fKtwo}{%
  \begin{tikzpicture}[inline]
    \draw (-1,0) -- (1,0);
    \node[fv] at (-1,0) {};
    \node[fv] at (1,0) {};
  \end{tikzpicture}%
}
\newcommand{\fKtwobull}{%
  \begin{tikzpicture}[inline]
    \draw (-1,0) -- (1,0);
    \node[fr] at (-1,0) {};
    \node[fv] at (1,0) {};
  \end{tikzpicture}%
}
\newcommand{\fbarKtwobull}{%
  \begin{tikzpicture}[inline]
    \node[fr] at (-1,0) {};
    \node[fv] at (1,0) {};
  \end{tikzpicture}%
}
% 3-vertex flags: (90:1)=top, (210:1)=bottom-left, (330:1)=bottom-right
\newcommand{\fbarKthree}{%
  \begin{tikzpicture}[inline]
    \node[fv] at (90:1) {};
    \node[fv] at (210:1) {};
    \node[fv] at (330:1) {};
  \end{tikzpicture}%
}
\newcommand{\fbarPthree}{%
  \begin{tikzpicture}[inline]
    \draw (210:1) -- (330:1);
    \node[fv] at (90:1) {};
    \node[fv] at (210:1) {};
    \node[fv] at (330:1) {};
  \end{tikzpicture}%
}
\newcommand{\fPthree}{%
  \begin{tikzpicture}[inline]
    \draw (210:1) -- (90:1) -- (330:1);
    \node[fv] at (90:1) {};
    \node[fv] at (210:1) {};
    \node[fv] at (330:1) {};
  \end{tikzpicture}%
}
\newcommand{\fKthree}{%
  \begin{tikzpicture}[inline]
    \draw (210:1) -- (90:1) -- (330:1) -- (210:1);
    \node[fv] at (90:1) {};
    \node[fv] at (210:1) {};
    \node[fv] at (330:1) {};
  \end{tikzpicture}%
}
% Labeled 3-vertex flags: labeled vertex (fr) at (90:1)
% label degree 0, no uu edge (all isolated)
\newcommand{\fbarKthreebull}{%
  \begin{tikzpicture}[inline]
    \node[fr] at (90:1) {};
    \node[fv] at (210:1) {};
    \node[fv] at (330:1) {};
  \end{tikzpicture}%
}
% label degree 0, uu edge between unlabeled vertices
\newcommand{\fEthreebull}{%
  \begin{tikzpicture}[inline]
    \draw (210:1) -- (330:1);
    \node[fr] at (90:1) {};
    \node[fv] at (210:1) {};
    \node[fv] at (330:1) {};
  \end{tikzpicture}%
}
\newcommand{\fbarPthreebull}{%
  \begin{tikzpicture}[inline]
    \draw (90:1) -- (330:1);
    \node[fr] at (90:1) {};
    \node[fv] at (210:1) {};
    \node[fv] at (330:1) {};
  \end{tikzpicture}%
}
\newcommand{\fPthreebull}{%
  \begin{tikzpicture}[inline]
    \draw (90:1) -- (330:1) -- (210:1);
    \node[fr] at (90:1) {};
    \node[fv] at (210:1) {};
    \node[fv] at (330:1) {};
  \end{tikzpicture}%
}
% Labeled 3-vertex flag: label at top connected to both (star/cherry)
\newcommand{\fKonetwobull}{%
  \begin{tikzpicture}[inline]
    \draw (90:1) -- (210:1);
    \draw (90:1) -- (330:1);
    \node[fr] at (90:1) {};
    \node[fv] at (210:1) {};
    \node[fv] at (330:1) {};
  \end{tikzpicture}%
}
% 5-vertex flags: pentagon layout
\newcommand{\fCfive}{%
  \begin{tikzpicture}[inline]
    \draw (90:1) -- (18:1) -- (-54:1) -- (-126:1) -- (162:1) -- cycle;
    \node[fv] at (90:1) {};
    \node[fv] at (18:1) {};
    \node[fv] at (-54:1) {};
    \node[fv] at (-126:1) {};
    \node[fv] at (162:1) {};
  \end{tikzpicture}%
}
\newcommand{\fCfivebull}{%
  \begin{tikzpicture}[inline]
    \draw (90:1) -- (18:1) -- (-54:1) -- (-126:1) -- (162:1) -- cycle;
    \node[fr] at (90:1) {};
    \node[fv] at (18:1) {};
    \node[fv] at (-54:1) {};
    \node[fv] at (-126:1) {};
    \node[fv] at (162:1) {};
  \end{tikzpicture}%
}

% Lean code style
\lstdefinelanguage{Lean4}{
  keywords={def,theorem,lemma,instance,structure,class,import,open,namespace,
             end,where,fun,let,have,show,by,match,with,if,then,else,
             return,do,for,in,noncomputable,abbrev,variable,section,
             native_decide,decide,norm_num,simp,ring,linarith,omega,
             exact,apply,rw,calc,intro,constructor,use,refine,obtain,
             rcases,cases,induction,assumption,contradiction,trivial,rfl,
             macro,elab,notation},
  keywordstyle=\color{blue!70!black}\bfseries,
  comment=[l]{--},
  commentstyle=\color{gray}\itshape,
  stringstyle=\color{orange!80!black},
  morestring=[b]",
  sensitive=true,
  basicstyle=\ttfamily\small,
  breaklines=true,
  columns=fullflexible,
  escapeinside={(*@}{@*)},
  literate=
    {->}{{$\to$}}2
    {<-}{{$\leftarrow$}}2
    {<=}{{$\leq$}}1
    {>=}{{$\geq$}}1
    {alpha}{{$\alpha$}}1
    {forall}{{$\forall$}}1
    {exists}{{$\exists$}}1
}
\lstset{language=Lean4, frame=single, framesep=4pt,
        rulecolor=\color{black!20}, backgroundcolor=\color{gray!5}}

% =========================================================================
\title{Formalizing Flag Algebras in Lean: Certified Proofs in Extremal Graph Theory and Root-Planting Semantics}
\date{}

\author[1,3]{Gyeongwon Jeong}
\author[1,3]{Seonghun Park}
\author[1]{Jihoon Hyun}
\author[2]{Sang-il Oum}
\author[3]{Hongseok Yang}
\affil[1]{School of Computing, KAIST, South Korea}
\affil[2]{Discrete Mathematics Group, Institute for Basic Science (IBS), South Korea}
\affil[3]{School of Computational Sciences, Korea Institute for Advanced Study (KIAS), South Korea}

% =========================================================================
\begin{document}
\maketitle

\begin{abstract}
Razborov's flag algebras are a powerful method for proving asymptotic inequalities in
extremal graph theory, often by reducing a graph-theoretic problem to a finite
semidefinite-programming certificate.  The method is especially useful in
constrained settings, where the desired inequality is required only for graphs
that avoid specified finite patterns, such as triangles.  We present a
machine-checked formalization of flag algebras for finite simple graphs,
together with a certificate-to-proof pipeline that reconstructs, in Lean,
algebraic proofs suggested by external certificate data.

The formalization covers the mathematical layer of the method: partially
labeled graphs, their densities in large graphs, the quotient algebra of
density expressions, graph-limit semantics through positive homomorphisms, and
the downward operators used to average out labels.  The external search is
used only to suggest finite data: lists of flags, density and multiplication
tables, and rational positive-semidefinite matrices.  Lean then checks these
data against computable graph enumerations, connects the computations to the
quotient definitions, verifies positive-semidefiniteness exactly over
$\mathbb{Q}$, and automates the algebraic normalization steps that occur in
flag-algebra proofs.  As case studies, we obtain machine-checked proofs of Mantel's
theorem, the Erd\H{o}s pentagon theorem, Goodman-type inequalities, and
a $C_4$-density bound for induced-triangle-free graphs.

The same formalization also led us to new theoretical results about
constrained flag algebras.  We distinguish the built-in constraint semantics approach,
which builds a hereditary graph constraint, such as induced-\(H\)-freeness, into the
flag algebra from the start, from the ensemble semantics approach, which
imposes constraints afterwards by testing inequalities on random labeled views
of constrained graph limits.  The two approaches often agree, but they need
not agree for all densities of finite labeled patterns.  We prove a
root-planting criterion characterizing when they do agree, establish positive
results for classes closed under vertex blow-ups or more general
substitutions, and exhibit a counterexample based on induced-\(C_4\)-free graphs.  This
post hoc ensemble viewpoint also points beyond Razborov's original hereditary
setting: it can impose constraints on graph limits that are not themselves
hereditary, such as restricting to an extremal-density or equality slice.
\end{abstract}

% =========================================================================
\section{Introduction}
\label{sec:intro}

\paragraph{Why flag algebras matter.}
Many central questions in extremal combinatorics ask how often one finite
pattern can occur when another pattern is forbidden.  Let $P$ be a finite
simple graph and $\mathcal{H}$ a collection of finite simple graphs.  We write
\[
  \mathrm{ex}(n,\,P;\,\mathcal{H})
  :=
  \max\bigl\{
    \text{number of induced copies of }P\text{ in }G
    \mid |V(G)|=n,\; G\text{ is induced-}\mathcal{H}\text{-free}
  \bigr\},
\]
where an \emph{induced copy} of $P$ in $G$ is a vertex subset
$S\subseteq V(G)$ with $|S|=|V(P)|$ and $G[S]\cong P$, and
\emph{induced-\(\mathcal{H}\)-free} means that $G$ contains no induced copy of any
member of $\mathcal{H}$.  The corresponding asymptotic
question normalizes by the number of possible $|V(P)|$-vertex subsets and
takes a limit; we call this the \emph{induced Tur\'an density} of $P$ with
respect to $\mathcal{H}$:
\[
  \tdensity{P}{\mathcal{H}}
  :=
  \lim_{n\to\infty}
  \frac{\mathrm{ex}(n,\,P;\,\mathcal{H})}{\binom{n}{|V(P)|}}.
\]
This convention is slightly nonstandard.  In the classical Tur\'an-density
convention, ``\(H\)-free'' usually means free in the ordinary subgraph order,
and the counted copies of the target graph are also usually non-induced.  In
this paper, \(\tdensity{P}{\mathcal{H}}\) always refers to the induced version:
hosts are induced-\(\mathcal{H}\)-free and the statistic being optimized is
the induced \(P\)-density.  After this point we sometimes omit the word
``induced'' when no confusion should arise.
For example, Mantel's theorem states that an induced-triangle-free graph on $n$
vertices has at most $\lfloor n^2/4\rfloor$ edges---that is,
$\mathrm{ex}(n,\,K_2;\,K_3)=\lfloor n^2/4\rfloor$---achieved by the
complete bipartite graph $K_{\lfloor n/2\rfloor,\lceil n/2\rceil}$.  The
asymptotic value
\[
  \tdensity{K_2}{K_3}
  = \lim_{n\to\infty}\frac{\lfloor n^2/4\rfloor}{\binom{n}{2}}
  = \frac{1}{2}
\]
records that, in any sufficiently large induced-triangle-free graph, at most half of
all vertex pairs are edges.  (Whenever $\mathcal{H}=\{H\}$ is a singleton we
write $H$ in place of $\mathcal{H}$; so $K_3$ above abbreviates $\{K_3\}$.)

Determining $\tdensity{P}{\mathcal{H}}$ is hard because it is a statement
about \emph{all} sufficiently large induced-\(\mathcal{H}\)-free graphs: a valid upper
bound must guarantee that the proportion of induced copies of $P$ never exceeds
the claimed value, asymptotically, no matter how large the host graph is.  No
finite collection of examples suffices, and the extremal graphs achieving the
maximum may change unpredictably as $n$ grows.

Razborov's flag algebra method~\cite{razborov2007flag} is a proof technology
for exactly this setting.  A flag is a finite graph together with a
distinguished labeled subgraph.  By treating the densities of such partially
labeled patterns as algebraic objects, the method turns a global asymptotic
problem into a finite inequality in a quotient algebra.  Crucially, many
strong inequalities can be obtained from positive semidefinite matrices: an
semidefinite-programming (SDP) solver searches for a certificate, and a flag algebra theorem explains
why that finite certificate implies an asymptotic bound for the whole graph
class.  This combination of mathematical abstraction and computer search has been unusually
effective.  It underlies, for example, the solution of the Erd\H{o}s pentagon
problem~\cite{hatami2012,grzesik2012}, which determines the maximum asymptotic
density of pentagons in induced-triangle-free graphs.

\paragraph{From certificates to formal proofs.}
This paper presents a machine-checked formalization of flag algebras for
finite simple graphs, together with a certificate-to-proof pipeline that
reconstructs, in Lean, algebraic proofs suggested by external certificate
data.  The aim is not only to verify a few final inequalities, but to
formalize enough of the method that certificates found by
semidefinite-programming search can be checked against the mathematical
definitions inside the proof assistant.  The formalization therefore covers the
mathematical layer of the method: partially labeled graphs, their densities in
large graphs, the quotient algebra of density expressions, graph-limit
semantics through positive homomorphisms, and the downward operator that
transfers labeled reasoning back to unlabeled graph statements.

\paragraph{What Lean has to check.}
A flag-algebra certificate is compact, but its soundness is spread across
several different kinds of obligation.  First, there is an abstract semantic
obligation: graphs and flags are used only up to isomorphism, density
expressions are identified when they have the same asymptotic meaning, and a
claimed inequality must hold under every positive homomorphism.  Second, there
is a data-heavy computation obligation: the external search proposes finite
data, namely, lists of flags, density and multiplication tables, and rational
positive-semidefinite matrices, but Lean must check those data against the
formal definitions rather than trust the scripts that produced them.  Third,
there is an algebraic bookkeeping obligation: the proof repeatedly expands
flags to common sizes, multiplies typed flags, applies forbidden-graph
assumptions, and normalizes large linear combinations.  Our bridge layer
separates these concerns.  It connects concrete computable graph
representations to quotient definitions in Lean by adequacy theorems, verifies
positive-semidefiniteness exactly over $\mathbb{Q}$, and packages the
normalization steps into tactics and elaboration macros used by the case
studies.

\paragraph{Constrained semantics.}
Constrained extremal problems also expose a mathematical choice about when a
constraint is imposed.  In the \emph{built-in constraint semantics} approach, one builds
a hereditary graph constraint, such as induced-\(H\)-freeness, into the flag algebra
from the start; this is the traditional way to work in an algebra of
induced-\(H\)-free graphs.  In the \emph{ensemble semantics approach}, one first works
with graph limits and then imposes the constraint post hoc by testing
inequalities on random labeled views of those constrained limits.  These two viewpoints agree
for many purposes, but they do not automatically agree on every density of a
finite labeled pattern.  Our investigation proves a root-planting criterion
characterizing when they do agree, gives positive results for classes closed
under vertex blow-ups or more general substitutions, and exhibits a
counterexample based on induced-\(C_4\)-free graphs.  The ensemble viewpoint is useful
beyond this comparison: because it imposes constraints post hoc on graph
limits, it can also express non-hereditary restrictions such as restricting
to an extremal-density or equality slice.  This is also the semantic viewpoint
made explicit by our formalization.

\paragraph{Contributions.}
The paper makes four main contributions.  First, it gives a Lean
formalization of the core flag-algebra objects for finite simple graphs:
flags, densities, the quotient algebra, positive homomorphisms, semantic
order, and the downward/random-homomorphism machinery used in flag-algebra
proofs.  Second, it builds a certificate-to-proof pipeline that turns
externally generated rational certificate data into Lean-checked algebraic
proofs.  Third, it demonstrates the pipeline on machine-checked proofs of
Mantel's theorem, the Erd\H{o}s pentagon theorem, Goodman-type inequalities,
and a $C_4$-density bound for induced-triangle-free graphs.  Fourth, it develops the
semantic comparison between built-in constraint and ensemble approaches to constrained
flag algebras, including the root-planting criterion and the resulting
generalization beyond hereditary constraints.

\paragraph{Organization.}
Section~\ref{sec:background} recalls the mathematical background of flag algebras.
Section~\ref{sec:abstract} describes the Lean specification layer and the
semantic treatment of forbidden-graph assumptions.
Section~\ref{sec:obstacles} describes the main formalization obstacles:
quotient eliminators, dependent-type transports, function extensionality, and
\lean{Fintype}-instance ambiguity.
Sections~\ref{sec:reflection} and~\ref{sec:bridge} describe the bridge from
specification to evaluable proof: reflection, SDP-certificate verification,
the tactic layer, the API layer, and elaboration-time theorem generation.
Section~\ref{sec:proof-style} records a recurring proof pattern from the case
studies: density equalities are proved by explicit bijections between finite
sample spaces.  Section~\ref{sec:results} presents the verified case studies.
Sections~\ref{sec:unsure} and~\ref{sec:lessons} discuss limitations and
lessons for future formalizations, and Sections~\ref{sec:related}
and~\ref{sec:conclusion} give related work and conclusions.


\section{Background: Flag Algebras}
\label{sec:background}

This section recalls the mathematical definitions underlying flag algebras,
following Razborov~\cite{razborov2007flag}, and sets up the notation used
throughout the paper.  No background in flag algebras or extremal
combinatorics is assumed.  For $n \in \mathbb{N}$ we write
$[n] = \{1,\ldots,n\}$. All graphs $G$ are undirected, finite and simple,
and their vertex and edge sets are written $V(G)$ and $E(G)$, respectively.
Throughout this section we fix a (possibly empty) finite set $\mathcal{H}$ of
forbidden induced subgraphs.  If $G$ is a graph and $S\subseteq V(G)$, the
\emph{induced subgraph} $G[S]$ is the graph on vertex set $S$ whose edges are
exactly the edges of $G$ with both endpoints in $S$.  Thus a graph is
\emph{induced-\(\mathcal{H}\)-free} if there is no set $S\subseteq V(G)$ for which
$G[S]$ is isomorphic to a member of $\mathcal{H}$.  All flag types, flags, and flag
algebras below are defined relative to this fixed $\mathcal{H}$.
% The goal is to make the proof pattern legible:
% what are the terms, what do they mean, and why can a finite certificate say
% something about all large graphs?

% For $n \in \mathbb{N}$ we write $[n] = \{1,\ldots,n\}$.
% All graphs are undirected, finite and simple.
% To keep roles separate, this section uses $P$ for a target graph, $H$ for a
% forbidden graph or intermediate flag, $F$ for a pattern flag, $G$ for a host graph or host flag,
% $Q$ for a positive semidefinite matrix, and $\phi$ for a positive
% homomorphism.

% The motivating question has the following shape.  Fix a target pattern $P$
% and a collection $\mathcal{H}$ of forbidden graphs.  Among all very large
% graphs that contain no induced subgraph isomorphic to any member of
% $\mathcal{H}$, how often can $P$ appear as an induced subgraph?  We call such
% graphs \emph{$\mathcal{H}$-free}.  For example, if $\mathcal{H}=\{K_3\}$ is
% the singleton containing the triangle and $P=C_5$ is a pentagon, the question
% asks for the largest possible pentagon density in triangle-free graphs.  Such
% questions are hard because the graph being optimized over is not fixed: it can
% have arbitrarily many vertices.

% Flag algebras are useful because they turn this infinite family of finite
% optimization problems into a finite certificate-checking problem.  One can
% read the method as a small domain-specific proof language:
% \emph{terms} are linear combinations of small graph patterns,
% \emph{equations} come from exact finite counting identities, and
% \emph{non-negativity proofs} come from finite quadratic certificates.  Once
% the definitions below are in place, a successful certificate will imply that
% every limiting $H$-free graph has target density at most the claimed constant.
% Flag algebras answer this question by exhibiting explicit finite certificates.
% The coefficients in such a certificate are rational numbers obtained from finite
% counting experiments; non-negativity is witnessed by a positive-semidefinite
% matrix over $\mathbb{Q}$.  The definitions below make these notions precise and
% explain why a finite certificate of this form implies a bound on
% $\pi(P;\,\mathcal{H})$.

% The rest of the section introduces this language one construct at a time.
% Types and flags are its syntax for local, rooted graph observations.  Densities
% give those observations a probabilistic meaning.  The quotient algebra installs
% the counting identities as definitional equalities of the proof language.
% We then introduce positive homomorphisms, which define what it means for one
% algebra expression to be non-negative or bounded above by another.  Finally,
% an unlabeling map transports labeled local reasoning back to ordinary
% unlabeled graph statements.

\subsection{Flag Types and Flags}

% The first move in a flag algebra proof is to look at a large graph from the
% perspective of a small labeled configuration.  This is what lets the method use
% conditional information: for example, instead of only asking how many edges a
% triangle-free graph has, we can ask how the graph looks around a distinguished
% vertex or a distinguished edge.  Types and flags are the notation for these
% local views.

\paragraph{Flag Types.}
A \emph{flag type} of size $k$ is an induced-\(\mathcal{H}\)-free graph \(\sigma\) with vertex set $[k]$.
The \emph{empty flag type} $\emptyset$ (size $k=0$) is the unique graph on the
empty vertex set.
We use the term \emph{flag type} rather than Razborov's original \emph{type}
to avoid confusion with Lean's type system in later sections.

\paragraph{Flags.}
Fix a flag type $\sigma$ of size $k$.
A \emph{$\sigma$-flag} is a pair $F = (M, \theta_F)$ where $M$ is an
induced-\(\mathcal{H}\)-free finite graph and $\theta_F : [k] \hookrightarrow V(M)$ is an injective map such
that $\sigma$ is isomorphic to $M[\operatorname{Im}(\theta_F)]$ via $\theta_F$
(i.e., $\theta_F$ is a graph embedding of $\sigma$ into $M$).
The integer $|V(M)|$ is the \emph{size} of $F$; we write $V(F)$ for $V(M)$.
The vertices in $\operatorname{Im}(\theta_F)$ are called the
\emph{labeled vertices} of $F$; the remaining $|V(F)|-k$ vertices are
\emph{unlabeled}.
For a set $S \subseteq V(F)$ with $\operatorname{Im}(\theta_F) \subseteq S$,
we write $F[S]$ for the $\sigma$-flag $(M[S], \theta_F)$ and call it the
\emph{$\sigma$-flag induced by $F$ on $S$}.
Note that flags over the empty flag type $\emptyset$ are just ordinary finite graphs.
\Cref{fig:flag-examples} illustrates several flags of different types.

Two $\sigma$-flags $(M_1,\theta_1)$ and $(M_2,\theta_2)$ are
\emph{isomorphic}, denoted by $(M_1,\theta_1) \cong_\sigma (M_2,\theta_2)$,
if there is a graph isomorphism $\alpha: V(M_1)\to V(M_2)$
with $\alpha \circ \theta_1 = \theta_2$.  The set of isomorphism classes of
$\sigma$-flags of size $n$ is written $\mathcal{F}^\sigma_n$, and
$\mathcal{F}^\sigma = \bigcup_{n \geq k} \mathcal{F}^\sigma_n$.
Each $\mathcal{F}^\sigma_n$ is finite, although
$\mathcal{F}^\sigma$ itself is infinite.

\begin{figure}[t]
  \centering
  \setlength{\tabcolsep}{1.2em}
  \renewcommand{\arraystretch}{1.3}
  \begin{tabular}{cccccc}
    % ∅-flags (ordinary graphs)
    \begin{tikzpicture}[scale=0.8, baseline=-0.5ex]
      \node[graph vertex] (a) at (0,0) {};
      \node[graph vertex] (b) at (0.8,0) {};
      \draw[graph edge] (a) -- (b);
    \end{tikzpicture}
    &
    \begin{tikzpicture}[scale=0.8, baseline=-0.5ex]
      \node[graph vertex] (a) at (90:0.65) {};
      \node[graph vertex] (b) at (210:0.65) {};
      \node[graph vertex] (c) at (330:0.65) {};
    \end{tikzpicture}
    &
    \begin{tikzpicture}[scale=0.8, baseline=-0.5ex]
      \node[graph vertex] (a) at (90:0.65) {};
      \node[graph vertex] (b) at (210:0.65) {};
      \node[graph vertex] (c) at (330:0.65) {};
      \draw[graph edge] (b) -- (c);
    \end{tikzpicture}
    &
    \begin{tikzpicture}[scale=0.8, baseline=-0.5ex]
      \node[graph vertex] (a) at (90:0.65) {};
      \node[graph vertex] (b) at (210:0.65) {};
      \node[graph vertex] (c) at (330:0.65) {};
      \draw[graph edge] (a) -- (b);
      \draw[graph edge] (a) -- (c);
    \end{tikzpicture}
    &
    \begin{tikzpicture}[scale=0.8, baseline=-0.5ex]
      \node[graph vertex] (a) at (90:0.65) {};
      \node[graph vertex] (b) at (210:0.65) {};
      \node[graph vertex] (c) at (330:0.65) {};
      \draw[graph edge] (a) -- (b);
      \draw[graph edge] (a) -- (c);
      \draw[graph edge] (b) -- (c);
    \end{tikzpicture}
    &
    \begin{tikzpicture}[scale=0.7, baseline=-0.5ex]
      \node[graph vertex] (v1) at (90:0.65) {};
      \node[graph vertex] (v2) at (18:0.65) {};
      \node[graph vertex] (v3) at (-54:0.65) {};
      \node[graph vertex] (v4) at (-126:0.65) {};
      \node[graph vertex] (v5) at (162:0.65) {};
      \draw[graph edge] (v1) -- (v2) -- (v3) -- (v4) -- (v5) -- (v1);
    \end{tikzpicture}
    \\
    $K_2$ & $\overline{K_3}$ & $\overline{P_3}$ & $P_3$ & $K_3$ & $C_5$ \\[0.8em]
    % 1-vertex type σ-flags
    \begin{tikzpicture}[scale=0.8, baseline=-0.5ex]
      \node[graph root] (a) at (0,0) {};
      \node[graph vertex] (b) at (0.8,0) {};
      \draw[graph edge] (a) -- (b);
    \end{tikzpicture}
    &
    \begin{tikzpicture}[scale=0.8, baseline=-0.5ex]
      \node[graph root] (a) at (0,0) {};
      \node[graph vertex] (b) at (0.8,0) {};
    \end{tikzpicture}
    &
    \begin{tikzpicture}[scale=0.8, baseline=-0.5ex]
      \node[graph vertex] (a) at (90:0.65) {};
      \node[graph root] (b) at (210:0.65) {};
      \node[graph vertex] (c) at (330:0.65) {};
      \draw[graph edge] (b) -- (c);
    \end{tikzpicture}
    &
    \begin{tikzpicture}[scale=0.8, baseline=-0.5ex]
      \node[graph vertex] (a) at (90:0.65) {};
      \node[graph root] (b) at (210:0.65) {};
      \node[graph vertex] (c) at (330:0.65) {};
      \draw[graph edge] (b) -- (a);
      \draw[graph edge] (a) -- (c);
    \end{tikzpicture}
    &
    &
    \begin{tikzpicture}[scale=0.7, baseline=-0.5ex]
      \node[graph root] (v1) at (90:0.65) {};
      \node[graph vertex] (v2) at (18:0.65) {};
      \node[graph vertex] (v3) at (-54:0.65) {};
      \node[graph vertex] (v4) at (-126:0.65) {};
      \node[graph vertex] (v5) at (162:0.65) {};
      \draw[graph edge] (v1) -- (v2) -- (v3) -- (v4) -- (v5) -- (v1);
    \end{tikzpicture}
    \\
    $K_2^\bullet$ & $\overline{K_2}^\bullet$ & $\overline{P_3}^\bullet$ & $P_3^\bullet$
    & & $C_5^\bullet$ \\
  \end{tabular}
  \caption{Examples of flags.  \emph{Top row}: $\emptyset$-flags (ordinary graphs)
  of sizes 2, 3, and 5.  \emph{Bottom row}: flags of the one-vertex flag type $\sigma$;
  the filled vertex is the labeled one.}
  \label{fig:flag-examples}
\end{figure}

\subsection{Subflag Densities}

In a density calculation, flags play two roles.  The \emph{pattern flag} is the
small local configuration whose occurrence is being measured, while the
\emph{host flag} is the larger flag in which we sample vertices.  These are
roles in the calculation, not different kinds of flags.  We use \(F\), or
\(F_i\), for pattern flags, \(G\) for host flags, and \(H\) for an intermediate
flag when a third flag variable is needed.

% Once pattern flags have been named, the basic question is simple: if we sample
% a few extra vertices in a larger host flag, what smaller flag do we see?
% The resulting densities are the bridge between finite host flags and
% asymptotic limits.  Along a convergent sequence of larger and larger flags,
% each fixed finite pattern has a limiting density.  Flag algebra manipulates
% those limiting densities symbolically, but every coefficient comes from the
% finite probability definitions below.  This is why the method is well suited
% to machine checking: the coefficients in a certificate are not analytic
% guesses, but rational numbers obtained by finite enumeration.

\paragraph{Single-flag density.}
Let $\sigma$ be a flag type of size $k$.  For a $\sigma$-flag $F$ of size $m$
(the pattern flag) and a $\sigma$-flag $G = (M, \theta_G)$ of size $n \geq m$
(the host flag), the \emph{density of $F$ in $G$}, written $\den{F}{G}$, is
the probability that a uniformly random $(m-k)$-element subset of
$V(G) \setminus \operatorname{Im}(\theta_G)$, together with the $k$ labeled
vertices of $G$, induces a $\sigma$-flag isomorphic to $F$.
% The \emph{density of $F$ in $G$}, written $\den{F}{G}$, is the probability
% that a uniformly random $(m-k)$-element subset $S$ of
% $V(G) \setminus \operatorname{Im}(\theta_G)$ satisfies
% $G[S \cup \operatorname{Im}(\theta_G)] \cong_\sigma F$.
Explicitly:
\[
  \den{F}{G}
  \;:=\;
  \frac{1}{\binom{n-k}{m-k}}
    \cdot {\Bigl|\Big\{S \subseteq V(G)\setminus \operatorname{Im}(\theta_G)
                  ~\Big|~ |S|=m-k,\;
                  G[S \cup \operatorname{Im}(\theta_G)]
                  \cong_\sigma F \Big\}\Bigr|}.
\]
This value lies in $[0,1]$ and is invariant under
isomorphism of both $F$ and $G$.
When $n < m$, we set $\den{F}{G} := 0$ by convention.


\begin{example}[Subflag densities in $C_5^\bullet$]
  Among the four unlabeled vertices of
  $\fCfivebull$, exactly two are adjacent to the labeled one.  Hence
  \[
    \den{\fKtwobull}{\fCfivebull} = \frac{2}{4} = \frac12,
    \qquad
    \den{\fbarKtwobull}{\fCfivebull} = \frac{2}{4} = \frac12.
  \]
  % This tiny example captures the reason flags are useful: a flag remembers
  % the local view from a labeled configuration.  The same unrooted pentagon can
  % be queried from the perspective of a distinguished vertex, and the algebra
  % keeps those conditional densities as named algebraic coordinates rather than
  % side comments in prose.
\end{example}

\paragraph{Chain rule.}
The density $\den{F}{G}$ can be computed in two equivalent ways: by sampling
$m-k$ unlabeled vertices directly from $G$, or by first sampling a larger
intermediate flag $H$ of size $n$ inside $G$ and then sampling $F$ inside $H$.
These two procedures are the same experiment, which gives the following identity.

\begin{lemma}[Chain rule]\label{lem:chain-rule}
  For $F \in \mathcal{F}^\sigma_m$, $n \geq m$, and a $\sigma$-flag $G$ of size
  $\geq n$,
  \[
    \den{F}{G}
    \;=\;
    \sum_{H \in \mathcal{F}^\sigma_n} \den{F}{H}\,\den{H}{G}.
  \]
\end{lemma}

\paragraph{Multi-flag density.}
For pattern $\sigma$-flags $F_1, \ldots, F_t$ of sizes $m_1, \ldots, m_t$ and
a host $\sigma$-flag $G = (M, \theta_G)$ of size $n$, we say $F_1, \ldots, F_t$
\emph{fit} in $G$ if
\[
  n - k \;\geq\; (m_1 - k) + \cdots + (m_t - k).
\]
When $F_1, \ldots, F_t$ fit in $G$, let $(S_1, \ldots, S_t)$ be a uniformly
random tuple of pairwise-disjoint subsets $S_i \subseteq V(G) \setminus
\operatorname{Im}(\theta_G)$ with $|S_i| = m_i - k$.
The \emph{density of $F_1, \ldots, F_t$ in $G$}, written $\den{F_1,\ldots,F_t}{G}$, is the
probability that $G[S_i \cup \operatorname{Im}(\theta_G)] \cong_\sigma F_i$
for every $i \in [t]$.
When $F_1, \ldots, F_t$ do not fit in $G$, we set
$\den{F_1,\ldots,F_t}{G} := 0$ by convention.
% Joint densities are what make multiplication possible: the product of two
% pattern flags is interpreted as the probability that two disjoint samples occur
% simultaneously in the same host.  This turns combinatorial double counting
% into algebra.

\subsection{The Flag Algebra}

Fix a flag type $\sigma$ of size $k$.
Informally, the flag algebra $\mathcal{A}^\sigma$ is an algebra\footnote{Here
``algebra'' means a vector space with bilinear multiplication and a unit,
satisfying the usual laws; the flag algebras in this paper are commutative
$\mathbb{R}$-algebras.} of density expressions generated
by $\sigma$-flags.  Each flag is a symbolic building block: it stands for the
density of that local labeled pattern in a large host graph.  The quotient
below records the standard sampling identities among these building blocks,
beginning with the expansion identity coming from
Lemma~\ref{lem:chain-rule}.
Formally, let
$\mathbb{R}[\mathcal{F}^\sigma]$ be the free $\mathbb{R}$-module with basis
$\mathcal{F}^\sigma$, i.e., the vector space of finite formal real linear
combinations of $\sigma$-flags.  We also write $F$ for the basis vector
corresponding to a flag $F$.
Define the \emph{zero space} $\mathcal{Z}^\sigma$ as the subspace generated
by all elements of the form
\begin{equation}
  \label{eq:zero-element}
  F - \sum_{H \in \mathcal{F}^\sigma_n} \den{F}{H} \cdot H,
  \qquad \text{for}\ F \in \mathcal{F}^\sigma_m \ \text{and}\ n \geq m.
\end{equation}
The \emph{flag algebra} is the quotient module
\[
  \mathcal{A}^\sigma \;:=\; \mathbb{R}[\mathcal{F}^\sigma] \,/\, \mathcal{Z}^\sigma.
\]
We write $[F]$ for the equivalence class of a flag $F$ in $\mathcal{A}^\sigma$.
By abuse of notation, we write $F$ for $[F]$ whenever the distinction between
a flag and its equivalence class is immaterial.

We call $\sum_{H \in \mathcal{F}^\sigma_n} \den{F}{H} \cdot H$ the
\emph{expansion of $F$ at size $n$}.
By Lemma~\ref{lem:chain-rule}, a flag $F \in \mathcal{F}^\sigma_m$ and
its expansion at any size $n \geq m$ have the same density in every host flag of size $\geq n$.
The quotient by $\mathcal{Z}^\sigma$ records this counting identity as the
equality
\[
  [F]
  =
  \sum_{H \in \mathcal{F}^\sigma_n} \den{F}{H}\,[H]
  \qquad\text{in }\mathcal{A}^\sigma .
\]
We may therefore freely replace $F$ by its expansion at any size $n \geq m$.
% The zero-space relations express the combinatorial identity that, in a
% sufficiently large host flag $G$,
% \[
%   \den{F}{G}
%   =
%   \sum_{H \in \mathcal{F}^\sigma_n} \den{F}{H}\,\den{H}{G}.
% \]
% That is, counting $F$ directly agrees with counting through an intermediate
% sample of size $n$.

\begin{example}[$\fKtwo$ expanded at size $3$ in the induced-triangle-free algebra]
  Working in the induced-triangle-free flag algebra (i.e., $\mathcal{H} = \{K_3\}$),
  the $\emptyset$-flags of size $3$ are the three induced-triangle-free graphs on three vertices:
  $\fbarKthree$, $\fbarPthree$, and $\fPthree$.
  The density of an edge in each of them is
  \[
    \den{\fKtwo}{\fbarKthree}=0,\quad
    \den{\fKtwo}{\fbarPthree}=\frac13,\quad
    \den{\fKtwo}{\fPthree}=\frac23,
  \]
  so the zero-space quotient identifies
  \begin{equation}
    \fKtwo
    \;=\;
    \frac13\,\fbarPthree
    + \frac23\,\fPthree
    \qquad\text{in } \mathcal{A}^{\emptyset}.
    \label{eq:k2-expansion}
  \end{equation}
\end{example}

\paragraph{Multiplication.}
For flags $F_1 \in \mathcal{F}^\sigma_{m_1}$ and
$F_2 \in \mathcal{F}^\sigma_{m_2}$, choose an auxiliary size
$\ell \geq m_1 + m_2 - k$ and define
\begin{equation}
  \label{eq:mul}
  [F_1] \cdot [F_2]
  \;:=\;
  \sum_{H \in \mathcal{F}^\sigma_\ell} \den{F_1,F_2}{H} \cdot [H]
  \qquad\text{in } \mathcal{A}^\sigma.
\end{equation}
At first glance this definition appears to depend on both the choice of $\ell$
and the choice of representatives for $[F_1]$ and $[F_2]$.  The zero-space
quotient ensures neither matters: the product is independent of $\ell$ and
respects isomorphism classes of the input flags.  Extending bilinearly to all
of $\mathcal{A}^\sigma$ gives a commutative, associative $\mathbb{R}$-algebra.
The unit of this algebra is the $\sigma$-flag $(\sigma, \operatorname{id}_{[k]})$ of size $k$,
in which the underlying graph is $\sigma$ itself and every vertex is labeled;
we write it as $\mathbf{1}_\sigma$, and simply as $\mathbf{1}$ when
$\sigma=\emptyset$.

\begin{example}[Multiplication of $K_2^\bullet$ and $\overline{K_2}^\bullet$]
  Let $\sigma$ be the one-vertex flag type.  Taking $\ell = 3$, the density
  $\den{\fKtwobull,\,\fbarKtwobull}{G}$ equals $\frac{1}{2}$ for each size-3
  $\sigma$-flag $G$ whose labeled vertex has degree exactly $1$:
  \[
    \den{\fKtwobull,\,\fbarKtwobull}{\fbarPthreebull}
    \;=\;
    \den{\fKtwobull,\,\fbarKtwobull}{\fPthreebull}
    \;=\; \frac{1}{2},
  \]
  and vanishes on all other size-3 $\sigma$-flags:
  \[
    \den{\fKtwobull,\,\fbarKtwobull}{G} = 0
    \qquad\text{for all } G \in \mathcal{F}_3^\sigma \setminus \{\fbarPthreebull,\,\fPthreebull\}.
  \]
  Hence
  \[
    \fKtwobull \cdot \fbarKtwobull
    \;=\;
    \frac{1}{2}\,\fbarPthreebull + \frac{1}{2}\,\fPthreebull
    \qquad\text{in } \mathcal{A}^\sigma.
  \]
\end{example}

\subsection{Positive Homomorphisms and Semantic Order}
\label{sec:background-semantic-order}

% The quotient algebra is syntax for density expressions.  To use such an
% expression in an extremal argument, we need to say how it is evaluated in a
% limiting graph sequence.  That role is played by positive homomorphisms.

The flag algebra $\mathcal{A}^\sigma$ provides algebraic syntax for expressing
densities of pattern $\sigma$-flags.  Flag algebra proofs, however, target
\emph{asymptotic} statements: rather than a single finite host flag, we
consider increasing sequences $\{G_n\}_{n\in\mathbb{N}}$ of host flags with $|V(G_n)|\to\infty$ and ask
what happens to flag densities in the limit.  Positive homomorphisms can be thought of as the limiting objects of convergent
sequences of host flags.  The semantic non-negativity cone
$\mathcal{C}^\sigma$ then consists of the pattern $\sigma$-flag expressions
whose density at every such limiting object is non-negative, and establishing
membership in $\mathcal{C}^\sigma$ is precisely what flag algebra proofs aim
to do.

\paragraph{Convergent sequences.}
Call a sequence $\{G_n\}_{n\in\mathbb{N}}$ of host $\sigma$-flags \emph{increasing} if
$|V(G_1)| < |V(G_2)| < \cdots$.  Each host $\sigma$-flag $G$ defines a point
$p^G \in [0,1]^{\mathcal{F}^\sigma}$ by $p^G(F) := \den{F}{G}$.
An increasing sequence is \emph{convergent} if the sequence of points $p^{G_n}$
converges in $[0,1]^{\mathcal{F}^\sigma}$ endowed with the product topology;
equivalently, if $\lim_{n\to\infty}\den{F}{G_n}$ exists for every $\sigma$-flag $F$.

\paragraph{Positive homomorphisms.}
Fix a flag type $\sigma$.  A \emph{positive homomorphism} on
$\mathcal{A}^\sigma$ is an $\mathbb{R}$-algebra homomorphism
$\phi:\mathcal{A}^\sigma \to \mathbb{R}$ such that
$\phi(F) \geq 0$ for every $\sigma$-flag $F$.  Thus $\phi$ assigns a real
number to each flag algebra expression, respects addition, multiplication, and the unit, and
never assigns a negative value to any flag.
We write $\poshom{\sigma}$ for the set of all positive homomorphisms on
$\mathcal{A}^\sigma$.

In $\mathcal{A}^\sigma$, the following identity holds for every $\ell \geq |V(\sigma)|$:
\begin{equation}
  \sum_{F \in \mathcal{F}^\sigma_\ell} F \;=\; \mathbf{1}_\sigma.
  \label{eq:sum-to-one}
\end{equation}
Applying any $\phi \in \poshom{\sigma}$ to both sides gives
$\sum_{F \in \mathcal{F}^\sigma_\ell} \phi(F) = 1$;
non-negativity then forces $\phi(F) \in [0,1]$ for every $\sigma$-flag $F$.
Thus, each $\phi \in \poshom{\sigma}$ determines a point in $[0,1]^{\mathcal{F}^\sigma}$
via the assignment $F \mapsto \phi(F)$, so $\poshom{\sigma}$ can be identified with
a subset of $[0,1]^{\mathcal{F}^\sigma}$.  % It is moreover compact.

The following theorem makes precise the sense in which positive homomorphisms
are exactly the limiting objects of convergent sequences of $\sigma$-flags.

\begin{theorem}[Lov\'{a}sz--Szegedy~{\cite{lovasz2006limits}}]\label{thm:convergent-hom}
  \begin{enumerate}
  \item[\textup{(a)}] For every convergent sequence $\{G_n\}_{n\in\mathbb{N}}$ of
    $\sigma$-flags, $\lim_{n\to\infty} p^{G_n} \in \poshom{\sigma}$.
  \item[\textup{(b)}] Conversely, every $\phi \in \poshom{\sigma}$ satisfies
    $\phi = \lim_{n\to\infty} p^{G_n}$ for some convergent sequence
    $\{G_n\}_{n\in\mathbb{N}}$ of $\sigma$-flags.
  \end{enumerate}
\end{theorem}

\paragraph{Semantic non-negativity cone and order.}
For a fixed flag type $\sigma$, define the \emph{semantic non-negativity cone}
\[
  \mathcal{C}^\sigma
  :=
  \{f \in \mathcal{A}^\sigma
    \mid \phi(f) \geq 0
    \text{ for every } \phi \in \poshom{\sigma}\}.
\]
We write
\[
  f \leq_\sigma g
  \qquad\text{when}\qquad
  g-f \in \mathcal{C}^\sigma.
\]
Equivalently, $f\leq_\sigma g$ means that
$\phi(f)\leq\phi(g)$ for every $\phi \in \poshom{\sigma}$.  When $\sigma=\emptyset$, we write $f \leq g$ in place of $f \leq_\emptyset g$, and $\mathcal{C}$ in place of $\mathcal{C}^\emptyset$.
Note that $f \in \mathcal{C}^\sigma$ is the same as saying $0 \leq_\sigma f$.

Two immediate examples of elements of $\mathcal{C}^\sigma$ are flags and squares.
Every $\sigma$-flag $F$ satisfies $0 \leq_\sigma F$,
because positive homomorphisms send every flag to a non-negative real by definition.
For any $f \in \mathcal{A}^\sigma$, the square $f^2$ satisfies $0 \leq_\sigma f^2$,
because $\phi(f^2) = \phi(f)^2 \geq 0$ for every $\phi \in \poshom{\sigma}$.
Moreover, if $0 \leq_\sigma f$ and $0 \leq_\sigma g$, then $0 \leq_\sigma f + g$,
since $\phi(f+g) = \phi(f) + \phi(g) \geq 0$ for every $\phi \in \poshom{\sigma}$.
Similarly, if $0 \leq_\sigma f$ and $c$ is a non-negative real number, then $0 \leq_\sigma c \cdot f$.

\paragraph{From semantic non-negativity to induced Tur\'an density.}
The problem of bounding the induced Tur\'an density $\tdensity{P}{\mathcal{H}}$ from above
reduces to establishing membership in $\mathcal{C}^\emptyset$, as the following
lemma shows.

\begin{lemma}\label{lem:semantic-bound}
  If $P \leq c \cdot \mathbf{1}$ for some $c \in \mathbb{R}$, then
  $\tdensity{P}{\mathcal{H}} \leq c$.
\end{lemma}
\begin{proof}
  Suppose for contradiction that $\tdensity{P}{\mathcal{H}} > c$.  Then
  there exist $\varepsilon > 0$ and an increasing sequence of
  induced-\(\mathcal{H}\)-free graphs $\{G_k\}_{k\in\mathbb{N}}$ such that
  $\den{P}{G_k} > c + \varepsilon$ for all $k$.  Since each $p^{G_k}$
  is a point in $[0,1]^{\mathcal{F}^\emptyset}$ and this product space is
  compact (Tychonoff), the sequence $\{p^{G_k}\}$ has a convergent
  subsequence $\{p^{G_{k_j}}\}$; by \Cref{thm:convergent-hom}(a) its limit is some
  $\phi \in \poshom{\emptyset}$.  Since $\phi$ is the pointwise limit of
  $\{p^{G_{k_j}}\}$, we have
  $\phi(P) = \lim_j \den{P}{G_{k_j}} \geq c + \varepsilon > c$.  But
  $P \leq c \cdot \mathbf{1}$ means $\phi(P) \leq c$ for every
  $\phi \in \poshom{\emptyset}$, a contradiction.
\end{proof}

A matching lower bound typically comes from an explicit induced-\(\mathcal{H}\)-free construction
whose $P$-density approaches $c$, pinning down $\tdensity{P}{\mathcal{H}}$ exactly.
Two classical results illustrate how this looks in flag algebra language.

\begin{example}[Mantel's theorem]
  Mantel's theorem states $\tdensity{K_2}{K_3} = \tfrac{1}{2}$, meaning
  asymptotically at most half of all vertex pairs in an induced-triangle-free graph
  are edges.  The upper bound translates into the flag algebra inequality
  $K_2 \leq \tfrac{1}{2}\cdot\mathbf{1}$; the balanced complete bipartite
  graphs $K_{\lfloor n/2\rfloor,\lceil n/2\rceil}$ witness the lower bound.
\end{example}

\begin{example}[Erd\H{o}s Pentagon Problem~{\cite{grzesik2012,hatami2012}}]
  The Erd\H{o}s pentagon problem asks for the maximum density of $C_5$ in
  induced-triangle-free graphs; the asymptotic answer is $\tdensity{C_5}{K_3} = \tfrac{24}{625}$.  The upper bound corresponds to the flag algebra inequality
  $C_5 \leq \tfrac{24}{625}\cdot\mathbf{1}$; an explicit induced-\(K_3\)-free construction
  (the balanced blow-up of $C_5$) witnesses the lower bound.
\end{example}

% \paragraph{Squares and PSD quadratic forms.}
% A real symmetric matrix $Q \in \mathbb{R}^{r\times r}$ is \emph{positive
% semidefinite} (PSD) if $v^\top Q v\geq 0$ for every real vector
% $v \in \mathbb{R}^r$.  The simplest source of elements of
% $\mathcal{C}^\sigma$ is a square: if $e\in\mathcal{A}^\sigma$ and $\phi$ is a
% positive homomorphism, then
% \[
%   \phi(e^2)=\phi(e)^2\geq 0.
% \]
% So $e^2$ is semantically non-negative.  A finite sum
% $\sum_i a_i e_i^2$ with $a_i\geq 0$ is non-negative for the same reason.  More
% generally, if $Q$ is PSD and $e_1,\ldots,e_r\in\mathcal{A}^\sigma$, then every
% positive homomorphism sends
% \[
%   \sum_{i,j=1}^r Q_{ij}\,e_i e_j
% \]
% to the real quadratic form $v^\top Q v$ for
% $v=(\phi(e_1),\ldots,\phi(e_r))$, which is always non-negative.  So the
% above quadratic expression is in $\mathcal{C}^\sigma$.
% This is the sense in which flag algebra
% proofs use ``sums of squares'': the proof data is a finite quadratic expression
% whose value is forced to be a non-negative real number by a PSD matrix.

\subsection{The Downward Operator}
\label{sec:background-downward}
Labeled flag algebras (i.e., $\mathcal{A}^\sigma$ with $\sigma \neq \emptyset$) are often
richer than the unlabeled algebra $\mathcal{A}^\emptyset$: establishing membership
in $\mathcal{C}^\sigma$ can be easier than working directly in $\mathcal{C}$.
The final extremal statement, however, is always about unlabeled graphs.
The downward operator bridges this gap by mapping $\mathcal{A}^\sigma$ to
$\mathcal{A}^\emptyset$ in a way that preserves semantic non-negativity
(\Cref{thm:downward-nonneg}), so valid inequalities proved in a labeled algebra
can be transferred to the unlabeled one.

For a $\sigma$-flag $F = (M, \theta_F)$, let $F|_\emptyset$ denote the
$\emptyset$-flag obtained from $F$ by forgetting the embedding $\theta_F$,
i.e., viewing $M$ as a $\emptyset$-flag.
For a $\sigma$-flag $F$ of size $m$ with $|V(\sigma)|=k$, let
\[
  q_\sigma(F)
  \;:=\;
  \frac{
    \bigl|\{\theta' : [k]\hookrightarrow V(M)
       \mid \theta' \text{ embeds }\sigma\text{ into }M
       \text{ and } (M,\theta')\cong_\sigma F\}\bigr|}
       {m(m-1)\cdots(m-k+1)}
\]
be the probability that a uniformly random injection $[k]\hookrightarrow V(M)$
recovers a flag isomorphic to $F$.
The \emph{downward operator}
$\llbracket\cdot\rrbracket_\sigma:\mathcal{A}^\sigma\to\mathcal{A}^\emptyset$
is then defined on flags by
\[
  \llbracket F \rrbracket_\sigma
  \;:=\;
  q_\sigma(F)\cdot F|_\emptyset,
\]
extended linearly to all of $\mathcal{A}^\sigma$.

\paragraph{Random homomorphisms.}
A positive homomorphism $\phi_0 \in \poshom{\emptyset}$ represents a limiting
object for unlabeled graph sequences.  Given such a $\phi_0$ and a non-empty
flag type $\sigma$, one can generate a \emph{random positive homomorphism} on
$\mathcal{A}^\sigma$ by sampling a random embedding of $\sigma$ into the
limiting object.
Whenever $\phi_0(\sigma) > 0$ (i.e., $\sigma$ has positive density as an unlabeled graph in the limiting object),
the following holds~\cite{razborov2007flag}:
there exists a unique probability distribution on $\poshom{\sigma}$,
which we denote $\Ext_\sigma(\phi_0)$,
such that for every $f \in \mathcal{A}^\sigma$, the random positive
homomorphism $\boldsymbol{\phi}^\sigma \sim \Ext_\sigma(\phi_0)$ satisfies
\begin{equation}\label{eq:ensemble}
  \mathbb{E}\bigl[\boldsymbol{\phi}^\sigma(f)\bigr]
  \;=\;
  \frac{\phi_0\bigl(\llbracket f \rrbracket_\sigma\bigr)}
       {\phi_0\bigl(\llbracket \mathbf{1}_\sigma \rrbracket_\sigma\bigr)}.
\end{equation}
Note that the denominator is positive by the assumption $\phi_0(\sigma) > 0$, since
$\phi_0\bigl(\llbracket \mathbf{1}_\sigma \rrbracket_\sigma\bigr) = q_\sigma(\mathbf{1}_\sigma)\cdot\phi_0(\sigma)$
and $q_\sigma(\mathbf{1}_\sigma) > 0$ always.

Equation~\eqref{eq:ensemble} implies that the downward operator $\llbracket \cdot \rrbracket_\sigma$ maps $\mathcal{C}^\sigma$ to $\mathcal{C}$:

\begin{theorem}[Razborov~{\cite{razborov2007flag}}]\label{thm:downward-nonneg}
  For any flag type $\sigma$ and any $f \in \mathcal{A}^\sigma$ with $0 \leq_\sigma f$,
  \[
    0 \;\leq\; \llbracket f \rrbracket_\sigma.
  \]
\end{theorem}
\begin{proof}[Proof sketch]
  Let $\phi_0 \in \poshom{\emptyset}$; we show $\phi_0(\llbracket f \rrbracket_\sigma) \geq 0$.
  If $\phi_0(\sigma) = 0$, then the flag type $\sigma$ has zero density in the limiting object, and every $\sigma$-flag density also
  vanishes, so $\phi_0(\llbracket f \rrbracket_\sigma) = 0$.
  Otherwise, let $\boldsymbol{\phi}^\sigma \sim \Ext_\sigma(\phi_0)$.
  Since $0 \leq_\sigma f$, every positive homomorphism on
  $\mathcal{A}^\sigma$ evaluates $f$ non-negatively, so
  $\mathbb{E}[\boldsymbol{\phi}^\sigma(f)] \geq 0$.
  Rearranging~\eqref{eq:ensemble}:
  \[
    \phi_0\bigl(\llbracket f \rrbracket_\sigma\bigr)
    \;=\;
    \phi_0\bigl(\llbracket \mathbf{1}_\sigma \rrbracket_\sigma\bigr)
    \cdot
    \mathbb{E}\bigl[\boldsymbol{\phi}^\sigma(f)\bigr]
    \;\geq\; 0.
  \]
  Since $\phi_0$ was arbitrary, $0 \leq \llbracket f \rrbracket_\sigma$.
\end{proof}

A flag algebra proof of $P \leq c\cdot\mathbf{1}$ typically proceeds as follows:
first, one exhibits $f \in \mathcal{A}^\sigma$ with $0 \leq_\sigma f$ for some
non-empty flag type $\sigma$ (e.g., a square $f = g^2$) and applies
\Cref{thm:downward-nonneg} to obtain $0 \leq \llbracket f \rrbracket_\sigma$;
then, one adds $\llbracket f \rrbracket_\sigma$ to $P$ or an expansion of $P$
at an appropriate size, and uses non-negativity of flags and
\eqref{eq:sum-to-one} to conclude $P \leq c\cdot\mathbf{1}$.
The following example illustrates this pattern.

\begin{example}[Mantel's theorem via flag algebras]\label{ex:mantel-fa}
We prove $\fKtwo \leq \tfrac{1}{2}\cdot\mathbf{1}$ in the induced-triangle-free flag algebra (i.e., the algebra $\mathcal{A}^\emptyset$ with $\mathcal{H} = \{K_3\}$).
Let $\sigma$ be the one-vertex flag type.

\textit{Step 1: A key inequality.}
Since $0 \leq_\sigma (\fbarKtwobull - \fKtwobull)^2$,
\Cref{thm:downward-nonneg} gives $0 \leq \llbracket(\fbarKtwobull - \fKtwobull)^2\rrbracket_\sigma$.
Expanding the square and computing each product at size~$3$,
then applying the downward operator:
\begin{align}
  0 &\leq \llbracket(\fbarKtwobull - \fKtwobull)^2\rrbracket_\sigma\notag\\
  &= \llbracket\fbarKtwobull\cdot\fbarKtwobull - 2\cdot\fbarKtwobull\cdot\fKtwobull + \fKtwobull\cdot\fKtwobull\rrbracket_\sigma\notag\\
  &= \left\llbracket\fbarKthreebull + \fEthreebull
    - \fbarPthreebull - \fPthreebull
    + \fKonetwobull\right\rrbracket_\sigma\notag\\
  &= \fbarKthree - \frac{1}{3}\,\fbarPthree - \frac{1}{3}\,\fPthree.
  \label{eq:mantel-key}
\end{align}

\textit{Step 2: Algebraic derivation.}
Starting from \eqref{eq:k2-expansion}, we add two non-negative terms and simplify using \eqref{eq:sum-to-one}:
\begin{align*}
  \fKtwo
  &= \frac{1}{3}\,\fbarPthree + \frac{2}{3}\,\fPthree
  &&\text{by \eqref{eq:k2-expansion}}\\
  &\leq \frac{1}{3}\,\fbarPthree + \frac{2}{3}\,\fPthree +
    \frac{1}{2}\left(\fbarKthree - \frac{1}{3}\,\fbarPthree - \frac{1}{3}\,\fPthree\right) +
    \frac{1}{3}\,\fbarPthree
  &&\text{by \eqref{eq:mantel-key} and $0 \leq \fbarPthree$}\\
  &= \frac{1}{2}\cdot\left(\fbarKthree + \fbarPthree + \fPthree\right)
  \;=\; \frac{1}{2}\cdot\mathbf{1}.
  &&\text{by \eqref{eq:sum-to-one}}
\end{align*}
\end{example}

% \begin{corollary}[Razborov~{\cite{razborov2007flag}}]
%   \label{thm:sdp-nonneg}
%   For any flag type $\sigma$, elements $e_1,\ldots,e_r \in \mathcal{A}^\sigma$,
%   and positive semidefinite matrix $Q \in \mathbb{R}^{r\times r}$,
%   \[
%     \Bigl\llbracket \sum_{i,j=1}^r Q_{ij}\, e_i\, e_j \Bigr\rrbracket_\sigma
%     \;\in\; \mathcal{C}.
%   \]
% \end{corollary}
% \begin{proof}
%   PSD quadratic forms are in $\mathcal{C}^\sigma$ by the squares argument
%   (\Cref{sec:background}), so \Cref{thm:downward-nonneg} applies directly.
% \end{proof}
%
% The corollary says that every positive homomorphism on
% $\mathcal{A}^\emptyset$ evaluates the downward-averaged quadratic expression
% as a non-negative real number.  This is the engine behind flag algebra bounds
% based on positive semidefinite matrices.  Given a target element
% $x_0 \in \mathcal{A}^\emptyset$ and a claimed bound $c \in \mathbb{R}$, if
% one exhibits a PSD matrix $Q$ and elements $e_i \in \mathcal{A}^\sigma$ such
% that
% \[
%   c \cdot \mathbf{1} - x_0
%   \;=\;
%   \Bigl\llbracket \sum_{i,j} Q_{ij}\, e_i\, e_j \Bigr\rrbracket_\sigma,
% \]
% then $c\cdot\mathbf{1}-x_0\in\mathcal{C}$, so
% $\phi(x_0) \leq c$ for every positive homomorphism
% $\phi:\mathcal{A}^\emptyset\to\mathbb{R}$.

% \subsection{Turán Densities and the Pentagon Problem}
%
% We can now connect the algebra back to the combinatorial problems that motivate
% it.  Here $H$ is an unlabeled forbidden graph and $P$ is an unlabeled target
% graph.  The finite graphs being optimized over are $H$-free, meaning that they
% contain no induced subgraph isomorphic to $H$, and the target statistic is the
% induced density of $P$.
%
% In the algebraic certificate used in this paper, the forbidden-graph
% assumption is represented as a condition on positive homomorphisms.  Being
% $H$-free is equivalent to having zero induced density of $H$, so we restrict
% attention to positive homomorphisms $\phi:\mathcal{A}^\emptyset\to\mathbb{R}$
% satisfying
% \[
%   \phi([H])=0.
% \]
% Every convergent sequence of $H$-free graphs gives such a homomorphism.  For
% $H=K_3$, the condition says precisely that the limiting triangle density is
% zero.
%
% To distinguish this from the type-indexed order $\leq_\sigma$, for
% $f,g\in\mathcal{A}^\emptyset$ write
% \[
%   f \leq g
% \]
% when $\phi(f)\leq\phi(g)$ for every positive homomorphism
% $\phi:\mathcal{A}^\emptyset\to\mathbb{R}$ satisfying
% $\phi([H])=0$.  Equivalently, $g-f$ evaluates to a non-negative real number in
% every positive homomorphism satisfying the forbidden-graph condition
% $\phi([H])=0$.
% Thus an upper-bound proof for the target graph $P$ aims to establish
% \[
%   [P] \leq c\cdot\mathbf{1}.
% \]
% The finite proof data usually combines three ingredients already defined
% above: zero-space relations expand expressions to a common size, PSD quadratic
% forms provide semantically non-negative terms, and the downward operator turns
% labeled quadratic terms into unlabeled ones.
%
% \paragraph{Generalized Turán density.}
% For an unlabeled target graph $P$ and an unlabeled forbidden graph $H$, the
% \emph{generalized Turán density} is
% \[
%   \tdensity{P}{H}
%   \;:=\;
%   \lim_{n\to\infty}
%   \frac{1}{\binom{n}{|V(P)|}} \cdot
%     {\max\Bigl\{|\text{induced copies of }P\text{ in }G|
%                    ~\Big|~ |V(G)|=n \text{ and } G \text{ is }H\text{-free}\Bigr\}}.
% \]
% For the complete forbidden graph used in the certificate, the semantic
% inequality $[P] \leq c\cdot\mathbf{1}$ implies
% $\tdensity{P}{H}\leq c$.  A matching construction gives the reverse
% inequality when the claimed value is sharp.
%
% \paragraph{The Erd\H{o}s pentagon problem.}
% The problem asks for $\tdensity{C_5}{K_3}$: the maximum density of
% pentagons in triangle-free graphs.  The answer $24/625$ was proved
% independently by Grzesik~\cite{grzesik2012} and by Hatami, Hladk\'y, Kr\'al,
% Norine, and Razborov~\cite{hatami2012}.  It is achieved by the blow-up of
% $C_5$ (partition $n$ vertices into five nearly-equal groups and add all edges
% between consecutive groups in the cycle).  We use this theorem as the main
% case study throughout the paper, as it involves the full proof-engineering
% stack: abstract structure, extensional finite counting, data-heavy density
% computation (thousands of rational values over 5-vertex graphs), and heavy
% algebraic bookkeeping.


\section{Formalizing Flag Algebras in Lean~4}
\label{sec:abstract}

% -- Original introduction (commented out) ------------------------------------
% The specification layer of our Lean formalization gives flag algebras a
% typed mathematical interface that is reusable across different extremal
% problems.  This section describes the choices we made for that interface.
% We are reporting choices, not contributions: each choice has a coherent
% alternative, and on at least one of them we do not know which is preferable.
%
% \paragraph{A design choice, not an enabling condition.}
% In our predicate-based design, the forbidden-subgraph assumption is deferred
% to a parameter on the semantic order rather than baked into a universal theory
% at the outset.  In Razborov's original presentation~\cite{razborov2007flag},
% the forbidden subgraph is specified by a universal first-order theory chosen
% when the algebra is set up; both designs can express the same mathematical
% statements.  A consequence---not an enabling condition---of our choice is
% that random-extension theorems become the definitional content of the
% forbidden-subgraph relation, rather than serving only as an analytic
% ingredient inside a fixed theory.  What differs between the two designs is
% which intermediate lemmas are reusable across choices of constraint, and how
% many compatibility theorems must be reproved when a new forbidden graph is
% introduced.  Whether predicate-deferral is preferable to Razborov's
% universal-theory specification is a tradeoff we report on, not one we resolve.
%
% \paragraph{What this section reports.}
% In informal arguments, the same symbol freely denotes a concrete graph, its
% isomorphism class, a formal linear combination of flags, an element of the
% quotient algebra, or a real value assigned by a positive homomorphism; the
% transitions between these roles are implicit.  In Lean they are explicit:
% each role is a distinct type, and each transition is a map that must be
% proved compatible with the relevant equivalence relations and algebraic
% operations.  The specification layer provides those maps and their
% compatibility theorems, so that every subsequent certificate proof can work
% at the algebra level without re-establishing the soundness of the surrounding
% framework.
%
% The structure of the specification follows Razborov's mathematical
% presentation.  Flags are quotients of labeled graphs by isomorphism.  Their
% densities are formalized as a sampling semantics: the probability that a
% uniformly random induced subgraph of a host graph is isomorphic to a given
% flag.  The flag algebra is a further quotient, of the free real vector space
% on flags, by expansion laws that identify a small flag with its expected
% density expansion to a collection of larger flags.  Positive homomorphisms
% are the models of that equational theory.  Razborov's random-extension
% theorem, formalized using Mathlib's Prokhorov compactness theorem, connects
% the standard empty-typed level with the typed levels: every positive
% homomorphism $\phi$ on the empty-typed algebra extends to a probability
% measure $\Ext_\sigma(\phi)$ on the space of $\sigma$-typed positive
% homomorphisms for any $\sigma$.  This measure-theoretic result is the
% foundation on which the forbidden-subgraph predicate $\leq_\sigma$
% rests.  The transfer theorem
% \lean{generalizedTuranDensity\_le\_of\_forbidLE} then reduces a Tur\'an
% density bound to a single empty-typed semantic inequality $[P] \leq
% c\cdot\mathbf{1}$ in the ambient algebra.  Any proof of that inequality,
% however it is assembled, immediately yields the bound $\tdensity{P}{H}\leq
% c$; and all the abstract-layer results about products, quadratic forms, and
% the downward operator carry over to every forbidden graph without reproof.
% Code snippets are lightly rendered in ASCII when the Lean source uses Unicode
% notation; the names and fields match the implementation.
% -- End original introduction -------------------------------------------------

Informal flag algebra proofs are compact partly because the notation is
deliberately overloaded.  The same symbol may denote a concrete labeled graph,
its isomorphism class, a basis vector in a free vector space, an element of the
quotient algebra, or the real number obtained by evaluating that element in a
limit.  On paper these shifts are usually harmless and are left implicit.  In
Lean they become the first formalization problem: each role must be represented
by a separate type, and each shift between roles must be justified by a theorem
saying that it respects isomorphism, quotienting, density expansion, and
multiplication.

Our first design choice is to separate this typed mathematical interface from
the proof-by-computation machinery used later to check flag-algebra
certificates.  Here computation means evaluation inside Lean as part of proof
checking, not a numerical experiment outside the prover.  For example, Lean may
evaluate a finite flag-density table or normalize an algebraic expression, and
the resulting equality can be used as a proof step only because previously
proved theorems connect the computation to the mathematical definitions.  We
call the interface supporting these proof-producing computations the
\emph{specification layer}.  It formalizes the objects from
Section~\ref{sec:background}---flags as isomorphism classes of labeled graphs,
densities as finite sampling probabilities, the flag algebra as a quotient by
expansion identities, and positive homomorphisms as semantic interpretations of
algebra elements---and proves their basic interface theorems.  Once this layer
has been established, certificate proofs can manipulate algebraic expressions
without repeatedly returning to graph representatives and finite sampling
definitions.

A second design choice concerns where forbidden-subgraph assumptions enter the
development.  In the background section we fixed $\mathcal{H}$ first and
defined all flags and algebras relative to the induced-\(\mathcal{H}\)-free
world.  In Lean we instead build the ambient flag algebra without fixing
$\mathcal{H}$, and impose induced-\(\mathcal{H}\)-freeness later as a
predicate on the semantic order.  Thus the core formalization implements the unconstrained
theory.  If every definition carried an arbitrary forbidden family
$\mathcal{H}$ from the start, then even the elementary constructions of flags,
densities, expansion identities, and multiplication would acquire an extra
parameter and extra side conditions.  By postponing induced-\(\mathcal{H}\)-freeness, we
develop the ambient algebra once and introduce the constraint only when it is
needed for extremal applications.  This keeps the algebraic infrastructure
reusable across different forbidden graphs.  The deferred treatment is still
sound for induced Tur\'an-density bounds: \Cref{sec:forbidden} defines the
ensemble semantic order used for forbidden graphs and states the transfer
theorem that turns an empty-type ensemble inequality into the corresponding
induced extremal-density bound.

One notational convention is specific to the formalization.  In
Section~\ref{sec:background}, the empty flag type was written simply as
$\emptyset$.  In Lean-related notation, we write it as $\emptyt$ to avoid
confusion with ordinary empty sets.

The displayed snippets below are also an expository version rather than a
literal transcript of the source.  In the actual Lean implementation, many of
the parameters are declared once as ambient variables and then inferred by
Lean; here we spell them out at the definitions where they are used.  When a
displayed definition marks a parameter as implicit, later applications omit
that parameter, following Lean's usual argument-passing convention.  We also
occasionally use paper-facing names for intermediate terms when the
implementation names are optimized for local proof context or tactic plumbing.
These deviations are meant to expose the dependency of each construction on
the label type, the vertex types, and the flag type \(\sigma\), without asking
the reader to remember a global Lean context or source-file naming convention.

\subsection{Flag Types and Flags}
\label{sec:flagtypes}

Our formalization builds on Mathlib, and in particular on Mathlib's graph
library.  Before giving the Lean definitions of flags, we recall the small
pieces of Mathlib notation used below.  In Section~\ref{sec:background} we
described a finite graph as having a vertex set $V(G)$.  In Lean, the vertex
set is represented by a type: a term of type \lean{SimpleGraph V} is a simple
graph whose vertices are the elements of \lean{V}.  Finiteness is supplied
separately by assumptions such as \lean{Fintype V}; for example,
\lean{SimpleGraph (Fin n)} is the Lean version of a graph on the vertex set
\(\{0,\ldots,n-1\}\).

Mathlib provides ${G\hookrightarrow_g H}$ for the type of graph embeddings
from $G$ to $H$.  Such an embedding is an injective vertex map that
preserves and reflects adjacency, so its image is an induced copy of $G$
in $H$.  Mathlib also provides $G\simeq_g H$ for graph isomorphisms.
With this notation in place, the concrete pair $(M,\theta_A)$ from
Section~\ref{sec:background} is represented by \lean{LabeledGraph}, while its
isomorphism class is represented by \lean{Flag}.

\begin{lstlisting}
abbrev FlagType (T : Type) := SimpleGraph T

structure LabeledGraph {T : Type} [Fintype T] ((*@$\sigma$@*) : FlagType T) (V : Type) where
  graph      : SimpleGraph V
  type_embed : (*@$\sigma$@*) (*@$\hookrightarrow$@*)g graph
\end{lstlisting}
Here \lean{FlagType T} is simply a name for \lean{SimpleGraph T}: a flag type
is represented by a simple graph on the label type \lean{T}.  The structure
\lean{LabeledGraph} packages the two pieces of a concrete $\sigma$-flag: the
field \lean{graph} stores the underlying graph $M$, and the field
\lean{type\_embed} stores the embedding
$\theta_A:\sigma\hookrightarrow_g M$.  In the signature, braces mark
parameters that Lean normally infers from context, and square brackets mark
typeclass assumptions.  Thus \lean{[Fintype T]} tells Lean that the label type
\lean{T} is finite; later signatures use assumptions such as
\lean{[DecidableEq V]} when equality on a vertex type must be decidable.

We define \lean{LabeledGraphIso} as a structure capturing isomorphisms between
two \lean{LabeledGraph} terms, and \lean{Flag} as the type obtained by
quotienting \lean{LabeledGraph} by this isomorphism relation:
\begin{lstlisting}
structure LabeledGraphIso {T V V' : Type} [Fintype T]
    ((*@$\sigma$@*) : FlagType T) (G : LabeledGraph (*@$\sigma$@*) V) (G' : LabeledGraph (*@$\sigma$@*) V') where
  graph_iso     : G.graph (*@$\simeq$@*)g G'.graph
  type_preserve : graph_iso (*@$\circ$@*) G.type_embed = G'.type_embed

infixl:50 " (*@$\simeq$@*)f " => LabeledGraphIso

instance labeledGraphSetoid {T : Type} [Fintype T] ((*@$\sigma$@*) : FlagType T) (V : Type)
    : Setoid (LabeledGraph (*@$\sigma$@*) V) where
  -- r : LabeledGraph (*@$\sigma$@*) V (*@$\to$@*) LabeledGraph (*@$\sigma$@*) V (*@$\to$@*) Prop
  r     := fun G G' => Nonempty (G (*@$\simeq$@*)f G')
  iseqv := ... -- proof that r defines an equivalence relation

def Flag {T : Type} [Fintype T] ((*@$\sigma$@*) : FlagType T) (V : Type) : Type :=
  Quotient (labeledGraphSetoid (*@$\sigma$@*) V)
\end{lstlisting}
\lean{Quotient} comes with bracket notation for equivalence classes.  In Lean
code, the bare notation \(\llbracket x\rrbracket\) means ``the quotient class
of the representative \(x\),'' with the relevant quotient inferred from the
expected type.  Thus, when the expected type is
\lean{Flag}~\(\sigma\)~\lean{V}, \(\llbracket G\rrbracket\) is the flag
represented by the concrete labeled graph \(G\).  Later, when the expected type
is \lean{FlagAlgebra}~\(\sigma\), the same bare brackets send a
\lean{FlagVector} representative to its class modulo the zero space.  This
bare quotient notation should be distinguished from the subscripted brackets
\(\llbracket-\rrbracket_\sigma\) or \(\llbracket-\rrbracket_0\), which denote
the downward operator.

\lean{LabeledGraphIso} represents an isomorphism of concrete $\sigma$-flags.
It consists of a graph isomorphism \lean{graph\_iso} between the underlying
graphs, together with a compatibility proof \lean{type\_preserve} saying that
this isomorphism sends the labeled copy of $\sigma$ in $G$ to the labeled copy
of $\sigma$ in $G'$.  This is the Lean form of the condition
$\alpha\circ\theta_G=\theta_{G'}$.  The notation $G\simeq_f G'$ abbreviates
this structure, with $\sigma$ inferred from the two labeled graphs.

Lean's quotient type does not take an arbitrary relation directly; it expects
a \emph{setoid}, namely a type together with an explicitly supplied equivalence
relation.  The instance \lean{labeledGraphSetoid} provides this package for
\lean{LabeledGraph}~$\sigma$~\lean{V}.  Its field \lean{r} identifies two
terms $G,G'$ when there exists labeled-graph isomorphism data
$G\simeq_f G'$, and its field \lean{iseqv} contains the proofs that this
relation is reflexive, symmetric, and transitive.  Finally,
\lean{Flag}~$\sigma$~\lean{V} is defined as the quotient of
\lean{LabeledGraph}~$\sigma$~\lean{V} by that setoid, so its elements are
precisely isomorphism classes of concrete $\sigma$-flags on the vertex type
\lean{V}.

We can now define $\mathcal{F}^\sigma_n$ in Lean as
\lean{Flag}~$\sigma$~\lean{(Fin}~$n$\lean{)}.  The type
\lean{FinFlag}~$\sigma$ is the all-size union $\mathcal{F}^\sigma =
\bigcup_n \mathcal{F}^\sigma_n$, defined as a dependent pair type in Lean; each element is a pair
$(n, F)$ where $n$ records the size and $F$ is the $\sigma$-flag of that size.
\begin{lstlisting}
def FinFlag {n(*@$_0$@*) : (*@$\mathbb{N}$@*)} ((*@$\sigma$@*) : FlagType (Fin n(*@$_0$@*))) : Type :=
  (*@$\Sigma$@*) (n : (*@$\mathbb{N}$@*)), Flag (*@$\sigma$@*) (Fin n)
\end{lstlisting}

Note that unlike Section~\ref{sec:background}, none of the definitions above carries the
induced-\(\mathcal{H}\)-free condition: the underlying graphs are arbitrary.  This is a
deliberate interface choice.  The core quotient and density machinery is
developed once for all finite graphs, while the induced-\(\mathcal{H}\)-free assumption
is imposed later through the semantic predicate on the final unlabeled
inequality.

\subsection{Subflag Densities}
\label{sec:formal-subflag-densities}

Formalizing the density function is not straightforward.  The density
$\den{F_1,\ldots,F_t}{G}$ from Section~\ref{sec:background} should be a
function taking $t$ pattern flags and a host flag to a value in $[0,1]$, but
\lean{Flag} is a quotient type of \lean{LabeledGraph}.  The standard way to
define a function from a quotient type in Lean is to first define it on the
underlying type, prove that it respects the equivalence relation, and lift it
via \lean{Quotient.lift}.  That is, one must first define the density with
\lean{LabeledGraph} inputs, and then prove that replacing any input by an
isomorphic \lean{LabeledGraph} gives the same result.

We first introduce \lean{LabeledGraphList} as an abbreviation for a
(\lean{Fin t})-indexed tuple of \lean{LabeledGraph}s.
Given such a tuple of pattern \lean{LabeledGraph}s and a host \lean{LabeledGraph},
\lean{labeledGraphListDensity} computes the density as the number of tuples
$(S_1, \ldots, S_t)$ satisfying the conditions in the definition of $\den{F_1,\ldots,F_t}{G}$,
divided by the total number of candidate tuples.  Rather than adopting the
probabilistic sampling formulation of Section~\ref{sec:background}, this direct ratio is easier to work with in Lean.
\begin{lstlisting}
abbrev LabeledGraphList {T : Type} [Fintype T]
    ((*@$\sigma$@*) : FlagType T) (t : (*@$\mathbb{N}$@*)) (Vl : Fin t (*@$\to$@*) Type) : Type :=
  (*@$\forall$@*) (i : Fin t), LabeledGraph (*@$\sigma$@*) (Vl i)

noncomputable def labeledGraphListDensity
    {T W : Type} [Fintype T] [Fintype W] [DecidableEq W]
    {t : (*@$\mathbb{N}$@*)} {Vl : Fin t (*@$\to$@*) Type} [FintypeList Vl] [DecidableEqList Vl]
    ((*@$\sigma$@*) : FlagType T) (Hl : LabeledGraphList (*@$\sigma$@*) t Vl) (G : LabeledGraph (*@$\sigma$@*) W) : (*@$\mathbb{Q}$@*) :=
  let r_list (i : Fin t) := (Hl i).size - (*@$\sigma$@*).size
  labeledGraphListCount Hl G / multinomialCoefficient r_list (G.size - (*@$\sigma$@*).size)
\end{lstlisting}
The first line defines \lean{LabeledGraphList} by spelling out what a finite
list means in this dependent setting: it is a function that assigns to each
index \lean{i : Fin t} a concrete $\sigma$-flag \lean{LabeledGraph}~$\sigma$~\lean{(Vl i)}.
The family \lean{Vl : Fin t -> Type} records the vertex type of each entry, so
the $t$ pattern flags are not forced to share one ambient vertex type.
The square-bracket assumptions in the density signature provide the finite
enumeration and equality tests needed for counting.  A \lean{Fintype} instance
for a vertex type gives Lean a finite list of all its vertices, and a
\lean{DecidableEq} instance gives a decision procedure for testing whether two
vertices are equal.  The typeclasses \lean{FintypeList Vl} and
\lean{DecidableEqList Vl} are componentwise versions of these assumptions:
they supply the corresponding enumeration and equality-test instances for
each vertex type \lean{Vl i} in the tuple.

The second definition is the finite ratio corresponding to
$\den{F_1,\ldots,F_t}{G}$.  Here \lean{Hl} is the tuple of pattern labeled
graphs and \lean{G} is the concrete labeled graph in which they are counted.
The local function \lean{r\_list} records, for each pattern, how many unlabeled
vertices it contributes:
\[
  r_i = |V(H_i)| - |\sigma|.
\]
The numerator \lean{labeledGraphListCount Hl G}, defined immediately before
this function in the Lean development, counts the induced subgraph lists of
\lean{G} that realize the patterns in \lean{Hl}.  The denominator applies
\lean{multinomialCoefficient} to \lean{r\_list} and
\(|V(G)| - |\sigma|\), written in Lean as the difference between \lean{G.size}
and the size of \(\sigma\); it counts the possible choices of
pairwise-disjoint vertex sets of those sizes among the unlabeled vertices of
\lean{G}.  Thus the quotient is a rational number.
The \lean{noncomputable} keyword signals that Lean cannot synthesize an
algorithm to evaluate this definition directly.  This poses no obstacle to
formalizing flag algebra itself, but becomes relevant when applying the algebra
to concrete extremal problems, a point we return to in Section~\ref{sec:reflection}.

The concrete function \lean{labeledGraphListDensity} has two inputs: the tuple
\lean{Hl} of pattern labeled graphs and the target labeled graph \lean{G}.  To
make this a function of flags, we lift these two inputs separately.  The target
input is lifted from \lean{LabeledGraph} to \lean{Flag}; for the pattern tuple,
the intermediate quotient is the quotient of concrete tuples by entrywise
isomorphism.  We first lift the target input \lean{G}, then lift the pattern
tuples, and finally package the result as a function on \lean{FlagList}.  The
congruence proofs, omitted below, show that the density is unchanged when the
target or any pattern flag is replaced by an isomorphic representative.
\begin{lstlisting}
def labeledGraphListEqv {T : Type} [Fintype T] {(*@$\sigma$@*) : FlagType T}
    {t : (*@$\mathbb{N}$@*)} {Vl : Fin t (*@$\to$@*) Type} (Gl Gl' : LabeledGraphList (*@$\sigma$@*) t Vl) : Prop :=
  (*@$\forall$@*) (i : Fin t), Nonempty (Gl i (*@$\simeq$@*)f Gl' i)

instance labeledGraphListSetoid {T : Type} [Fintype T]
    ((*@$\sigma$@*) : FlagType T) (t : (*@$\mathbb{N}$@*)) (Vl : Fin t (*@$\to$@*) Type)
    : Setoid (LabeledGraphList (*@$\sigma$@*) t Vl) where
  r     := labeledGraphListEqv
  iseqv := ... -- proof that the entrywise isomorphism is an equivalence

def QuotLabeledGraphList {T : Type} [Fintype T]
    ((*@$\sigma$@*) : FlagType T) (t : (*@$\mathbb{N}$@*)) (Vl : Fin t (*@$\to$@*) Type) : Type :=
  Quotient (labeledGraphListSetoid (*@$\sigma$@*) t Vl)

abbrev FlagList {T : Type} [Fintype T]
    ((*@$\sigma$@*) : FlagType T) (t : (*@$\mathbb{N}$@*)) (Vl : Fin t (*@$\to$@*) Type) : Type :=
  (*@$\forall$@*) (i : Fin t), Flag (*@$\sigma$@*) (Vl i)

noncomputable def labeledGraphListDensityLifted
    {T W : Type} [Fintype T] [Fintype W] [DecidableEq W]
    {t : (*@$\mathbb{N}$@*)} {Vl : Fin t (*@$\to$@*) Type} [FintypeList Vl] [DecidableEqList Vl]
    {(*@$\sigma$@*) : FlagType T} (Hl : LabeledGraphList (*@$\sigma$@*) t Vl)
    : Flag (*@$\sigma$@*) W (*@$\to$@*) (*@$\mathbb{Q}$@*) :=
  Quotient.lift
    (fun G => labeledGraphListDensity Hl G)
    (...) -- proof that target equivalence preserves density

noncomputable def quotLabeledGraphListDensity
    {T W : Type} [Fintype T] [Fintype W] [DecidableEq W]
    {t : (*@$\mathbb{N}$@*)} {Vl : Fin t (*@$\to$@*) Type} [FintypeList Vl] [DecidableEqList Vl]
    {(*@$\sigma$@*) : FlagType T}
    : QuotLabeledGraphList (*@$\sigma$@*) t Vl (*@$\to$@*) Flag (*@$\sigma$@*) W (*@$\to$@*) (*@$\mathbb{Q}$@*) :=
  Quotient.lift
    labeledGraphListDensityLifted
    (...) -- proof that pattern-list equivalence preserves density

noncomputable def FlagList.coe
    {T : Type} [Fintype T] {(*@$\sigma$@*) : FlagType T}
    {t : (*@$\mathbb{N}$@*)} {Vl : Fin t (*@$\to$@*) Type}
    (Fl : FlagList (*@$\sigma$@*) t Vl)
    : QuotLabeledGraphList (*@$\sigma$@*) t Vl :=
  (*@$\llbracket$@*)fun i => (Fl i).out(*@$\rrbracket$@*)

noncomputable def flagListDensity
    {T W : Type} [Fintype T] [Fintype W] [DecidableEq W]
    {t : (*@$\mathbb{N}$@*)} {Vl : Fin t (*@$\to$@*) Type} [FintypeList Vl] [DecidableEqList Vl]
    {(*@$\sigma$@*) : FlagType T}
    : FlagList (*@$\sigma$@*) t Vl (*@$\to$@*) Flag (*@$\sigma$@*) W (*@$\to$@*) (*@$\mathbb{Q}$@*) :=
  fun Fl => quotLabeledGraphListDensity (FlagList.coe Fl)
\end{lstlisting}
Here \lean{labeledGraphListEqv} says that two concrete tuples are equivalent exactly
when their entries are isomorphic as labeled graphs: for each index \lean{i},
there must exist an isomorphism \lean{Gl i}~\(\simeq_f\)~\lean{Gl' i}.  The
instance \lean{labeledGraphListSetoid} supplies the setoid required by
\lean{Quotient}.  Its quotient is \lean{QuotLabeledGraphList}.  The
abbreviation \lean{FlagList} is the more convenient tuple of already-quotiented
flags.

The three density definitions then lift the original concrete density one
argument at a time.  For a fixed concrete pattern tuple \lean{Hl},
\lean{labeledGraphListDensityLifted Hl} is a function
\lean{Flag}~$\sigma$~\lean{W}~\(\to \mathbb{Q}\), obtained by quotienting the
target argument.  The next definition, \lean{quotLabeledGraphListDensity},
lifts the pattern tuple as well, so its first input is a
\lean{QuotLabeledGraphList} and its second input is a target \lean{Flag}.  In
the final definition, the input \lean{Fl} has type
\lean{FlagList}~$\sigma$~\lean{t Vl}, that is, \lean{Fl i} is already a flag
for each \lean{i : Fin t}.  The term
\lean{FlagList.coe Fl} has type
\lean{QuotLabeledGraphList}~$\sigma$~\lean{t Vl}; concretely, it chooses a
representative \lean{(Fl i).out} for each flag in the tuple and forms the
quotient class of the concrete tuple \lean{fun i => (Fl i).out}.  Thus
\lean{flagListDensity} can call
\lean{quotLabeledGraphListDensity} after converting the
tuple-of-flags input into the quotient-of-tuples input expected by the
preceding definition.

The following \lean{flagDensity}$_1$ and \lean{flagDensity}$_2$ are specializations of \lean{flagListDensity}
to singleton and pair lists:
\begin{lstlisting}
noncomputable def flagDensity(*@$_1$@*)
    {T U W : Type} [Fintype T] [Fintype U] [DecidableEq U] [Fintype W] [DecidableEq W]
    {(*@$\sigma$@*) : FlagType T} (F : Flag (*@$\sigma$@*) U) (G : Flag (*@$\sigma$@*) W) : (*@$\mathbb{Q}$@*) :=
  flagListDensity [F](*@$^\mathtt{f}$@*) G

noncomputable def flagDensity(*@$_2$@*)
    {T U(*@$_1$@*) U(*@$_2$@*) W : Type} [Fintype T] [Fintype U(*@$_1$@*)] [DecidableEq U(*@$_1$@*)] [Fintype U(*@$_2$@*)] [DecidableEq U(*@$_2$@*)]
    [Fintype W] [DecidableEq W]
    {(*@$\sigma$@*) : FlagType T} (F(*@$_1$@*) : Flag (*@$\sigma$@*) U(*@$_1$@*)) (F(*@$_2$@*) : Flag (*@$\sigma$@*) U(*@$_2$@*)) (G : Flag (*@$\sigma$@*) W) : (*@$\mathbb{Q}$@*) :=
  flagListDensity [F(*@$_1$@*), F(*@$_2$@*)](*@$^\mathtt{f}$@*) G
\end{lstlisting}
Here the superscripted bracket notation is a tuple constructor for
\lean{FlagList}.  Thus \lean{[F]}$^{\mathtt{f}}$ is the one-entry
\lean{FlagList} whose unique entry is \lean{F}, and
\lean{[F}$_1$\lean{, F}$_2$\lean{]}$^{\mathtt{f}}$ is the two-entry
\lean{FlagList} sending the first index of \lean{Fin 2} to \lean{F}$_1$ and
the second to \lean{F}$_2$.  The superscript \(\mathtt{f}\) distinguishes this
notation from ordinary list notation and from quotient brackets; it does not
add another quotienting operation.

Building on the definitions above, we prove the following chain rule (Lemma~\ref{lem:chain-rule}) and several of its variants.
Although the statement mirrors the elementary sampling identity from
Section~\ref{sec:background}, the formal proof must also commute finite counts
with quotient lifts and the coercions between rational densities and real-valued
algebra expressions.
\begin{lstlisting}
theorem flagDensity_eq_sum_density_prods
    {(*@$\ell_0$@*) (*@$\ell_1$@*) (*@$\ell$@*) (*@$\ell'$@*) : (*@$\mathbb{N}$@*)}
    ((*@$\sigma$@*) : FlagType (Fin (*@$\ell_0$@*))) (F(*@$_1$@*) : Flag (*@$\sigma$@*) (Fin (*@$\ell_1$@*))) (G : Flag (*@$\sigma$@*) (Fin (*@$\ell$@*)))
    (h(*@$\ell_1$@*) : (*@$\ell_0$@*) (*@$\leq$@*) (*@$\ell_1$@*)) (h(*@$\ell'$@*) : (*@$\ell_1$@*) (*@$\leq$@*) (*@$\ell'$@*)) (h(*@$\ell$@*) : (*@$\ell'$@*) (*@$\leq$@*) (*@$\ell$@*))
    : flagDensity(*@$_1$@*) F(*@$_1$@*) G = (*@$\sum$@*) G' : Flag (*@$\sigma$@*) (Fin (*@$\ell'$@*)), flagDensity(*@$_1$@*) F(*@$_1$@*) G' * flagDensity(*@$_1$@*) G' G := ...
\end{lstlisting}

\subsection{The Flag Algebra}

To define the flag algebra, we first construct the free $\mathbb{R}$-module with basis
$\mathcal{F}^\sigma$.  Each element in this module is represented as a finitely supported
function from $\mathcal{F}^\sigma$ to $\mathbb{R}$, using Mathlib's \lean{Finsupp}
type.  \lean{FlagVector} $\sigma$ is the resulting type, and \lean{basisVector} is
the constructor that produces the basis vector for each flag:
\begin{lstlisting}
abbrev FlagVector {n(*@$_0$@*) : (*@$\mathbb{N}$@*)} ((*@$\sigma$@*) : FlagType (Fin n(*@$_0$@*))) : Type :=
  FinFlag (*@$\sigma$@*) (*@$\to_0$@*) (*@$\mathbb{R}$@*)

def basisVector {n(*@$_0$@*) : (*@$\mathbb{N}$@*)} {(*@$\sigma$@*) : FlagType (Fin n(*@$_0$@*))} (F : FinFlag (*@$\sigma$@*)) : FlagVector (*@$\sigma$@*) :=
  Finsupp.single F 1
\end{lstlisting}
Here $\to_0$ denotes the type of finitely supported functions, and \lean{Finsupp.single F 1} constructs the basis vector for $F \in \mathcal{F}^\sigma$: the function taking value $1$ at $F$ and $0$ elsewhere.

Next, we define the zero space $\mathcal{Z}^\sigma$.  The auxiliary function
\lean{flagExpansion F}~$\ell$ is the Lean version of
$\sum_{G \in \mathcal{F}^\sigma_\ell} \den{F}{G}\cdot G$, the expansion of
$F$ at size $\ell$.  We encode the generating element of
Equation~\eqref{eq:zero-element} as \lean{zeroElement}, \lean{zeroSet} $\sigma$
collects all elements of this form, and \lean{ZeroSpace} $\sigma$ is the
submodule they span:
\begin{lstlisting}
noncomputable def flagExpansion {n(*@$_0$@*) : (*@$\mathbb{N}$@*)} {(*@$\sigma$@*) : FlagType (Fin n(*@$_0$@*))} (F : FinFlag (*@$\sigma$@*)) ((*@$\ell$@*) : (*@$\mathbb{N}$@*))
    : FlagVector (*@$\sigma$@*) :=
  (*@$\sum$@*) F' : Flag (*@$\sigma$@*) (Fin (*@$\ell$@*)), flagDensity(*@$_1$@*) F.2 F' (*@$\leansmul$@*) basisVector (*@$\langle$@*)(*@$\ell$@*), F'(*@$\rangle$@*)

noncomputable def zeroElement {n(*@$_0$@*) : (*@$\mathbb{N}$@*)} {(*@$\sigma$@*) : FlagType (Fin n(*@$_0$@*))} (F : FinFlag (*@$\sigma$@*)) ((*@$\ell$@*) : (*@$\mathbb{N}$@*))
    : FlagVector (*@$\sigma$@*) :=
  basisVector F - flagExpansion F (*@$\ell$@*)

noncomputable def zeroSet {n(*@$_0$@*) : (*@$\mathbb{N}$@*)} ((*@$\sigma$@*) : FlagType (Fin n(*@$_0$@*)))
    : Set (FlagVector (*@$\sigma$@*)) :=
  {k | (*@$\exists$@*) (F : FinFlag (*@$\sigma$@*)) ((*@$\ell$@*) : (*@$\mathbb{N}$@*)), F.1 (*@$\leq$@*) (*@$\ell$@*) (*@$\wedge$@*) k = zeroElement F (*@$\ell$@*)}

noncomputable def ZeroSpace {n(*@$_0$@*) : (*@$\mathbb{N}$@*)} ((*@$\sigma$@*) : FlagType (Fin n(*@$_0$@*)))
    : Submodule (*@$\mathbb{R}$@*) (FlagVector (*@$\sigma$@*)) :=
  Submodule.span (*@$\mathbb{R}$@*) (zeroSet (*@$\sigma$@*))
\end{lstlisting}

Following the same pattern as the definition of \lean{Flag} from
\lean{LabeledGraph} in Section~\ref{sec:flagtypes}, we now define
the flag algebra $\mathcal{A}^\sigma$ as a quotient of \lean{FlagVector} $\sigma$.
The equivalence relation \lean{flagVectorEqv} declares two elements of \lean{FlagVector} $\sigma$ equivalent
when their difference lies in \lean{ZeroSpace} $\sigma$, and \lean{flagVectorSetoid} $\sigma$
packages this relation into a setoid so that Lean's quotient type can be applied.
\begin{lstlisting}
def flagVectorEqv {n(*@$_0$@*) : (*@$\mathbb{N}$@*)} {(*@$\sigma$@*) : FlagType (Fin n(*@$_0$@*))} (f g : FlagVector (*@$\sigma$@*)) : Prop :=
  f - g (*@$\in$@*) ZeroSpace (*@$\sigma$@*)

instance flagVectorSetoid {n(*@$_0$@*) : (*@$\mathbb{N}$@*)} ((*@$\sigma$@*) : FlagType (Fin n(*@$_0$@*))) : Setoid (FlagVector (*@$\sigma$@*)) where
  r     := flagVectorEqv
  iseqv := ... -- proof that flagVectorEqv is an equivalence relation

abbrev FlagAlgebra {n(*@$_0$@*) : (*@$\mathbb{N}$@*)} ((*@$\sigma$@*) : FlagType (Fin n(*@$_0$@*))) : Type :=
  Quotient (flagVectorSetoid (*@$\sigma$@*))
\end{lstlisting}
The quotient therefore imposes a family of expansion identities: for every
flag \lean{F : FinFlag}~$\sigma$ and every size $\ell \geq F.1$, the class of
\lean{basisVector F} is identified with the class of
\lean{flagExpansion F}~$\ell$, i.e. with the linear combination of all
size-$\ell$ flags weighted by their densities over \(F\).

Formalizing multiplication in $\mathcal{A}^\sigma$ also requires several steps.
The term \lean{flagMulWithSize} defines the product
of two flags at a chosen output size $\ell$ as in
Equation~\eqref{eq:mul}. The term \lean{flagMul} specializes this product to the minimal admissible
size $\ell = F.1 + F'.1 - n_0$. Then, \lean{bilinearExtension} extends
\lean{flagMul} linearly in each argument to give a product on
\lean{FlagVector} $\sigma$.  This intermediate product is only a bilinear
operation on formal sums; the quotient is where the usual algebra laws are
established.  Finally, \lean{Quotient.map$_2$} lifts the product to
\lean{FlagAlgebra} $\sigma$.
\begin{lstlisting}
noncomputable def flagMulWithSize {n(*@$_0$@*) : (*@$\mathbb{N}$@*)} {(*@$\sigma$@*) : FlagType (Fin n(*@$_0$@*))} (F F' : FinFlag (*@$\sigma$@*)) ((*@$\ell$@*) : (*@$\mathbb{N}$@*))
    : FlagVector (*@$\sigma$@*) :=
  (*@$\sum$@*) G : Flag (*@$\sigma$@*) (Fin (*@$\ell$@*)), flagDensity(*@$_2$@*) F.2 F'.2 G (*@$\leansmul$@*) basisVector (*@$\langle$@*)(*@$\ell$@*), G(*@$\rangle$@*)

noncomputable def flagMul {n(*@$_0$@*) : (*@$\mathbb{N}$@*)} {(*@$\sigma$@*) : FlagType (Fin n(*@$_0$@*))} (F F' : FinFlag (*@$\sigma$@*)) : FlagVector (*@$\sigma$@*) :=
  flagMulWithSize F F' (F.1 + F'.1 - n(*@$_0$@*))

noncomputable instance {n(*@$_0$@*) : (*@$\mathbb{N}$@*)} ((*@$\sigma$@*) : FlagType (Fin n(*@$_0$@*))) : Mul (FlagVector (*@$\sigma$@*)) where
  mul := bilinearExtension flagMul

noncomputable instance {n(*@$_0$@*) : (*@$\mathbb{N}$@*)} ((*@$\sigma$@*) : FlagType (Fin n(*@$_0$@*))) : Mul (FlagAlgebra (*@$\sigma$@*)) where
  mul := by
    apply Quotient.map(*@$_2$@*) ((*@$\cdot$@*) * (*@$\cdot$@*))
    -- proof omitted: for f f' g g' : FlagVector (*@\textcolor{gray}{$n_0$}@*) (*@\textcolor{gray}{$\sigma$}@*),
    --   flagVectorEqv f f' and flagVectorEqv g g'
    --   imply flagVectorEqv (f * g) (f' * g')
\end{lstlisting}
Here \lean{Quotient.map$_2$} is a Lean built-in function that lifts a binary
operation to a quotient type given a proof that the operation respects the
equivalence relation.  In the displayed definition, the expression
\lean{($\cdot$ * $\cdot$)} is not a new operation: since its two arguments are
flag vectors, Lean resolves \lean{*} using the immediately preceding
\lean{Mul (FlagVector}~$\sigma$\lean{)} instance, whose multiplication is
\lean{bilinearExtension flagMul}.  The quotient instance then lifts exactly
this vector-level multiplication to \lean{FlagAlgebra} $\sigma$.  The proof
obligation amounts to showing that this multiplication respects
\lean{flagVectorEqv}, as noted in the comment above.  The load-bearing lemma
is that \lean{ZeroSpace} $\sigma$ is
closed under multiplication by arbitrary flag vectors:
\begin{lstlisting}
theorem flagVector_mul_zeroSpace {n(*@$_0$@*) : (*@$\mathbb{N}$@*)} ((*@$\sigma$@*) : FlagType (Fin n(*@$_0$@*)))
    (f : FlagVector (*@$\sigma$@*)) {k : FlagVector (*@$\sigma$@*)} (hk_zero : k (*@$\in$@*) ZeroSpace (*@$\sigma$@*))
    : f * k (*@$\in$@*) ZeroSpace (*@$\sigma$@*) := ...
\end{lstlisting}
Together with commutativity, this says that \lean{ZeroSpace} $\sigma$ is an
ideal in the algebraic sense: multiplying a zero-space element by any flag
vector still gives a zero-space element.  This is exactly the property needed
to define multiplication on the quotient.  If two representatives differ by a
zero-space element, then multiplying both by the same vector still gives
representatives that differ by a zero-space element, so the product does not
depend on the chosen representatives.  This is the compatibility condition
that allows the Lean definition to use \lean{Quotient.map$_2$}.  A separate
theorem handles the choice of output size:
\begin{lstlisting}
theorem flagMulWithSize_indep_on_size {n(*@$_0$@*) : (*@$\mathbb{N}$@*)} ((*@$\sigma$@*) : FlagType (Fin n(*@$_0$@*)))
    {F(*@$_1$@*) F(*@$_2$@*) : FinFlag (*@$\sigma$@*)} {(*@$\ell_1$@*) (*@$\ell_2$@*) : (*@$\mathbb{N}$@*)} ...
    : flagVectorEqv (flagMulWithSize F(*@$_1$@*) F(*@$_2$@*) (*@$\ell_1$@*)) (flagMulWithSize F(*@$_1$@*) F(*@$_2$@*) (*@$\ell_2$@*)) := ...
\end{lstlisting}
The ellipsis hides the admissibility assumptions on the two output sizes:
\(\ell_1\) and \(\ell_2\) must both be large enough to contain the two input
flags after identifying their common type \(\sigma\).  In Lean this is written
as
\[
  F_1.1 + F_2.1 \leq \ell_i + n_0
  \qquad (i=1,2),
\]
where \(F_j.1\) is the total number of vertices of \(F_j\), and \(n_0\) is the
number of labeled vertices in \(\sigma\).  Equivalently, the intended output
size is at least \(F_1.1 + F_2.1 - n_0\): the unlabeled vertices outside the
type are sampled disjointly in the two flags, while the \(n_0\) type vertices
are shared.  The theorem then says that any two such admissible output sizes
give equivalent flag vectors, justifying the choice of the minimal size in
\lean{flagMul}.

With these definitions in place, we prove that \lean{FlagAlgebra} $\sigma$ is
a commutative associative $\mathbb{R}$-algebra, completing the formalization of the algebraic structure of $\mathcal{A}^\sigma$.

\subsection{Positive Homomorphisms and Semantic Order}

This subsection describes how we encode the semantic side of flag algebras in
Lean: sequences of finite flags, their limiting density functions, positive
homomorphisms, and the order induced by testing against all positive
homomorphisms.  The convergence theorem from
Theorem~\ref{thm:convergent-hom} is the main mathematical bridge between finite
flag sequences and positive homomorphisms, but the formalization first has to
make each of these surrounding notions explicit.

\lean{FlagSeq}~$\sigma$ is the type of sequences of $\sigma$-flags indexed by
$\mathbb{N}$.  The predicate \lean{ConvergesTo} says that such a sequence has a
prescribed limiting density function.
\begin{lstlisting}
abbrev FlagSeq {n(*@$_0$@*) : (*@$\mathbb{N}$@*)} ((*@$\sigma$@*) : FlagType (Fin n(*@$_0$@*))) :=
  (*@$\mathbb{N}$@*) (*@$\to$@*) FinFlag (*@$\sigma$@*)

def ConvergesTo {n(*@$_0$@*) : (*@$\mathbb{N}$@*)} {(*@$\sigma$@*) : FlagType (Fin n(*@$_0$@*))} (s : FlagSeq (*@$\sigma$@*)) (a : FinFlag (*@$\sigma$@*) (*@$\to$@*) (*@$\mathbb{R}$@*)) : Prop :=
  Increases s (*@$\wedge$@*) Tendsto (flagDensitySeq s) atTop ((*@$\mathcal{N}$@*) a)
\end{lstlisting}
Here the candidate limit \(a\) is a real-valued function on finite
\(\sigma\)-flags.  The first conjunct, \lean{Increases s}, asserts that the
underlying sizes of the flags \(s(0),s(1),\ldots\) strictly increase.  The
second conjunct uses Lean's filter notation for convergence:
\lean{atTop} is the limit \(n\to\infty\) on \(\mathbb{N}\), and
\(\mathcal{N} a\) is the neighbourhood filter of the function \(a\).  The term
\lean{flagDensitySeq s} is the sequence of density functions associated to
\(s\); its \(n\)-th value sends a finite \(\sigma\)-flag \(F\) to
\(\den{F}{s(n)}\).  Thus the \lean{Tendsto} line is the Lean statement that,
for every finite \(\sigma\)-flag \(F\),
\[
  \lim_{n\to\infty}\den{F}{s(n)}=a(F).
\]

Mathematically, a positive homomorphism is an
\(\mathbb{R}\)-algebra homomorphism
\(\phi:\mathcal{A}^\sigma\to\mathbb{R}\) such that
\(\phi(F)\ge 0\) for every finite \(\sigma\)-flag \(F\), viewed as a basis
element of the flag algebra.  The Lean definition follows this sentence almost
literally:
\begin{lstlisting}
abbrev Hom {n(*@$_0$@*) : (*@$\mathbb{N}$@*)} ((*@$\sigma$@*) : FlagType (Fin n(*@$_0$@*))) :=
  FlagAlgebra (*@$\sigma$@*) (*@$\to_a$@*)[(*@$\mathbb{R}$@*)] (*@$\mathbb{R}$@*)

def PositiveHom {n(*@$_0$@*) : (*@$\mathbb{N}$@*)} ((*@$\sigma$@*) : FlagType (Fin n(*@$_0$@*))) : Type :=
  { (*@$\phi$@*) : Hom (*@$\sigma$@*) // (*@$\forall$@*) F : FinFlag (*@$\sigma$@*), (*@$\phi$@*) (*@$\llbracket$@*)basisVector F(*@$\rrbracket$@*) (*@$\geq$@*) 0 }
\end{lstlisting}
The arrow \lean{$\to_a$[$\mathbb{R}$]} in the first definition is Lean's
notation for the type of \(\mathbb{R}\)-algebra homomorphisms:
functions preserving addition, multiplication, scalar multiplication by real
numbers, and the unit.  Thus \lean{Hom} abbreviates the type of algebra
homomorphisms from \lean{FlagAlgebra} to \(\mathbb{R}\), with the flag type
\(\sigma\) supplied as its argument.  The second definition is Lean's subtype
construction.  Mathematically, we read it in the usual set-comprehension form
\[
  \{\phi\in \operatorname{Hom}(\mathcal{A}^\sigma,\mathbb{R})
    \mid \phi(F)\ge 0
    \text{ for every }F\in\mathcal{F}^{\sigma}\}.
\]
In other words, the only extra condition on an algebra homomorphism is the
positivity condition.  In Lean syntax, this is the
\lean{$\forall$ F : FinFlag $\sigma$} field after the subtype separator.
Here \lean{basisVector F} is the formal basis vector representing \(F\), and
the bare quotient brackets send that flag vector to its class in the flag
algebra before \(\phi\) evaluates it.

Using these definitions, we formalize Theorem~\ref{thm:convergent-hom} as two separate theorems:
\begin{lstlisting}
(*@\textcolor{gray}{\itshape -- Theorem~\ref{thm:convergent-hom} (a)}@*)
theorem flagSeq_limit_mem_positiveHom {n(*@$_0$@*) : (*@$\mathbb{N}$@*)} {(*@$\sigma$@*) : FlagType (Fin n(*@$_0$@*))}
    (s : FlagSeq (*@$\sigma$@*)) {a : FlagDensitySpace (*@$\sigma$@*)} (hs_conv : ConvergesTo s a)
    : (*@$\exists$@*) ((*@$\phi$@*) : PositiveHom (*@$\sigma$@*)), (*@$\phi$@*).coe = a := ...

(*@\textcolor{gray}{\itshape -- Theorem~\ref{thm:convergent-hom} (b)}@*)
theorem positiveHom_as_flagSeq_limit {n(*@$_0$@*) : (*@$\mathbb{N}$@*)} {(*@$\sigma$@*) : FlagType (Fin n(*@$_0$@*))}
    ((*@$\phi$@*) : PositiveHom (*@$\sigma$@*))
    : (*@$\exists$@*) (s : FlagSeq (*@$\sigma$@*)), ConvergesTo s (*@$\phi$@*).coe := ...
\end{lstlisting}
Here \lean{FlagDensitySpace} with argument \(\sigma\) is the set of functions
\(\mathcal{F}^\sigma\to\mathbb{R}\) whose values lie in \([0,1]\), so its
points are density profiles of finite \(\sigma\)-flags.  The expression
\lean{$\phi$.coe} denotes the map \lean{PositiveHom.coe}, which forgets the
algebra-homomorphism structure of \(\phi\) and remembers only the density
profile it induces:
\[
  F\longmapsto \phi\bigl(\llbracket\lean{basisVector F}\rrbracket\bigr).
\]
Thus the first theorem says that the limiting density profile \(a\) of a
convergent flag sequence is equal to this profile of some positive
homomorphism \(\phi\).  Conversely, the second theorem says that every
positive homomorphism \(\phi\), after passing to this profile, arises as the
limit of some flag sequence.
Both directions require nontrivial formalization work.  In part (a), one must
show that the limit \(a\) of a flag sequence \(s\) satisfies all algebraic
properties of a positive homomorphism, including details that informal proofs
leave implicit.  Part (b) goes in the opposite direction.  Starting from a
positive homomorphism \(\phi\), the proof uses the values \(\phi(G)\) to put a
probability distribution on the finite \(\sigma\)-flags of each sufficiently
large size: a flag \(G\) is sampled with weight \(\phi(G)\).  These weights are
indeed probabilities because \(\phi\) is non-negative on flags and satisfies the normalization
identity
\[
  \sum_{G\in\mathcal{F}^{\sigma}_{\ell}} \phi(G) 
  = 
  \phi\left(\sum_{G\in\mathcal{F}^{\sigma}_{\ell}} G\right) 
  = 
  \phi(\mathbf{1}_\sigma)
  = 
  1.
\]
One then samples an entire sequence
\[
  X=(X_0,X_1,\ldots),
\]
where \(X_n\) is chosen from the distribution on flags of size \(n^2+n_0\).
The key estimate says that, for each fixed test flag \(F\), the density
\(\den{F}{X_n}\) is very likely to be close to \(\phi(F)\) once \(n\) is large.
Thus almost every sampled sequence has the limiting density profile prescribed
by \(\phi\).  Selecting one such realization and calling it
\(G_0,G_1,\ldots\) gives the deterministic flag sequence required by
Theorem~\ref{thm:convergent-hom}(b).  The Lean proof formalizes this argument
by building the probability distributions, proving the concentration estimate,
and then extracting one sequence whose density profile is the profile obtained
from the Lean expression \lean{PositiveHom.coe}.

We also define the semantic non-negativity cone $\mathcal{C}^\sigma$
and the order it induces on $\mathcal{A}^\sigma$.
\lean{semanticCone}~$\sigma$ is the set of elements in
\lean{FlagAlgebra}~$\sigma$ that every positive homomorphism
evaluates non-negatively.  The induced order is installed as a Lean
\lean{LE} instance, with $f \leq g$ defined to mean
$g - f \in \mathcal{C}^\sigma$.
\begin{lstlisting}
def semanticCone {n(*@$_0$@*) : (*@$\mathbb{N}$@*)} ((*@$\sigma$@*) : FlagType (Fin n(*@$_0$@*))) : Set (FlagAlgebra (*@$\sigma$@*)) :=
  { f | (*@$\forall$@*) (*@$\phi$@*) : PositiveHom (*@$\sigma$@*), (*@$\phi$@*) f (*@$\geq$@*) 0 }

instance {n(*@$_0$@*) : (*@$\mathbb{N}$@*)} ((*@$\sigma$@*) : FlagType (Fin n(*@$_0$@*))) : LE (FlagAlgebra (*@$\sigma$@*)) where
  le := fun f g => g - f (*@$\in$@*) semanticCone (*@$\sigma$@*)
\end{lstlisting}

\subsection{The Downward Operator}

The downward operator from Section~\ref{sec:background-downward} forgets the
labels of a typed flag-algebra expression while keeping the correct averaging
factor.  On a single flag the intended formula is
\[
  \llbracket F \rrbracket_\sigma = q_\sigma(F)\cdot F|_\emptyset.
\]
The formalization follows this formula, but it has to pass through the two
quotients already introduced: first the quotient from concrete labeled graphs
to flags, and then the zero-space quotient from flag vectors to the flag
algebra.

The single-flag part defines the two ingredients of the formula separately.
The rational number \(q_\sigma(F)\) is implemented as
\lean{downwardNormalizingFactor}; before it can be used on flags, it is defined
on concrete \lean{LabeledGraph}s and proved invariant under labeled-graph
isomorphism.  The operation \(F|_\emptyset\), which forgets the type embedding,
is implemented first as \lean{unlabel\_labeledGraph} on concrete labeled graphs and
then lifted to \lean{unlabel} on \lean{Flag}s.  The map \lean{downwardFlag}
combines the two: it sends a typed flag to the basis vector of its unlabeling,
scaled by the normalizing factor.

The normalizing factor is computed by counting label placements.  For a
concrete typed labeled graph \(G\) on \(n\) vertices, \lean{isomorphismCount G}
is the number of ways to place the \(n_0\) labels on the same underlying
unlabeled graph so that the resulting typed labeled graph is flag-isomorphic to
\(G\).  The denominator \(n!/(n-n_0)!\) is the total number of injections of the
\(n_0\) labels into the \(n\) vertices, so their ratio is exactly
\(q_\sigma(G)\).  In the last line of the snippet, the black dot is Lean's
scalar multiplication operation on \lean{FlagVector}s: it scales the basis
vector indexed by the pair consisting of \(n\) and \lean{unlabel F} by
\lean{downwardNormalizingFactor F}.
\begin{lstlisting}
noncomputable def downwardNormalizingFactor_labeledGraph {n(*@$_0$@*) n : (*@$\mathbb{N}$@*)} {(*@$\sigma$@*) : FlagType (Fin n(*@$_0$@*))}
    (G : LabeledGraph (*@$\sigma$@*) (Fin n)) : (*@$\mathbb{Q}$@*) :=
  let num_of_all_injections := n.factorial / (n - n(*@$_0$@*)).factorial
  isomorphismCount G / num_of_all_injections

noncomputable def downwardNormalizingFactor {n(*@$_0$@*) n : (*@$\mathbb{N}$@*)} {(*@$\sigma$@*) : FlagType (Fin n(*@$_0$@*))}
    : Flag (*@$\sigma$@*) (Fin n) (*@$\to$@*) (*@$\mathbb{Q}$@*) :=
  Quotient.lift (fun G : LabeledGraph (*@$\sigma$@*) (Fin n) => downwardNormalizingFactor_labeledGraph G) (...)

def unlabel_labeledGraph {n(*@$_0$@*) : (*@$\mathbb{N}$@*)} {V : Type} {(*@$\sigma$@*) : FlagType (Fin n(*@$_0$@*))} (G : LabeledGraph (*@$\sigma$@*) V)
    : LabeledGraph (*@$\emptyt$@*) V := ...

noncomputable def unlabel {n(*@$_0$@*) : (*@$\mathbb{N}$@*)} {V : Type} {(*@$\sigma$@*) : FlagType (Fin n(*@$_0$@*))}
    : Flag (*@$\sigma$@*) V (*@$\to$@*) Flag (*@$\emptyt$@*) V :=
  Quotient.lift (fun G : LabeledGraph (*@$\sigma$@*) V => (*@$\llbracket$@*)unlabel_labeledGraph G(*@$\rrbracket$@*)) (...)

noncomputable def downwardFlag {n(*@$_0$@*) n : (*@$\mathbb{N}$@*)} {(*@$\sigma$@*) : FlagType (Fin n(*@$_0$@*))} (F : Flag (*@$\sigma$@*) (Fin n))
    : FlagVector (*@$\emptyt$@*) :=
  downwardNormalizingFactor F (*@$\leansmul$@*) basisVector (*@$\langle$@*)n, unlabel F(*@$\rangle$@*)
\end{lstlisting}

The algebra-level operator is built in three steps.  First,
\lean{downwardFlagVector} extends the single-flag map linearly from basis
flags to arbitrary \(\sigma\)-typed \lean{FlagVector} terms.  Second,
\lean{downwardFlagVectorQuot} sends the resulting empty-type flag vector into
\(\emptyt\)-typed \lean{FlagAlgebra}.  Third, \lean{downward} turns this
representative-level construction into a function whose input is an element of
the \(\sigma\)-typed quotient algebra, rather than a chosen
\lean{FlagVector} representative.

The compatibility theorem needed for the third step is the essential point.  A
flag vector may be changed by an element of the zero space without changing the
corresponding flag-algebra element; the respect lemma displayed below proves
that the downward image is unchanged whenever two flag-vector representatives
differ by an element of the zero space.  Thus the formalized operator is not
merely a linear map on representatives: it is a well-defined map on the
quotient algebra.
\begin{lstlisting}
noncomputable def downwardFlagVector {n(*@$_0$@*) : (*@$\mathbb{N}$@*)} {(*@$\sigma$@*) : FlagType (Fin n(*@$_0$@*))}
    : FlagVector (*@$\sigma$@*) (*@$\to$@*) FlagVector (*@$\emptyt$@*) :=
  linearExtension (fun F : FinFlag (*@$\sigma$@*) => downwardFlag F.2)

noncomputable def downwardFlagVectorQuot {n(*@$_0$@*) : (*@$\mathbb{N}$@*)} {(*@$\sigma$@*) : FlagType (Fin n(*@$_0$@*))}
    : FlagVector (*@$\sigma$@*) (*@$\to$@*) FlagAlgebra (*@$\emptyt$@*) :=
  fun f : FlagVector (*@$\sigma$@*) => (*@$\llbracket$@*)downwardFlagVector f(*@$\rrbracket$@*)

lemma downwardFlagVectorQuot_respect_eqv {n(*@$_0$@*) : (*@$\mathbb{N}$@*)} {(*@$\sigma$@*) : FlagType (Fin n(*@$_0$@*))}
    {f f' : FlagVector (*@$\sigma$@*)} (h : f - f' (*@$\in$@*) ZeroSpace (*@$\sigma$@*))
    : downwardFlagVectorQuot f = downwardFlagVectorQuot f' := ...

noncomputable def downward {n(*@$_0$@*) : (*@$\mathbb{N}$@*)} {(*@$\sigma$@*) : FlagType (Fin n(*@$_0$@*))}
    : FlagAlgebra (*@$\sigma$@*) (*@$\to$@*) FlagAlgebra (*@$\emptyt$@*) :=
  Quotient.lift (fun f : FlagVector (*@$\sigma$@*) => downwardFlagVectorQuot f) (...)

notation "(*@$\llbracket$@*)" f "(*@$\rrbracket_0$@*)" => (downward f)
\end{lstlisting}
The last definition is the quotient-level downward operator, and the final
line of the snippet is the Lean notation declaration saying that
\(\llbracket f\rrbracket_0\) expands to \lean{downward f}.  In mathematical
prose we write this operator as \(\llbracket-\rrbracket_\sigma\), following
Section~\ref{sec:background-downward}.  In Lean code, the subscript \(0\)
emphasizes that the codomain is the
\(\emptyt\)-typed algebra.  The notation is compact, but the implementation
contains the full sequence of quotient lifts: from concrete labeled graphs to
flags, from flags to flag vectors, and finally from flag-vector
representatives to elements of the quotient algebra.

It remains to connect this algebraic map to the semantics of positive
homomorphisms.  The theorem below is the Lean form of the random-extension
identity from Equation~\eqref{eq:ensemble}, so we first spell out the objects
appearing in the statement.  The type \lean{PositiveHomSpace}, with argument
\(\sigma\), is the
measurable space whose points are positive homomorphisms on
\(\mathcal{A}^\sigma\), represented in Lean by their density profiles.  A term
of type \lean{PositiveHomSpace} applied to \(\sigma\) can be converted back to
the corresponding positive homomorphism by
\lean{PositiveHomSpace.toPosHom} applied to that term.
At this point \(\phi\) is just a point of this measurable space, not yet a
random sample.  The Lean expression
\lean{PositiveHomSpace.toPosHom $\phi$ f} is the value of the positive
homomorphism represented by \(\phi\) on \(f\).

The hypothesis named \lean{h$\sigma$} in the theorem is the Lean assumption
that \(\phi_0(\langle\sigma\rangle_0)>0\).  It says that the type \(\sigma\),
viewed as an empty-type flag-algebra element, has positive density under
\(\phi_0\).  Under this hypothesis Lean constructs a probability measure
\(\mathbb{P}\) on \lean{PositiveHomSpace} with argument \(\sigma\).  This
\(\mathbb{P}\) is the formal
random extension \(\Ext_\sigma(\phi_0)\): it is the distribution obtained by
randomly choosing an occurrence of the type \(\sigma\) in the limiting object
represented by \(\phi_0\) and using that occurrence as the labeled type.  Once
\(\mathbb{P}\) is present, a point of this measurable space may be regarded as
a random positive homomorphism drawn from this distribution.  The
displayed theorem states the defining identity of that distribution:
\begin{lstlisting}
theorem (*@\texttt{exists\_probMeasure\_extend\_emptyType\_positiveHom}@*) {n(*@$_0$@*) : (*@$\mathbb{N}$@*)} {(*@$\sigma$@*) : FlagType (Fin n(*@$_0$@*))}
    {(*@$\phi_0$@*) : PositiveHom (*@$\emptyt$@*)} (h(*@$\sigma$@*) : (*@$\phi_0$@*) ((*@$\langle\sigma\rangle_0$@*)) > 0)
    : (*@$\exists$@*) (*@$\mathbb{P}$@*) : ProbabilityMeasure (PositiveHomSpace (*@$\sigma$@*)), (*@$\forall$@*) f : FlagAlgebra (*@$\sigma$@*),
        (*@$\int$@*) ((*@$\phi$@*) : PositiveHomSpace (*@$\sigma$@*)), PositiveHomSpace.toPosHom (*@$\phi$@*) f (*@$\partial$@*)(*@$\mathbb{P}$@*)
        = (*@$\phi_0$@*) (*@$\llbracket$@*)f(*@$\rrbracket_0$@*) / (*@$\phi_0$@*) (*@$\llbracket$@*)(1 : FlagAlgebra (*@$\sigma$@*))(*@$\rrbracket_0$@*)
  := ...
\end{lstlisting}
The integral line uses Mathlib's binder notation for integration with respect
to a specified measure.  In general, Lean writes
\lean{$\int$ x, g x $\partial\mu$} for the integral of the function
\(x\mapsto g(x)\) with respect to the measure \(\mu\).  The comma separates
the bound variable from the integrand, and \(\partial\mu\) supplies the
measure; the symbol \(\partial\) is only notation for ``with respect to'', not
a derivative.  In the theorem above, the bound variable is
\lean{$\phi$ : PositiveHomSpace $\sigma$}, the integrand is
\lean{PositiveHomSpace.toPosHom $\phi$ f}, and the measure is
\(\mathbb{P}\).  Since \(\mathbb{P}\) is a probability measure, this integral
is mathematically an expectation:
\[
  \mathbb{E}_{\phi\sim \mathbb{P}}[\phi(f)].
\]
With this reading, the theorem says exactly that \(\mathbb{P}\) satisfies
Equation~\eqref{eq:ensemble}:
\[
  \mathbb{E}_{\phi\sim \mathbb{P}}[\phi(f)]
  =
  \frac{\phi_0(\llbracket f\rrbracket_0)}
       {\phi_0(\llbracket \mathbf{1}_\sigma\rrbracket_0)} .
\]
The denominator is the normalizing value of the typed unit after applying the
downward operator, and it is positive under the hypothesis
\(\phi_0(\langle\sigma\rangle_0)>0\).  Thus the theorem gives a concrete way
to compute the unlabeled value \(\phi_0(\llbracket f\rrbracket_0)\) from typed
evaluations: randomly label the limiting object represented by \(\phi_0\) with
a copy of \(\sigma\), evaluate \(f\) in the resulting \(\sigma\)-typed positive
homomorphism, average these values, and multiply by the positive normalizing
denominator.

The payoff is the Lean transfer theorem corresponding to
\Cref{thm:downward-nonneg}: if \(f\) is non-negative under every positive
homomorphism in the
\(\sigma\)-typed algebra, then its downward image is non-negative under every
empty-type positive homomorphism.
\begin{lstlisting}
theorem downward_preserve_semanticCone {n(*@$_0$@*) : (*@$\mathbb{N}$@*)} {(*@$\sigma$@*) : FlagType (Fin n(*@$_0$@*))}
    (f : FlagAlgebra (*@$\sigma$@*)) (hf : f (*@$\in$@*) semanticCone (*@$\sigma$@*))
    : (*@$\llbracket$@*)f(*@$\rrbracket_0$@*) (*@$\in$@*) semanticCone (*@$\emptyt$@*) := ...
\end{lstlisting}
The proof is the formal use of Equation~\eqref{eq:ensemble}.  To show that
\(\llbracket f\rrbracket_0\) lies in the empty-type \lean{semanticCone}, Lean
tests it against an arbitrary empty-type positive homomorphism \(\phi_0\).  If
\(\phi_0(\langle\sigma\rangle_0)=0\), a separate zero-density lemma shows that
every downward value evaluates to \(0\) at \(\phi_0\).  If
\(\phi_0(\langle\sigma\rangle_0)>0\), the random-extension theorem provides
the probability measure \(\mathbb{P}=\Ext_\sigma(\phi_0)\), obtained by randomly
labeling an occurrence of \(\sigma\), and rewrites
\(\phi_0(\llbracket f\rrbracket_0)\) as a positive normalizing factor times
the expectation of \(\phi(f)\) over \(\phi\sim \mathbb{P}\).  Since
\(f\in\mathcal{C}^\sigma\), every typed positive homomorphism \(\phi\) has
\(\phi(f)\ge 0\).  The expectation is therefore non-negative, and multiplying
by the positive normalizing factor preserves non-negativity.  This is the
transfer from typed inequalities to the unlabeled order used later in density
bounds.

\subsection{Forbidden Graphs as Assumptions, Not New Theories}
\label{sec:forbidden}

Section~\ref{sec:background} presented flag algebras in the usual constrained
form: a finite forbidden family \(\mathcal{H}\) is fixed first, and the flag
types, flags, algebras, and positive homomorphisms are all built inside the
induced-\(\mathcal{H}\)-free world.  In that setup, an inequality
\(f\leq_\sigma g\) is already an inequality in the constrained
\(\sigma\)-typed algebra: the forbidden configurations have been removed
before the semantic order is tested.

Our Lean formalization deliberately separates these two steps.  The core
development constructs the ambient flag algebra \(\mathcal{A}^\sigma\) of all
finite simple graphs, with no forbidden family as a parameter.  The definitions
of flags, densities, products, the downward operator, and positive
homomorphisms are therefore reusable across extremal problems.  A forbidden
graph is introduced only later, as an assumption on the semantic comparison
between two ambient algebra elements \(f,g\in\mathcal{A}^\sigma\): we test only
base positive homomorphisms \(\phi_0\) in which the forbidden graph has induced
density zero, and then require the desired typed inequality to hold almost surely
after randomly labeling an occurrence of \(\sigma\).  Thus forbiddenness is not a new
flag-algebra universe in the implementation; it is an a posteriori condition on
which positive homomorphisms are allowed to interpret the ambient algebra.

The order used for this a posteriori condition is the
\emph{ensemble semantic order}.  The current API implements the case of one
forbidden graph \(H\); a forbidden family would replace the single zero-density
assumption below by one such assumption for every member of the family.

\begin{definition}[Ensemble semantic order]\label{def:ensemble-semantic-order}
Fix a forbidden graph \(H\), a flag type \(\sigma\), and elements
\(f,g\in\mathcal{A}^\sigma\) in the ambient flag algebra.  We write
\(f\leq^{\mathrm{ens}}_{H,\sigma} g\) if, for every empty-type positive
homomorphism \(\phi_0\) such that, viewing \(H\) and \(\sigma\) as
empty-type flag-algebra elements,
\[
  \phi_0(H)=0
  \qquad\text{and}\qquad
  \phi_0(\sigma)>0,
\]
the random extension \(\phi\sim\Ext_\sigma(\phi_0)\) satisfies
\[
  \mathbb{P}\bigl[\phi(f)\leq\phi(g)\bigr]=1.
\]
\end{definition}

\begin{remark}[Relation to built-in induced-\(H\)-free semantics]
There is another standard way to impose the same forbidden-graph condition:
one builds the induced-\(H\)-free restriction into the flag algebra from the
beginning, so that the algebra itself is already an algebra of
induced-\(H\)-free flags.  Semantic inequalities are then tested inside that
separately built constrained algebra.  Our ensemble order does something
different.  It keeps the ambient algebra of all finite graphs, takes a limit
of large induced-\(H\)-free graphs, and then chooses the labels randomly in
that limit.  Consequently, it tests what is true for almost every labeled
view of a constrained limit.

For unlabelled density bounds, this distinction is harmless: the final
induced Tur\'an-density statements are recovered by the transfer theorem below.
For labelled intermediate inequalities, however, the two viewpoints are not
automatically identical.  The built-in induced-\(H\)-free algebra may allow us to put
the labels at a specially chosen, very rare vertex configuration, whereas the
ensemble order only sees configurations that occur with nonzero probability
when the labels are sampled at random.  Later in the paper we analyze exactly
when this difference disappears, and we also give examples showing that it can
be real.
\end{remark}

The first condition says that \(\phi_0\) assigns induced density zero to the
forbidden graph \(H\); the second says that \(\sigma\) has positive density under
\(\phi_0\), so a random occurrence of \(\sigma\) can be chosen.  The last line
then compares \(f\) and \(g\) not against all positive homomorphisms on a
separately built induced-\(H\)-free algebra, but against almost every \(\sigma\)-typed positive
homomorphism obtained by randomly labeling from such a \(\phi_0\).  The Lean
predicate \lean{forbidLE} is the formalization of
\Cref{def:ensemble-semantic-order}:
\begin{lstlisting}
def forbidLE {n(*@$_0$@*) : (*@$\mathbb{N}$@*)} {(*@$\sigma$@*) : FlagType (Fin n(*@$_0$@*))} (H : FinFlag (*@$\emptyt$@*)) (f g : FlagAlgebra (*@$\sigma$@*)) : Prop :=
  (*@$\forall\;$@*)((*@$\phi_0$@*) : PositiveHom (*@$\emptyt$@*)),
    (*@$\phi_0$@*) (*@$\langle\sigma\rangle_0$@*) > 0
    -> (*@$\,$@*) (*@$\phi_0$@*) (*@$\llbracket$@*)basisVector H(*@$\rrbracket_0$@*) = 0
    -> (*@$\,$@*) (*@$\mathbb{P}$@*)[(*@$\phi_0$@*)] {(*@$\phi$@*) : PositiveHomSpace (*@$\sigma$@*) | (*@$\phi$@*) f (*@$\leq\;$@*)(*@$\phi$@*) g} = 1
\end{lstlisting}
Read the header first.  The implicit arguments
\lean{n$_0$ : $\mathbb{N}$} and
\lean{$\sigma$ : FlagType (Fin n$_0$)} fix the flag type, while the explicit
arguments \lean{H : FinFlag $\emptyt$} and
\lean{f g : FlagAlgebra $\sigma$} name the forbidden graph and the two
ambient \(\sigma\)-typed algebra elements being compared.  In particular,
\lean{f} and \lean{g} do not live in a separately built induced-\(H\)-free algebra.

The body then ranges over empty-type positive homomorphisms
\lean{$\phi_0$ : PositiveHom $\emptyt$}.  The hypothesis
\lean{$\phi_0$ $\langle\sigma\rangle_0$ > 0} says that the type \(\sigma\)
has positive density after its labels are forgotten, which is the condition
needed to form the random-extension measure \(\mathbb{P}[\phi_0]\).  The
hypothesis \lean{$\phi_0$ $\llbracket$basisVector H$\rrbracket_0$ = 0}
imposes the induced forbidden-graph assumption: \lean{basisVector H} is the basis
element for \(H\), viewed in the empty-type flag algebra.

The conclusion says that the event
\lean{\{$\phi$ : PositiveHomSpace $\sigma$ | $\phi$ f $\leq$ $\phi$ g\}} has
probability one under this random extension.  Here a sampled
\lean{$\phi$} is a \(\sigma\)-typed positive
homomorphism, so the event is exactly that the sampled homomorphism evaluates
\lean{f} at most \lean{g}.

What remains is to connect
the semantic predicate above to the corresponding asymptotic bound on induced densities.
In the theorem below, the variables have types
\lean{H : SimpleGraph (Fin n)} and \lean{F : SimpleGraph (Fin m)}. Here \(H\) is
the forbidden graph and \(F\) is the target graph.  The expression \lean{F.toFlagAlgebra} converts
the simple graph \(F\) into the corresponding element of the empty-type flag
algebra \lean{FlagAlgebra $\emptyt$}, whose evaluation is the induced
\(F\)-density.  Likewise, \lean{H.toFinFlag} converts \(H\) into a finite
empty-type flag \lean{FinFlag $\emptyt$}, so the zero-density hypothesis in
\lean{forbidLE} with \lean{H.toFinFlag} is the semantic form of being induced-\(H\)-free.

The formal bridge packages this passage from semantic order to the finite
asymptotic statement:
\begin{lstlisting}
theorem generalizedTuranDensity_le_of_forbidLE
    {n m : (*@$\mathbb{N}$@*)} {H : SimpleGraph (Fin n)} {F : SimpleGraph (Fin m)} {c : (*@$\mathbb{R}$@*)} (hc : 0 <= c)
    (h : F.toFlagAlgebra <=[H.toFinFlag] c (*@$\leansmul$@*) (1 : FlagAlgebra (*@$\emptyt$@*)))
    : generalizedTuranDensity H F <= c := ...
\end{lstlisting}
The hypothesis is an \(\emptyt\)-typed \lean{forbidLE} inequality: the target
graph \(F\), viewed as an element of the empty-type flag algebra, is at most
\(c\cdot \mathbf{1}\) under the ensemble semantic order for the forbidden
induced subgraph \(H\).  The conclusion is the corresponding induced
Tur\'an-density bound, written in Lean as
\lean{generalizedTuranDensity H F <= c}; mathematically, it is the assertion
\[
  \tdensity{F}{H}
  =
  \lim_{N\to\infty}
  \max_{\substack{|V(G)|=N\\ G\text{ induced-}H\text{-free}}}
  \den{F}{G}
  \leq c .
\]
Thus a certificate proof only has to
establish one semantic inequality in the ambient algebra; the passage from
that inequality to the induced extremal-density statement is handled once and
reused.

\section{The Reflection Layer}
\label{sec:reflection}

The definitions in Section~\ref{sec:abstract} are intentionally close to the
mathematics.  Many of them are marked \lean{noncomputable}: this does not mean
that they are unusable as specifications, but that Lean should not be expected
to reduce them as executable programs.  This is a normal situation when
formalizing mathematics.  A definition is often written to express the
mathematical object directly, not to specify the data structures and loops by
which one should compute it.  For instance, we may identify flags up to
isomorphism, define a finite collection by a predicate, or use an existence
proof to choose a representative.  These choices are harmless, and often
preferable, at the level of theorems: they keep statements invariant under
isomorphism and avoid committing the theory to an arbitrary encoding.

The cost is that the computational content is only implicit.  To evaluate such
a definition, Lean would have to turn propositions such as adjacency or
isomorphism into data, enumerate finite sets described abstractly, and choose
representatives of quotient classes.  In proofs, this information can be
supplied by lemmas or by classical reasoning.  As an executable program,
however, there is no direct instruction saying which edge list to inspect,
which vertex permutations to try, or which representative to return.  Thus
\lean{noncomputable} marks a boundary: the definition is a faithful
mathematical specification, but not the program we want Lean to run repeatedly.

To apply the framework to concrete graphs, however, we need exactly the finite
computations that the specification layer leaves implicit.  Given the named finite
flags appearing in a certificate, Lean must decide whether two concrete flags
represent the same isomorphism class, compute subflag densities, and reduce
products among named flags.  These tasks are finite, but the abstract
definitions do not provide executable instructions for carrying them out.  They
describe equality through quotients and the existence of isomorphisms, and they
describe counts using abstract finite sets.  For actual evaluation, we therefore
provide separate Boolean algorithms and prove that their results agree with the
abstract definitions.  We call this pattern \emph{reflection}: Lean runs a
Boolean program on concrete data, then uses a correctness theorem saying that
the result has the intended meaning in the specification layer.

This section develops the reflection layer as an executable mirror of the
specification layer.  It introduces concrete graph and flag representations,
optimized Boolean procedures for isomorphism and density computations, and
adequacy theorems proving that these procedures agree with the quotient-based
definitions.  This is the reflection pattern used throughout the rest of the
development: a concrete representation, an evaluator, and a theorem connecting
the evaluator to the abstract specification. \gw{I'm here.}

\subsection{Concrete Graph Representations}

Mathlib's \lean{SimpleGraph} represents a graph as a \lean{Prop}-valued
adjacency relation:
\begin{lstlisting}
structure SimpleGraph (V : Type*) where
  Adj      : V -> V -> Prop
  symm     : Symmetric Adj
  loopless : Irreflexive Adj
\end{lstlisting}
This design is mathematically natural but not directly computable: checking
adjacency requires evaluating a \lean{Prop}, which \lean{native\_decide}
cannot reduce.

We introduce \lean{Sym2Graph n} in
\lean{LeanFlagAlgebras.FlagAlgebra.Compute.Basic}:
\begin{lstlisting}
structure Sym2Graph (n : N) where
  edges       : Finset (Sym2 (Fin n))
  edges_valid : forall e in edges, not e.IsDiag
\end{lstlisting}
This represents a graph on $\{0,\ldots,n-1\}$ as a finite set of unordered
non-diagonal pairs.  Every operation on \lean{Sym2Graph n}
is computable: membership in \lean{edges} is decidable by \lean{Finset}
membership.  The same file derives the \lean{Fintype} instances needed to
make graph isomorphism decidable:
\begin{lstlisting}
instance : Fintype (G1 ->g G2) := ...   -- graph embeddings
instance : Fintype (G1 =~g G2) := ...   -- graph isomorphisms
instance : Decidable (Nonempty (G =~f G')) := ...  -- flag isomorphism
\end{lstlisting}
These instances enumerate candidates (all injections from $V(G_1)$ to
$V(G_2)$, or all bijections) and filter those that preserve adjacency and
type embeddings.  This reduces ``are these two labeled graphs isomorphic?'' to
a finite, decidable search.

\lean{Sym2Graph n} connects to the abstract world via a coercion to
\lean{SimpleGraph (Fin n)}:
\begin{lstlisting}
def Sym2Graph.toSimpleGraph (G : Sym2Graph n) : SimpleGraph (Fin n) where
  Adj i j   := s(i, j) in G.edges
  symm      := ...
  loopless  := ...
\end{lstlisting}
The key line is \lean{Adj i j := s(i, j) \(\in\) G.edges}: \lean{Finset}
membership replaces the \lean{Prop}-valued adjacency relation.  Symmetry
follows from the symmetry of \lean{Sym2}, and \lean{loopless} follows from
the \lean{edges\_valid} field.  Registering this as a \lean{Coe} instance
means that wherever a \lean{SimpleGraph (Fin n)} is expected, a
\lean{Sym2Graph n} can be used directly.

The same concrete representation also supports a more usable isomorphism
test.  The generic \lean{Fintype}-based instances above are enough for
decidability, but repeated density checks need a faster decision procedure.
The reflected checker first rejects impossible cases by simple isomorphism
invariants such as edge counts; for typed flags, it also fixes the labeled
vertices and only permutes the remaining unlabeled vertices.  Correctness
theorems for the Boolean outcomes justify using this checker as the
\lean{DecidableEq} instance for reflected flags.  Equality still means equality
in the quotient setoid; only the decision procedure has changed.

\subsection{Correctness as Adequacy}

For density computations, we define a computable density function
\lean{sym2EmptyTypeFlagDensity} entirely over \lean{Finset} operations.  The
coercion above makes abstract definitions \emph{applicable} to
\lean{Sym2Graph} inputs, but applicability is not computability:
\lean{flagDensity} composed with the coercion still unfolds to non-computable
quotient machinery.  The concrete density function sidesteps this by computing
directly over the \lean{Finset} of edges, producing a value that
\lean{native\_decide} can evaluate.

The reflected density computation also replaces abstract subgraph enumeration
with data-oriented enumeration.  Rather than enumerate abstract
\lean{LabeledSubgraph}s directly, it enumerates \lean{Sym2InducedSubgraph}s,
represented only by finite vertex sets, and computes their edge sets by
filtering the host graph's \lean{Finset} of edges.  For typed flags,
\lean{Sym2InducedLabeledSubgraph} additionally stores the proof that the type
vertices are included.  A count-level adequacy theorem proves that this
enumeration has the same cardinality as the abstract \lean{LabeledSubgraph}
enumeration.  Thus the cost of quotient and subgraph reasoning is paid once in
the proof of adequacy, not in every generated density lemma.

We use \emph{adequacy theorem} for this form of correctness theorem: a theorem
stating that the reflected computation agrees with the abstract definition it
implements.  For the density evaluator, the central adequacy theorems are:
\begin{lstlisting}
theorem flagDensity1_eq_sym2FlagDensity1
    (F : Sym2EmptyTypedFlag m) (G : Sym2EmptyTypedFlag n) :
    flagDensity1 F.toFlag G.toFlag =
    sym2EmptyTypeFlagDensity1 F G := ...

theorem flagDensity2_eq_sym2FlagDensity2
    (F0 : Sym2EmptyTypedFlag m0) (F1 : Sym2EmptyTypedFlag m1)
    (G : Sym2EmptyTypedFlag n) :
    flagDensity2 F0.toFlag F1.toFlag G.toFlag =
    sym2EmptyTypeFlagDensity2 F0 F1 G := ...
\end{lstlisting}
The right-hand sides are defined entirely over \lean{Sym2Graph} and are fully
computable.  After rewriting a density goal with these theorems, the goal
becomes a decidable proposition about a finite computation over
\lean{Finset}, which \lean{native\_decide} can evaluate directly.

For a PL audience, the important point is that the adequacy theorem is a
refinement theorem: the executable program computes exactly the denotation of
the abstract specification.  The proof factors through count-level adequacy
(\lean{labeledSubgraphListCount\_eq\_sym2InducedSubgraphListCount}) and
density-level adequacy
(\lean{labeledSubgraphListDensity\_eq\_sym2InducedSubgraphListDensity}).
Consequently, the generated theorem
\[
  \lean{flagDensity2 F0 F1 G = q}
\]
is not a certificate that the external JSON table is trusted; it is a kernel
proof that the reflected evaluator, connected to the abstract definition by
adequacy, computes the rational number $q$.

\subsection{Elaboration-Time Theorem Generation}

Having established the adequacy theorems, we need to generate and prove the
density and multiplication lemmas required for each specific theorem.  Rather
than writing these by hand, we use Lean~4's elaboration monad to generate
them automatically from precomputed data.

The elaboration macro \lean{load\_flag\_pair\_density\_theorems} reads a JSON
file containing density values.  It is implemented in
\lean{Flags/Densities/DensityLoader.lean}; for each entry
\lean{(F1, F2, G, p/q)} it generates and immediately proves the theorem:
\begin{lstlisting}
@[simp]
theorem flagDensity2_Flag_5_1_0_3_Flag_5_1_0_7_Flag_5_0_0_2
    : flagDensity2 Flag_5_1_0_3 Flag_5_1_0_7 Flag_5_0_0_2 = 3/10
  := by
  delta Flag_5_1_0_3 Flag_5_1_0_7 Flag_5_0_0_2
  rw [flagDensity2_eq_sym2FlagDensity2]
  native_decide
\end{lstlisting}
The three proof steps are: unfold the concrete flag definitions (making their
\lean{Sym2Graph} structure explicit), rewrite with the adequacy theorem, then
call \lean{native\_decide} to evaluate the resulting finite computation.  The
same elaboration pattern is used by \lean{MulLoader.lean} for flag
multiplication tables, and by family-specific loaders such as
\lean{load\_triangle\_density\_theorems} and
\lean{load\_k4\_density\_theorems} that combine the generic adequacy lemma
with a forbidden-graph filter.

The macro \lean{generate\_basisVector\_lemmas} in \lean{API/Basic.lean} plays
a similar role for unit-vector identities.  Given \lean{n} and a count, it
emits simp lemmas of the form
$\lean{[[basisVector <n, Flag\_n\_0\_0\_i>]] = FlagAlgebra\_n\_0\_0\_i}$ for
each index in range, so that subsequent tactics can rewrite quotient classes
into named constants without invoking \lean{Quotient.out\_inj.mp rfl}
explicitly.

This design also changes the proof-engineering scale.  The JSON files are not
used as proof certificates in the logical sense; they are build inputs that
determine which theorems Lean should attempt to prove and what rational value
they should state.  If a table entry is wrong, the generated theorem fails at
\lean{native\_decide}.  If the table entry is right, the theorem is stored as
a named lemma and later proofs use it by rewriting rather than by recomputing
the density from scratch.  In effect, the elaborator performs a verified
specialization pass: it turns a general reflective evaluator into thousands
of small, reusable theorem constants specialized to the flags that appear in
the certificate.

\section{From a Recurring Proof Pattern to a Certificate-to-Lean Compiler}
\label{sec:flagmatic}
\label{sec:bridge}% legacy alias: this section absorbed the former ``Bridge'' section

Software such as Flagmatic~\cite{vaughan2013flagmatic} searches for
flag-algebra certificates automatically, but a certificate is numerical
data---lists of flags, density and multiplication tables, and rational
positive semidefinite matrices---not a machine-checked proof.  Turning one
into a verified theorem means transcribing every number into the proof
assistant, reproducing the bookkeeping that links subflags to host flags, and
checking positive-semidefiniteness, densities, and product expansions inside
the kernel.

The specification of Section~\ref{sec:abstract} and the reflection layer of
Section~\ref{sec:reflection} make it possible to carry this out by hand.  We
did so for our first case studies---computing the expansions, products, and
normalizations step by step---and found the proofs nearly identical in
structure.  Mantel's theorem, the $C_4$-density bound, and the Erd\H{o}s
pentagon theorem differ in their flags, their coefficients, and their positive
semidefinite matrices, but each follows the same sequence of moves, in the same
order, with only the sizes and the data changing.  That regularity is what
convinced us the manual work could be mechanized, and it showed which recurring
step each tactic would later have to automate.

The section is organized around that shared shape.
Section~\ref{sec:flagmatic-shape} isolates it as a four-part skeleton;
Sections~\ref{sec:flagmatic-gen}--\ref{sec:flagmatic-close} show how each part
is discharged by a reusable, kernel-checked component; and
Section~\ref{sec:flagmatic-script}
describes the translator that assembles the whole proof from a Flagmatic
certificate alone.  The translator realizes the second contribution of
Section~\ref{sec:intro}: a certificate-to-proof pipeline whose output the Lean
kernel checks and whose own computation never enters the trusted base
(Section~\ref{sec:flagmatic-trust}), validated on six certificates
(Section~\ref{sec:flagmatic-scenarios}).

\subsection{The recurring four-part shape}
\label{sec:flagmatic-shape}

Each of these proofs has the same four-part skeleton.  We list the parts in the
order a finished proof presents them:

\begin{enumerate}
  \item \textbf{Flags and density data.}  The flags of the sizes the proof needs
    (only those that are induced-free for the chosen forbidden graph), together with the pair-density
    and typed-multiplication identities among them that the body rewrites with.
  \item \textbf{Certificate matrices and flag vectors.}  For each subflag type
    $\sigma_t$, a rational matrix $M_t$ paired with a vector $v_t$ of
    $\sigma_t$-flags; together they form the quadratic form
    $\langle v_t, M_t\, v_t\rangle$.
  \item \textbf{Lifting the objective (when needed).}  If the objective is
    smaller than the host size, one auxiliary lemma first lifts it to host size,
    dropping the contributions that contain an induced copy of the forbidden graph.
  \item \textbf{The main theorem.}  The bound $P \leq_{\mathcal{H}}
    c\cdot\mathbf{1}$, discharged by adding each block's quadratic form,
    expanding the unit at host size, reducing every downward-averaged product,
    normalizing, and closing by non-negativity.
\end{enumerate}

Each part of the shape is discharged by a verified, reusable component, taken up
in turn by the next four subsections.


\subsection{Generating the induced forbid-free flags and their data}
\label{sec:flagmatic-gen}

% TODO/REVISIT: How much detail to give here (pruned generation, completeness
% via native_decide, adequacy-based density/product proofs) depends on how much
% the reflection layer (Section~\ref{sec:reflection}) ends up covering about the
% same machinery.  That  issection not finalized yet, so this subsection is left
% as-is for now -- revisit once Section~\ref{sec:reflection}'s level of detail is
% settled, to avoid over/under-explaining or duplicating it.

The forbidden graph is an explicit \lean{Sym2Graph} term, and the flag,
density, and product data are \emph{generated inside Lean}---by a handful of
elaboration commands, each parametric in the forbid term \lean{K}---rather than
read from precomputed tables.  The commands emit exactly the constants the proof
needs:
\begin{lstlisting}
def K3 : Sym2Graph 3 := completeSym2Graph 3
generate_pruned_forbid_free_empty_typed_flags 3 K3
generate_pruned_forbid_free_flags 3 1 0 K3
generate_pruned_flag_pair_density_theorems 2 3 1 0 K3
generate_pruned_forbid_free_mul_theorems 2 3 1 0 K3
\end{lstlisting}
These commands enumerate only the induced-\(\mathcal{H}\)-free flags: they filter the
small underlying-graph enumeration by an isomorphism-invariant induced-freeness
predicate \emph{before} building the typed flags, so the
flags containing an induced forbidden graph are pruned before they are ever
materialized---this is what keeps the typed enumeration tractable at host sizes
where the full typed \lean{native\_decide} would otherwise wall.  Completeness is
a lemma of the form \lean{flagSetHfree\_N\_k\_m\_<K>}~$=$~\lean{univ.filter p},
proved by a single \lean{native\_decide} that references only the surviving
constants; the pair densities and product expansions are each proved by the
adequacy rewrite of Section~\ref{sec:reflection} followed by an evaluator call.
All numerical content is therefore recomputed and re-checked in Lean every
build; the certificate never asserts a density value the kernel does not
reproduce.  Crucially, the generators take the forbidden graph as a \emph{term
argument} (an \lean{ident} naming a \lean{Sym2Graph}) and build the induced
freeness predicate from it, so a new forbidden graph is a new
\lean{Sym2Graph} definition rather than a new per-forbid branch in the
meta-commands.

\subsection{The certificate matrices and their flag vectors}
\label{sec:flagmatic-psd}

Each matrix is the part of the certificate that carries the bound.  The
flag-algebra method does not prove $P \leq_{\mathcal{H}}
c\cdot\mathbf{1}$ by manipulating $P$ directly; it adds to $P$ a quantity that is
provably non-negative and chosen so that the sum is bounded by $c$ term by term.
That quantity is a quadratic form $\langle v_t, M_t\, v_t\rangle$ in the
$\sigma_t$-flags, and a quadratic form with $M_t$ positive semidefinite is the
flag-algebra analogue of a sum of squares: it is non-negative for every
assignment, so adding it can only raise $P$, never lower it.  The SDP solver
(Flagmatic) is exactly the search for such matrices---it returns the $M_t$ for
which the slack $c\cdot\mathbf{1} - P - \sum_t \langle v_t, M_t\, v_t\rangle$ has
only non-negative coefficients.  We take the $M_t$ as given; the one property Lean
must verify is that each is positive semidefinite.

We establish that property through an $LDL^\top$ factorization: a symmetric $M$
that factors as $M = L D L^\top$ with $D$ diagonal and its diagonal entries
$d \geq 0$ is positive semidefinite, because
\[
  x^\top M x \;=\; x^\top L D L^\top x \;=\; \sum_i d_i\,(L^\top x)_i^2 \;\geq\; 0 .
\]
The right-hand side is a sum of squares---the concrete form of the analogy from
the intro.  We use this factorization rather than a spectral test because
checking it needs only rational arithmetic: with exact rational entries
$M = L D L^\top$ is an equality verified over $\mathbb{Q}$---no eigenvalues, no
floating point---and the result lifts to $\mathbb{R}$, where the flags live.  Each
block is a few rational definitions and a one-line proof, paired with the vector
of $\sigma$-flags it is contracted against:
\begin{lstlisting}
-- the rational matrix M, its factorization data (dM, LM), and the real lift
def M  : Matrix (Fin 2) (Fin 2) Q := !![(1/2 : Q), (-1/2 : Q); (-1/2 : Q), (1/2 : Q)]
def dM : Fin 2 -> Q := ![(1/2 : Q), 0]
def LM : Matrix (Fin 2) (Fin 2) Q := !![(1 : Q), 0; (-1 : Q), (1 : Q)]
noncomputable def M_real : Matrix (Fin 2) (Fin 2) R := ratMatrixToReal M

-- positive semidefiniteness, proved from the factorization
theorem M_real_posSemidef : M_real.PosSemidef := by
  psd_real_ldlt M LM dM

-- the vector of sigma-flags M is contracted against
def sigma : FlagType (Fin 1) := FlagType_1_0
noncomputable def v : FlagAlgebraVec sigma 2 := ![FlagAlgebra_2_1_0_0, FlagAlgebra_2_1_0_1]
\end{lstlisting}
The positive-semidefiniteness proof is a single call to
\lean{psd\_real\_ldlt M LM dM} (in \lean{API/Matrix/PosSemiDef.lean}): given the
precomputed factorization, the tactic proves \lean{M\_real.PosSemidef}
automatically.  It refines the goal
through the generic lemma \lean{posSemidef\_real\_of\_LDLt} and discharges its two
side conditions: the diagonal non-negativity $d \geq 0$ by
\lean{fin\_cases}/\lean{norm\_num}, and the factorization equality
$M = L D L^\top$ by \lean{decide +kernel}---a decidable
rational matrix equality the kernel evaluates entry by entry, using no compiled
native code.  The emitter therefore supplies only the rational triple
$(M, L, d)$.  Because the
equality is kernel-checked, a wrong $L$ or $d$ surfaces as a compile error
on this line, never as a silently accepted non-PSD matrix---the one place we
deliberately pay for in-kernel evaluation (Section~\ref{sec:flagmatic-trust}).
The vector $v$ alongside it merely lists the $\sigma$-flags by name;
together $(M, v)$ form the quadratic form
$\llbracket\langle v, M\,v\rangle\rrbracket$---one per $\sigma$-type---that the
main theorem adds to the target.

\subsection{Lifting the objective to host size}
\label{sec:flagmatic-expand}

The host size $N$ is the world the densities live in, and the objective is not
always a flag of that size.  When it is smaller we simply lift it to host
size before the rest of the proof runs; this is the proof's one branch point.  If the objective is already host-size
(\emph{Branch~A}: $n_{\mathit{obj}}=N$; the $C_4$-density bound, Erd\H{o}s
pentagon) the
step is empty.  If it is smaller (\emph{Branch~B}: $n_{\mathit{obj}}<N$; Mantel,
$K_4$- and $K_5$-Tur\'an) a single call to \lean{flag\_expand\_hfree N K} does
it, rewriting the objective (via \lean{basisVector\_quot\_forbidEq\_sum}) over the
$N$-vertex induced-$\mathcal{H}$-free flags and dropping the host terms that
contain an induced copy of $\mathcal{H}$.  Its closing arithmetic needs the
relevant \lean{flagDensity\textsubscript{1}} values, which the translator emits as
\lean{@[simp]} lemmas.

\subsection{The main theorem}
\label{sec:flagmatic-close}
\sh{I'm here.}

\paragraph{Adding quadratic forms and closing.}
% Phrased in terms of the established forbid order $\leq_{\mathcal{H}}$ to avoid
% introducing the undefined term "forbid-cone" (used nowhere else).
The standard upper-bound move---add one or more PSD quadratic forms to the
target, then check that the augmented target is
$\leq_{\mathcal{H}} c\cdot\mathbf{1}$---rests on
\lean{forbidLE\_add\_QuadraticForm}, which does the adding (from
$f \leq_{\mathcal{H}} g$ it gives
$f \leq_{\mathcal{H}} g + \llbracket\langle v, M\,v\rangle\rrbracket$):
\begin{lstlisting}
theorem forbidLE_add_QuadraticForm
    (M : Matrix (Fin n) (Fin n) R) (hM : M.PosSemidef) (v : FlagAlgebraVec sigma n)
    : (f <=[F_forbid] g) -> f <=[F_forbid] g + (*@$\llbracket$@*)flagQuadraticForm M v(*@$\rrbracket_0$@*) := ...
\end{lstlisting}
iterated once per type to build a transitive chain of additions.  The rest of the
body performs the check: the unit is
expanded to host size by \lean{expand\_one\_hfree\_at N K} (rewriting
$\mathbf{1}$ as the sum of $N$-vertex induced-\(\mathcal{H}\)-free flags); the
downward-averaged products are discharged by \lean{reduce\_downward\_flagmul},
which first right-associates the sum with
\lean{simp only [downward\_add, add\_assoc]} and then rewrites each head
summand $\llbracket c\,(A\cdot B)\rrbracket$ by its \lean{flagMul\_<A>\_<B>}
identity, moving each rewritten term across the inequality; the resulting
equality between two reordered combinations of flag terms is normalized by
\lean{flagsum\_ac\_sort\_rhs\_pipeline}, which sorts the summands by the index in
each flag's name and proves the reordering up to add-AC (an index-guided
substitute for \lean{ring}/\lean{ac\_rfl}, which grow too slow past ${\sim}20$
terms); and the final non-negativity goal is closed by \lean{flag\_nonneg}, which
distributes an arbitrary positive homomorphism over the sum, splits with
\lean{add\_nonneg}, and discharges each leaf by
\lean{positiveHom\_basisVector\_ge\_zero}.

\paragraph{Metaprogramming robustness.}
These tactics must tolerate the syntactic variation Lean's elaborator
introduces: after unification and implicit-argument resolution, the same flag
term can appear in different syntactic forms.  We normalize the expression tree
with a \lean{collectPrefixConstants} traversal before inspecting names, and
compare flag terms with \lean{Expr.eqv} rather than structural equality.  This
required substantial experimentation with the Lean~4 metaprogramming API and is
among the least portable parts of the development.

\subsection{The translation script}
\label{sec:flagmatic-script}

The pieces of the previous subsections---generated flags, certificate matrices,
the optional objective lift, and the main theorem---are assembled from a
Flagmatic certificate by one script, \lean{flagmatic\_to\_lean.py}.  A single
invocation
\begin{lstlisting}[language={}]
python flagmatic_to_lean.py gen-skeleton <cert>.json <target>.lean
\end{lstlisting}
emits a self-contained Lean file that compiles under \lean{lake build} as a
complete proof, with no \lean{sorry} gaps.  Two diagnostic
subcommands help before any compilation:
\lean{inspect} resolves every certificate string to its Lean identifier, and
\lean{check-deps} validates the certificate and previews the
\lean{generate\_pruned\_*} commands it will emit.

Every later step refers to flags by their canonical index, so the script's
first job is to translate.  A Flagmatic certificate names graphs by compact
strings (\lean{"3:1213"} is the $3$-vertex path, \lean{"3:12(2)"} a $3$-vertex
flag on a size-$2$ type), whereas Lean names them \lean{Flag\_n\_k\_m\_i}.  The
script parses each string into its vertex count, edge set, and type labels, then
isomorphism-matches it against a canonical enumeration that reproduces the order
of Lean's flag generators to recover the index $i$.  Because the script
reproduces that order exactly, the $i$ it computes is already the $i$ in
\lean{Flag\_n\_k\_m\_i}---so no separate flag-to-index table has to be kept in
sync with the Lean side.

With the indices in hand, the script fills the four parts of
Section~\ref{sec:flagmatic-shape}.  Three of them are emitted mechanically, as
fixed templates keyed to the sizes: the \lean{generate\_pruned\_*} commands for
the flags and density data (Part~1, parametric in the forbid term \lean{K}); the
\lean{flag\_expand\_hfree N K} lemma and its \lean{native\_decide} coefficient
table when the objective must be lifted (Part~3, Branch~B only); and the proof
body of Section~\ref{sec:flagmatic-close} (Part~4).  The one part that takes real
computation is the certificate matrices (Part~2): the certificate stores each
block as a pair $(Q_t', R_t)$, and the script reconstructs
$M_t = R_t\,Q_t'\,R_t^\top$ in exact rationals, factors it as $L_t D_t L_t^\top$
(an elimination that tolerates the usual rank deficiency), and emits $M_t$,
$L_t$, $d_t$, the flag vector $v_t$, and the one-line \lean{psd\_real\_ldlt} proof
of Section~\ref{sec:flagmatic-psd}.

For the Mantel certificate ($K_3$-free, $N=3$, $T=1$, Branch~B) the emitted main
theorem reads:
\begin{lstlisting}
theorem mantel_flagAlgebra
    : FlagAlgebra_2_0_0_1
      (*@$\leq$@*)[((*@$\langle$@*)_, Sym2EmptyTypedFlag.toFlag (*@$\llbracket$@*)K3(*@$\rrbracket$@*)(*@$\rangle$@*) : FinFlag (*@$\emptyt$@*))] (1 / 2 : (*@$\mathbb{R}$@*)) (*@$\bullet$@*) (1 : FlagAlgebra (*@$\emptyt$@*)) := by
  have quadraticForm_trans : ... := by
    apply forbidLE_add_QuadraticForm M_real M_real_posSemidef v
    exact forbidLE_refl _ FlagAlgebra_2_0_0_1
  apply forbidLE_trans quadraticForm_trans
  apply forbidLE_trans_forbidEq_right ?_
        (forbidEq_smul (forbidEq_symm (one_forbidEq_forbidExpand_one _ 3)))
  rw [forbidLE_rw_left_add_right mantel_flagAlgebra_expand_under_forbid]
  simp [flagQuadraticForm, v, M_real, ratMatrixToReal, M, Fin.sum_univ_two, add_assoc]
  reduce_downward_flagmul
  expand_one_hfree_at 3 K3
  simp [smul_smul, downward_add, downward_smul]
  flagsum_ac_sort_rhs_pipeline
  apply forbidLE_of_le
  flag_nonneg
\end{lstlisting}
Here the \lean{rw} line is the objective lift of Part~3 (Mantel is Branch~B); a
Branch~A certificate omits it, and a certificate with several SDP blocks simply
repeats the \lean{forbidLE\_add\_QuadraticForm} step once per block.
Choosing the branch, the per-block parameters, and the order of the additions is
the entire ``intelligence'' of the generator; none of its computation is trusted,
and everything it emits is re-checked by Lean
(Section~\ref{sec:flagmatic-trust}).

\subsection{Trust model}
\label{sec:flagmatic-trust}

The script computes a great deal---it isomorphism-matches certificate strings
to indices, assembles $R_t\,Q_t'\,R_t^\top$, factors it as
$L D L^\top$, enumerates induced subgraphs to fill the
Branch~B coefficient table, and decides which generators and imports to
emit---and \emph{none} of this enters the trusted base.  Every emitted
obligation is closed by one of: \lean{decide +kernel} (the PSD matrix
equality, kept in the kernel because it is load-bearing and small enough to
finish there), \lean{native\_decide} (the density evaluations, large but
individually simple rational equalities, checked by compiled code outside the
kernel), or a verified tactic of Sections~\ref{sec:flagmatic-gen}--\ref{sec:flagmatic-close}.  A
miscomputation in the script---a wrong factorization, a wrong density, a
mis-resolved identifier---surfaces as a compile error, never as a silently
accepted false theorem.  The certificate and the solver behind it are
candidate generators; the kernel is the authority.

We report the resulting two-level profile (compiled \lean{native\_decide} for
the bulk density bookkeeping, in-kernel \lean{decide +kernel} for the PSD
step) as a fact about the development rather than a designed trust hierarchy.
Section~\ref{sec:lessons} returns to how a reader who depends on the artifact
should read it.

\subsection{Verified scenarios}
\label{sec:flagmatic-scenarios}

Six certificates have been carried through \lean{gen-skeleton} and produce
files the kernel accepts with no \lean{sorry}
(Table~\ref{tab:flagmatic-scenarios}).  Between them they exercise both
branches, host sizes $N\in\{3,4,5\}$, block counts $1$ through $4$, and the
forbids $K_3$, $K_4$, $K_5$.

\begin{table}[t]
  \centering
  \caption{Certificates carried end-to-end through the translator.  ``Branch''
  is A when $n_{\mathit{obj}}=N$ and B when $n_{\mathit{obj}}<N$; $T$ is the
  number of SDP blocks.}
  \label{tab:flagmatic-scenarios}
  \begin{tabular}{lcccccc}
    \toprule
    Certificate & Forbid & $N$ & $n_{\mathit{obj}}$ & Branch & $T$ & Bound \\
    \midrule
    Mantel         & $K_3$ & 3 & 2 & B & 1 & $1/2$    \\
    K3forbidP3     & $K_3$ & 3 & 3 & A & 1 & $3/4$    \\
    K3forbidC4     & $K_3$ & 4 & 4 & A & 2 & $3/8$    \\
    K4turan        & $K_4$ & 4 & 2 & B & 2 & $2/3$    \\
    ErdosPentagon  & $K_3$ & 5 & 5 & A & 3 & $24/625$ \\
    K5turan        & $K_5$ & 5 & 2 & B & 4 & $3/4$    \\
    \bottomrule
  \end{tabular}
\end{table}

The two largest cases show that a single tactic call scales over a whole
block structure: ErdosPentagon's three quadratic forms of dimensions
$8,6,5$ leave one \lean{reduce\_downward\_flagmul} to dispatch
$8^2+6^2+5^2 = 125$ product terms, and K5turan's four dimension-$8$ forms
leave it $4\times 8^2 = 256$.  None of the six required hand transcription of
the certificate into Lean.

\subsection{What is not automated}
\label{sec:flagmatic-open}

Two caveats bound the claim.  First, the Branch~B host-size expansion
(\lean{flag\_expand\_hfree}) ends in a \lean{ring\_nf} normalization that we
have not formally characterized; it closes every case we have met, but we have
not proved that no certificate produces a form it fails to normalize.  Second, the identifier and forbid
machinery is written to be generic in the forbidden graph---nothing in the
proof body or the generators assumes a complete graph---but only $K_3$,
$K_4$, $K_5$ have been exercised end to end, so a non-$K_n$ forbid such as
$C_4$ or $P_4$ would be the first empirical test of the generality we claim.
We list these alongside the open questions of Section~\ref{sec:unsure}
rather than presenting the pipeline as complete.


\section{Proof-Style Observations}
\label{sec:proof-style}

The case studies produced one recurring proof-style observation worth
reporting on its own, separate from the engineering choices and the
certificate-to-Lean construction.  We do not claim it generalizes beyond extremal combinatorics
formalization; we report it because we noticed it repeatedly.

\paragraph{Density equalities want bijective justification.}
The most common informal move in flag-algebra calculations is to say that
two densities are equal because the two sampling procedures count the same
objects after a reindexing.  In Lean we make this extensional argument
explicit.  The proof first names the two sample spaces, fixes the
\lean{Fintype} instances that will be used to count them, constructs an
equivalence between the sample spaces, and only then simplifies the resulting
rational equality.  A typical proof has the following shape:
\begin{lstlisting}
let S0 := setOfLabeledSubgraphListIsoHl G Hl0
let S1 := setOfLabeledSubgraphListIsoHl G Hl1
let hS0 : Fintype S0 := Fintype.ofFinite S0
let hS1 : Fintype S1 := Fintype.ofFinite S1
let h_iso : Equiv S0 S1 := ...
have card_eq : Fintype.card S0 = Fintype.card S1 :=
  Fintype.card_congr h_iso
simp_all only [Set.toFinset_card]
\end{lstlisting}

This is the pattern used in \lean{flagDensity\_permute} to show that
permuting the order of sampled flags does not change density, and in
\lean{flagDensity\_insert\_empty} to show that inserting an empty flag
leaves the density unchanged.  It also appears in the adequacy proof that
the concrete \lean{Sym2InducedSubgraph} enumeration has the same count as
the abstract \lean{LabeledSubgraph} enumeration.  The invariant the pattern
enforces is that a numerical equality is justified by an explicit
isomorphism between the finite sets being counted; arithmetic simplification
is the last step, not the proof idea.

We initially expected to use Mathlib's existing counting tactics for these
identities and reach for explicit equivalences only as a fallback.  The
opposite happened.  Once the \lean{Fintype}-instance ambiguity of
Section~\ref{sec:fintype} is in play, generic counting tactics often fail to
elaborate because they pick up the ``wrong'' enumeration term for one of the
two sets.  Constructing the equivalence makes the intended enumeration
explicit and turns the cardinality equality into a one-line application of
\lean{Fintype.card\_congr}.  The pattern is, in a sense, the same kind of
move combinatorialists make when they search for a combinatorial
interpretation of an algebraic identity: the equivalence is the
combinatorial interpretation, and the formalization makes its presence
mandatory rather than optional.


\section{Case Studies}
\label{sec:results}

This section presents the verified results.  Each case study is annotated
with which layer of Sections~\ref{sec:abstract}--\ref{sec:bridge} does the
work, so the section doubles as a worked example of the report.  All proof
paths below are machine-checked in Lean~4 and contain no \lean{sorry}
placeholders.  Some commented auxiliary generalization lemmas outside those
proof paths remain unfinished (see Section~\ref{sec:obstacles}); we flag the
one in-progress case study (a $K_4$-free $P_4$ density bound) at the end of
Section~\ref{sec:unsure}.

\subsection{Scale of the Formalization}

Table~\ref{tab:metrics} summarizes the scale of the active formalization
(excluding archived experiments).  The line counts were computed from the
current repository rather than copied from an external report.

\begin{table}
  \small
  \caption{Scale of the Lean~4 formalization.}
  \label{tab:metrics}
  \begin{tabular}{lrl}
    \toprule
    Component & Lines & Primary files \\
    \midrule
    Specification layer (algebra, density, limits) & 15\,896 &
      \lean{GraphAlgebra/}, \lean{FlagAlgebra/}, \lean{Forbid/}, \lean{Turan/} \\
    Reflection and flag-loading layer & 4\,361 &
      \lean{FlagAlgebra/Compute/}, \lean{Flags/} \\
    Case studies, SDP certificates, API, tactics & 3\,483 &
      \lean{MantelTheorem/}, \lean{ErdosPentagon/}, \lean{Logic/}, \lean{API/} \\
    Shared utilities & 2\,451 &
      \lean{Utils/} \\
    \midrule
    Total active Lean (excluding \lean{Archive/}) & 26\,191 & \\
    \bottomrule
  \end{tabular}

  \medskip
  \begin{tabular}{lr}
    \toprule
    Computationally verified artifact & Count \\
    \midrule
    Auto-generated density lemmas (pentagon proof) & 2\,847 \\
    Multiplication table entries verified & $>$\,300 \\
    \lean{decide~+kernel} SDP matrix equalities (pentagon) &
      3 (sizes $8{\times}8$, $6{\times}6$, $5{\times}5$ over $\Q$) \\
    \lean{native\_decide} density table checks & $>$\,2\,800 \\
    \bottomrule
  \end{tabular}
\end{table}

\subsection{Mantel's Theorem}

Mantel's theorem states that an induced-triangle-free graph on $n$ vertices has at
most $\lfloor n^2/4 \rfloor$ edges, achieved by the complete bipartite graph
$K_{\lfloor n/2\rfloor, \lceil n/2\rceil}$.  The repository contains two
checked statements.  The first is the flag-algebra inequality used as the
upper-bound certificate:
\begin{lstlisting}
theorem Mantel_theorem : K2 <= (1/2 : R) * 1 + K3 := ...
\end{lstlisting}
Here $K_2$ and $K_3$ are elements of \lean{FlagAlgebra} $\emptyt$ representing
the single edge and the triangle.  The inequality asserts that in every
sequence of graphs whose induced triangle density tends to zero, the edge
density is at most $1/2$.  The second is the Tur\'an-density theorem, proved
in the same module by combining the algebraic upper bound with a formal
complete-bipartite lower-bound construction:
\begin{lstlisting}
theorem Turan_density_K3 : turanDensity (completeGraph (Fin 3)) = 1 / 2 := ...
\end{lstlisting}

\paragraph{Which layer does what.}
The proof exercises every layer.  The specification layer
(Section~\ref{sec:abstract}) supplies the algebraic statement and the
transfer theorem.  The forbid-conditioned generators (Section~\ref{sec:flagmatic-gen})
produce and discharge the density-table and multiplication lemmas used to expand $K_2$
and $(a-b)^2$ at size $3$.  The proof exhibits an explicit square in the
algebra of 1-vertex-typed flags.  Let $a = \lean{FlagAlgebra\_2\_1\_0\_0}$
and $b = \lean{FlagAlgebra\_2\_1\_0\_1}$ be the two typed flags of size 2
with a 1-vertex type.  Then
\[
  \bigl\llbracket (a - b)^2 \bigr\rrbracket_1
  \;=\;
  2 \cdot K_2 - \mathbf{1}
  \;\geq\; 0
  \;\;\text{in the semantic non-negativity cone,}
\]
which suffices to conclude $K_2 \leq \frac{1}{2} \cdot \mathbf{1}$ in the
presence of $K_3 = 0$.  The SDP certificate here is small enough to write
explicitly (no external solver needed), making Mantel's theorem a clean
end-to-end illustration of the method before the heavier machinery of the
pentagon proof.

The formal proof in \lean{LeanFlagAlgebras.MantelTheorem.MantelTheorem} uses
the generated expansion lemma \lean{expand\_K2\_on\_three\_vertex\_graphs}
and the squared downward lemma for $(a-b)^2$.  Those supporting lemmas are
proved in \lean{MantelTheorem/Lemmas.lean} and
\lean{MantelTheorem/FlagMul.lean} by the Mantel-specific tactics
\lean{prove\_flag\_expand}, \lean{prove\_flag\_expand\_with\_restriction},
and \lean{prove\_flag\_mul}.  The more uniform API-style version
\lean{Mantel\_flagAlgebra\_API} in \lean{API/MantelTheoremAPI.lean} proves
\[
  \lean{FlagAlgebra\_2\_0\_0\_1 <=[K3.toFinFlag] (1/2) * 1}
\]
using the generic API tactics: \lean{forbidLE\_add\_QuadraticForm} for the
PSD certificate, \lean{reduce\_downward\_flagmul} for the multiplication
rewrites, \lean{expand\_one\_at 3} for the unit-element expansion,
\lean{ac\_sort\_rhs\_pipeline} for the linear normalization, and
\lean{flag\_nonneg} for the closing step.  The API version is roughly thirty
lines including the matrix definitions.

\paragraph{Additional Goodman-style consequences.}
The active \lean{MantelTheorem/} directory also contains two small
flag-algebraic inequalities proved with the same basic infrastructure:
\lean{Goodman\_bound\_on\_triangle\_density} proves the algebraic lower bound
$K_3 \geq K_2(2K_2-\mathbf{1})$, and
\lean{Goodman\_theorem\_on\_Ramsey\_multiplicity} proves
$O_3 + K_3 \geq \frac14\mathbf{1}$, a Ramsey-multiplicity style statement.
These results are not needed for the pentagon proof, but they demonstrate
that the formalized algebraic layer is not specialized solely to the two
headline theorems.

\subsection{A $C_4$-Density Bound for Induced-Triangle-Free Graphs}

A second small case study, in \lean{API/C4TuranAPI.lean}, proves an
inequality about $C_4$ density (\lean{FlagAlgebra\_4\_0\_0\_8}) in
induced-triangle-free graphs:
\begin{lstlisting}
theorem C4_flagAlgebra_API : FlagAlgebra_4_0_0_8 <=[K3.toFinFlag] (3 / 8 : R) * 1 := ...
\end{lstlisting}
The certificate is a sum of two PSD quadratic forms, one for each of two
size-2 types.  The matrices $M_1$ and $M_2$ (sizes $4 \times 4$ and
$3 \times 3$ over $\Q$) and their $LDL^\top$ decompositions are written
explicitly in the same file.  The full proof is approximately twenty-five
lines, almost all of which is identical in structure to the
\lean{Mantel\_flagAlgebra\_API} proof: assemble the PSD certificates,
iterate \lean{forbidLE\_add\_QuadraticForm}, \lean{reduce\_downward\_flagmul},
\lean{expand\_one\_at 4}, \lean{ac\_sort\_rhs\_pipeline}, and
\lean{flag\_nonneg}.  This case study is the smallest non-trivial test of
the API layer: it uses two types rather than one and exercises the
multiplication-table machinery at size $4$, but does not require an
externally generated SDP certificate.

\subsection{The Erd\H{o}s Pentagon Theorem}

The main result is:
\begin{lstlisting}
theorem ErdosPentagon_Turan : generalizedTuranDensity K3 C5 = 24/625 := ...
\end{lstlisting}
proved as the conjunction of an upper bound and a lower bound.

\paragraph{Upper bound.}
\begin{lstlisting}
theorem ErdosPentagon_Turan_upperBound
    : generalizedTuranDensity K3 C5 <= 24/625 :=
  generalizedTuranDensity_le_of_forbidLE (by norm_num)
    ErdosPentagon_flagAlgebra
\end{lstlisting}
The key lemma \lean{ErdosPentagon\_flagAlgebra} establishes the flag-algebra
inequality $C_5 \leq \frac{24}{625} \cdot \mathbf{1}$, i.e., that
$C_5$-density is at most $24/625$ under the $K_3$-free assumption.  The proof
exhibits three PSD matrices $P, Q, R$ over $\mathbb{Q}$ (one per 1-vertex
type, of sizes $8\times 8$, $6\times 6$, and $5\times 5$) and shows that the
sum of the corresponding quadratic forms, after downward-averaging, equals
$\frac{24}{625} \cdot \mathbf{1} - [C_5]$.

\emph{Which layer does what.}  The PSD verification (\lean{P\_posSemidef},
\lean{Q\_posSemidef}, \lean{R\_posSemidef}) lives in
Section~\ref{sec:flagmatic-psd}: $LDL^\top$ decomposition over $\Q$, with the matrix
equalities discharged by \lean{decide +kernel}.  The expansion and
multiplication identities are discharged by the tactic layer
(\lean{prove\_flag\_expand\_with\_forbidden\_flag 5} and the corresponding
multiplication tactic) on top of the elaboration-time-generated density and
multiplication tables.  The linear normalization is done by
\lean{ac\_sort\_pipeline}.  The final assembly into a flag-algebra inequality
is done in \lean{ErdosPentagon/ErdosPentagon.lean}; the API-style restatement
\lean{ErdosPentagon\_flagAlgebra\_API} in \lean{API/ErdosPentagonAPI.lean}
shows the same proof in the API style of the previous two case studies.

The final one-line proof \lean{ErdosPentagon\_Turan\_upperBound} connects
\lean{ErdosPentagon\_flagAlgebra} to the combinatorial statement via
\lean{generalizedTuranDensity\_le\_of\_forbidLE}.

\paragraph{Lower bound.}
The lower bound
\begin{lstlisting}
theorem ErdosPentagon_Turan_lowerBound : generalizedTuranDensity K3 C5 >= 24/625 := ...
\end{lstlisting}
is proved by an explicit construction in
\lean{LeanFlagAlgebras.ErdosPentagon.ErdosPentagon}.  The blow-up
construction is:
\begin{lstlisting}
def blowUp
    {V : Type} [Fintype V] (G : SimpleGraph V) (n : Nat)
    : SimpleGraph (Prod V (Fin n)) :=
  { Adj v w := G.Adj v.1 w.1
    symm v w := by apply G.symm }
\end{lstlisting}
The $n$-fold blow-up of $C_5$ (replacing each vertex by an independent set
of $n$ vertices and each edge by a complete bipartite graph) is
induced-triangle-free.  Since \(K_3\) is complete, this is the same condition
as the Lean predicate appearing below: any triangle subgraph is automatically
an induced triangle.  The repository proves the statement by transporting any
hypothetical triangle in the blow-up back to a triangle in the base graph:
\begin{lstlisting}
theorem blowUp_K3_free
    {m : Nat} {G : SimpleGraph (Fin m)}
    (n : Nat) (hfree : K3.Free G) :
    K3.Free (blowUp G n) := ...
\end{lstlisting}
The construction has $5n$ vertices and at least $n^5$ copies of $C_5$:
\begin{lstlisting}
lemma subgraphCount_blowUp_C5_ge
    (n : Nat) :
    subgraphCount C5 (blowUp C5 n) >= n ^ 5
\end{lstlisting}
The proof of this counting lemma constructs an injection
\lean{(Fin 5 -> Fin n) -> (blowUp C5 n).Subgraph}: choosing one vertex in
each of the five parts determines a copy of $C_5$, and distinct choices give
distinct subgraphs.  After transferring the blow-up to the vertex type
\lean{Fin (5 * n)}, the repository obtains the finite extremal lower bound
\lean{generalizedExtremalNumber\_K3\_C5\_div\_choose\_ge}.  The asymptotic
arithmetic is the elementary estimate
\[
  \frac{n^5}{\binom{5n}{5}}
  \;\geq\;
  \frac{n^5}{(5n)^5 / 5!}
  \;=\;
  \frac{24}{625},
\]
combined with the convergence theorem \lean{tendsto\_generalizedTuranDensity}
in \lean{Turan/GeneralizedTuran.lean}.  Mathematically, the lower bound uses
$\lim_{n\to\infty} n^5/\binom{5n}{5} = 24/625$; in the Lean proof the limit
is replaced by the displayed lower bound on the subsequence, plus general
convergence machinery.

\paragraph{Reading the trust profile of this proof.}
The proof of \lean{ErdosPentagon\_Turan} depends on two classes of
computational verification.  \lean{native\_decide} is used for all density
table equalities; it relies on the correctness of the Lean-to-native
compiler, which is outside the kernel.  \lean{decide +kernel} is used for the
three SDP matrix equalities; it uses no compiled native code, so it
introduces no trust beyond the standard Lean axioms (\lean{propext},
\lean{Quot.sound}, \lean{Classical.choice}).  As we noted in
Section~\ref{sec:flagmatic-trust}, the choice of \lean{native\_decide} for the density
tables was forced by feasibility, while the choice of \lean{decide +kernel}
for the SDP equalities was deliberate: the SDP step is the trust-critical
component, and the small matrices allow kernel evaluation.  There are no
\lean{sorry}s or new \lean{axiom}s in the proof paths of either
\lean{Mantel\_theorem} or \lean{ErdosPentagon\_Turan}.


\section{Engineering Obstacles}
\label{sec:obstacles}

The specification layer of Section~\ref{sec:abstract} exposed several
engineering obstacles that we report here as obstacles rather than design
principles.  Three of them follow from Lean~4's intensional type theory and
its typeclass mechanism: quotient eliminators carry recurring compatibility
proofs; size-changing combinatorial operations produce dependent-type
transports; and \lean{Fintype} instances synthesized through different
elaboration paths are extensionally but not definitionally equal.  A fourth
obstacle, function extensionality, is mathematically benign but appears at
every step that relates a function-level equality to a pointwise one.  We
describe the obstacles, the mitigations we used, and where the open ends
remain.  The point is not the particular workarounds; it is that an
intensional type theory makes these obstacles visible at the boundaries of
quotient, dependent-type, and finite-enumeration constructions, and someone
formalizing similar material should expect to face them.

\subsection{Quotient Eliminators and Their Recurring Cost}

Both \lean{Flag} and \lean{FlagAlgebra} are quotients, so any function
\emph{from} them or equality \emph{in} them must go through
\lean{Quotient.lift} or \lean{Quotient.sound}.  \lean{Quotient.lift} requires
a proof that the function respects the equivalence relation;
\lean{Quotient.sound} requires producing a witness of the relation for each
equality goal.

The adequacy theorems in Section~\ref{sec:reflection} illustrate the cost.
To state that \lean{Sym2Graph.toFlag} is inverse to
\lean{Flag.toSym2EmptyTypedFlag}, one cannot simply compute both sides to the
same normal form; the map out of the quotient must go through an eliminator:
\begin{lstlisting}
def Sym2EmptyTypedFlag.toFlag (F : Sym2EmptyTypedFlag n) : Flag (*@$\emptyt$@*) (Fin n) :=
  Quotient.lift Sym2Graph.toFlag Sym2Graph.toFlag_respect_eqv F
\end{lstlisting}
The \lean{toFlag\_respect\_eqv} proof---showing that isomorphic
\lean{Sym2Graph} values map to the same \lean{Flag}---is not automatic; it
required \lean{Quotient.sound} applied to a graph isomorphism witness
assembled from the concrete data.  This is exactly the kind of compatibility
proof that ordinary mathematical notation suppresses.

The recurring cost of carrying \lean{Quotient} eliminators through the
development is real.  A design built around canonical representatives from the
start would not pay this cost, but it would attach representation-specific
side conditions to every algebraic lemma; we chose the quotient design to keep
the specification close to Razborov's mathematics and accept the eliminator
overhead.  Whether the alternative would have been more economical overall is
a tradeoff we cannot evaluate from this vantage point.

\subsection{Dependent-Type Transports from Size-Changing Operations}

Flag algebra operations change the vertex count of flags: the product of a
size-$m_1$ and a size-$m_2$ flag lives at size $\ell \geq m_1 + m_2 - k$.  In
the formalization, the product of basis elements is computed as a sum over
all size-$\ell$ flags.  To show this is well-defined on the quotient (i.e.,
independent of $\ell$), one proves \lean{flagMulWithSize\_indep\_on\_size}.
Its hypotheses say that both output sizes are large enough, and its proof has
to relate sums over \lean{Flag sigma (Fin ell)} and
\lean{Flag sigma (Fin ell')}, which are definitionally distinct types even
when the corresponding mathematical sets are connected by expansion.

The same issue arises internally when collecting the flags for the
multiplication sum.  A \lean{FlagList sigma t Vl} (a $t$-tuple of flags with
type family $Vl$) changes type when a new flag is inserted: the result has
type \lean{FlagList sigma (t+1) (listTypeInsert Vl W)}.  Because
\lean{listTypeInsert Vl W} is not definitionally equal to any pre-existing
type family, relating the old and new lists requires \emph{heterogeneous
equality} (\lean{HEq}):
\begin{lstlisting}
theorem flagList_HEq
    (h_Vl_eq : Vl' = Vl)
    (h_Fl_eq : forall i, Fl i = cast (Flag.type_eq h_Vl_eq i) (Fl' i))
    : HEq Fl Fl'
\end{lstlisting}
The \lean{cast} calls are explicit coercions through propositional equality
proofs; they are the formalization's acknowledgment that two types are equal
only up to a proof, not by computation.

Similarly, even if two indices $i$ and $i'$ into a flag list are
propositionally equal ($i = i'$), the types \lean{Flag sigma (Vl i)} and
\lean{Flag sigma (Vl i')} are not definitionally equal, forcing the use of
heterogeneous equality:
\begin{lstlisting}
theorem flaglist_heq_of_idx_eq {i i' : Fin t} (h : i = i')
    : HEq (Fl i) (Fl i') := by subst h; rfl
\end{lstlisting}
This pervasive use of \lean{HEq} and \lean{cast} propagates through the entire
density computation infrastructure: every operation that crosses a size
boundary must carry explicit propositional equality evidence.  Our
\emph{mitigation} was to isolate each size-change obligation in a small named
lemma---examples are \lean{flagList\_HEq},
\lean{flaglist\_heq\_of\_idx\_eq}, and \lean{flagListDensity\_HEq\_eq}---so
that high-level proofs invoke clean invariance statements while the transport
details remain in the small layer where they belong.

\subsection{Function Extensionality and Proof-Valued Fields}

In Lean~4's intensional type theory, function extensionality
($(\forall x,\, f\,x = g\,x) \Rightarrow f = g$) is not definitional but is
available as an axiom (\lean{funext}, derived from \lean{propext}).  This
creates friction wherever the formalization must relate a function-level
equality (such as a type embedding) to a pointwise equality (the embedding's
behavior on individual vertices).  For example, constructing a labeled-graph
isomorphism from a vertex map $\zeta$ requires reassembling pointwise
information back into a function:
\begin{lstlisting}
have h_emb : forall t : T, zeta (G0'.type_embed t) = G1'.type_embed t := ...
let iso : G0'.coe =~f G1'.coe := { ..., type_preserve := funext h_emb }
\end{lstlisting}
Without \lean{funext} as an axiom, the \lean{type\_preserve} field could not
be filled.

Proof-valued structure fields create a related difficulty.  The
\lean{Sym2Graph} type carries a field
\lean{edges\_valid : forall e in edges, not e.IsDiag}; two \lean{Sym2Graph}
values with the same edge set but different proof terms for
\lean{edges\_valid} are propositionally equal (by proof irrelevance) but not
definitionally equal.  Comparing \lean{Sym2Graph} values in the adequacy
theorems therefore requires \lean{proof\_irrel\_heq} at every such field:
\begin{lstlisting}
theorem Sym2Graph.toLabeledGraph_toSym2Graph_eq (G : Sym2Graph n) :
    G.toLabeledGraph.toSym2Graph = G := by
  congr  -- reduces to field equalities
  ...
  exact proof_irrel_heq _ _  -- edge validity field
\end{lstlisting}

\subsection{Fintype Instance Conflicts}
\label{sec:fintype}

The fourth obstacle is the most pervasive in the counting layer.  Sets defined
by set-builder notation $\{x \mid P\ x\}$ have type \lean{Set T}, which is
\lean{Finite} whenever \lean{T} is a \lean{Fintype}, but does not
automatically carry a \lean{Fintype} instance.  Many Mathlib lemmas about
cardinality (\lean{Fintype.card}, \lean{Finset.card\_univ},
\lean{Fintype.card\_congr}) require \lean{Fintype}, not merely \lean{Finite}.
The elaborator infers \lean{Fintype} instances automatically, but in a
development with multiple overlapping constructions (quotient types, embedded
subgraphs, partial-function spaces), the same type can receive different
\lean{Fintype} instances from different elaboration paths.  When two
instances do not reduce to the same term definitionally, the kernel rejects
the goal even if the instances are provably equal.

The mitigation is to bind the intended \lean{Fintype} instance immediately
after naming a finite set:
\begin{lstlisting}
let hS0 : Fintype S0 := Fintype.ofFinite S0
let hS1 : Fintype S1 := Fintype.ofFinite S1
have card_eq : Fintype.card S0 = Fintype.card S1 :=
  Fintype.card_congr h_iso_S0_S1
\end{lstlisting}
This pattern appears repeatedly in \lean{SubflagListDensity.lean}: for the
quotient-respect lemma for list densities, for permutation of sampled flags,
and for insertion of an empty flag.  Without the explicit \lean{let} bindings,
Lean can synthesize a \lean{Fintype} for each set through a different path at
each use, and \lean{Fintype.card\_congr} no longer sees the two cardinalities
as referring to the same enumeration terms.  A similar explicit
\lean{Fintype.ofFinite} construction is used in \lean{FlagOperators.lean} for
the isomorphism set with a fixed underlying graph.

For parameterized types such as injections \lean{Embedding V W} and
equivalences \lean{Equiv V W}, Lean can synthesize \lean{Fintype} instances in
multiple ways depending on which prior instances are in scope.  In
\lean{Compute/Basic.lean}, where \lean{Fintype} instances for these types are
defined, the elaborator had to be given explicit guidance via
\lean{Finset.univ} with a type ascription:
\begin{lstlisting}
instance : Fintype (Embedding V W) :=
  { elems := (Finset.univ : Finset (V -> W)).filterMap ..., ... }

instance : Fintype (G1 ->g G2) :=
  { elems := (Finset.univ : Finset (Embedding V W)).filterMap ..., ... }
\end{lstlisting}
The graph-embedding and graph-isomorphism instances are built by filtering
finite enumerations of simpler function or equivalence types.  The explicit
\lean{Finset.univ} type ascription keeps the intended enumeration visible to the elaborator,
instead of relying on typeclass search to rediscover it in a later goal.
A related issue arose for dependent function types indexed by a small finite
type, where Lean's typeclass search does not look through \lean{match}
expressions inside type-valued functions; the \lean{FintypeList} typeclass
threads \lean{Fintype} instances per component explicitly to compensate.

\paragraph{What remains open.}
Not all intensionality obstacles were resolved.  The most significant gap is
in \lean{Forbid/Basic.lean}: several strengthening lemmas about recovering
representative-level information from quotient-level forbidden equalities
remain as commented proof sketches, not as imported declarations or trusted
axioms.  For example, one would like a general theorem saying that if a
basis vector is zero under the forbidden equality relation, then the
corresponding unlabeled flag has positive induced density of the forbidden
graph.
Proving such a statement requires careful reasoning about how zero-space
relations interact with individual representatives.  This is not needed for
the main case studies---\lean{Mantel\_theorem} and \lean{ErdosPentagon\_Turan}
are proved without those sketches---but it marks a natural direction for
making the \lean{forbidLE} framework more convenient for future applications.



\section{What We Are Unsure About}
\label{sec:unsure}

This is the section that distinguishes this paper from a single-thesis
formalization paper.  Several questions surfaced during the work; we do not
resolve them here, and we believe that being explicit about them is more
honest than implying that our final design settled them.

\paragraph{Predicate-deferral vs.\ universal-theory specification.}
The forbidden-subgraph predicate \lean{forbidLE}
(Section~\ref{sec:forbidden}) defers the choice of forbidden graph from the
universal theory at the top of the algebra to a parameter on the semantic
order at use sites.  Razborov's original presentation does the opposite.
Both designs can express the same theorems.  We do not know which design is
preferable in the long run.  Reusable lemmas about the $\emptyt$-typed algebra
are easier to share across forbidden graphs in our design; lemmas that
exploit a specific forbidden graph from the outset may be cleaner in
Razborov's.  We invite future work to compare the two on a development that
needs both styles.

\paragraph{Are the intensionality obstacles inherent or accidental?}
The mitigations of Section~\ref{sec:obstacles}---explicit \lean{Fintype}
bindings, \lean{HEq}-and-\lean{cast} discipline for size-changing operations,
proof-irrelevance bridges for \lean{Sym2Graph} fields---are responses to
specific Lean~4 facts: typeclass coherence by convention rather than by
enforcement, definitional inequality for dependently-typed families that are
propositionally equal, decidable definitional equality that does not include
function extensionality.  We do not know which of these would survive a port
to a different ITP (Rocq with SSReflect, Agda, Isabelle/HOL with a different
typeclass story).  In particular, the \lean{Fintype} ambiguity issue and the
\lean{HEq} discipline are tied to specific elaboration heuristics; a
proof assistant with stricter coherence or with proof-irrelevant elaboration
might dissolve some of these obstacles.  We report them as Lean~4 obstacles
without claiming they are foundational.

\paragraph{Is the bridge layer essential or an artifact of our spec choices?}
Section~\ref{sec:bridge}'s reflection, tactic, and API layers are
load-bearing in our development: the pentagon proof does not go through
without them.  But a different specification layer might collapse some of
this work into the spec side.  For instance, a spec layer that committed to
canonical representatives in the quotient (rather than carrying isomorphism
classes as quotients) would not need the same adequacy theorems, though it
would attach representation-specific side conditions to every algebraic
lemma.  Whether the bridge layer is essential to the flag-algebra method or
a consequence of our specific design choices is an open question that
deserves a different formalization to test.

\paragraph{Does the elaboration-time generation pattern scale beyond
pentagon size?}
The pentagon proof generates roughly 2,847 density lemmas at elaboration
time.  This is large but well within Lean's elaborator budget on contemporary
hardware.  Larger flag-algebra results in the literature, particularly those
that work with $7$-vertex flags or beyond, would generate far more.  We have
not stress-tested the elaboration-time generation pattern at that scale and
do not know which part of the pipeline---density-table loading,
multiplication-table loading, or the named-lemma cache itself---would be the
first bottleneck.

\paragraph{Work in progress: $K_4$-free $P_4$ density.}
The file \lean{API/K4freeP4.lean} contains a skeleton for the
inequality
\[
  \lean{P4\_density <=[K4.toFinFlag] (96/27) * 1}
\]
as a further test of the API layer.  The proof structure compiles---it uses
the same API tactics as Mantel, $C_4$, and the pentagon---but several
non-negativity facts (\lean{f}$_1$\lean{\_nonneg}, the main combined
inequality) currently remain as \lean{sorry} placeholders.  We mention it
because it demonstrates that the API generalizes beyond induced-triangle-free
forbids ($K_3$) to a $4$-clique forbid ($K_4$); we do \emph{not} count it
among the proved results.


\section{Lessons for the Next Person}
\label{sec:lessons}

This section reports lessons we found transferable beyond flag algebras.
Each is defended by a specific artifact in the formalization.  None is
presented as a methodological discovery.  The first four are concrete
engineering lessons.  The fifth is a meta-level lesson whose generality is
conjectural---we mark it as such in the body.  The sixth is a reporting
lesson about how to read the development's trust profile; it is included
because the trust profile is otherwise easy to misread.

\paragraph{Lesson 1 (Quotient objects are best left as quotients in the
specification).}
Flag algebras are naturally quotient objects: graphs up to isomorphism, and
flag-algebra elements modulo the equational theory of the algebra.  A
formalization could pick canonical representatives in the specification
layer to make every operation computable.  We do not.  The specification
keeps the quotient, and the executable layer in \texttt{FlagAlgebra/Compute/}
chooses representatives over the concrete \lean{Sym2Graph} type.  The two
layers are connected by adequacy theorems such as
\lean{flagDensity1\_eq\_sym2FlagDensity1} and
\lean{flagDensity2\_eq\_sym2FlagDensity2}, whose proofs factor through
count-level adequacy
(\lean{labeledSubgraphListCount\_eq\_sym2InducedSubgraphListCount}).  The
benefit is that quotient and subgraph reasoning is paid once, inside
adequacy, rather than thousands of times in the generated density lemmas.
The cost is the adequacy proof itself, which is non-trivial.  The
transferable claim is narrow: when the textbook object is a quotient and
computability comes from a refinement, keeping the quotient in the
specification localizes the awkward proof obligations rather than spreading
them across the development.

\paragraph{Lesson 2 (Computable is not the same as usable).}
The naive reflective design is to leave the abstract quotient definitions in
place, give the relevant finite types \lean{Fintype} instances, and ask
\lean{native\_decide} to evaluate density goals by unfolding
\lean{flagDensity}.  This type-checks.  It is not a viable compilation
strategy for the pentagon proof: each density check repeatedly constructs
finite sets of abstract subgraphs, compares quotient classes by searching
for flag isomorphisms, synthesizes overlapping \lean{Fintype} instances, and
carries proof-valued fields through the evaluator.  The graphs are small but
the proof terms are not---the evaluator is reducing a proof-relevant
encoding of finite combinatorics, not a graph algorithm.

What worked was to treat reflection as verified compilation.  The abstract
definitions remain the specification; the executable code in
\texttt{Compute/FastIso.lean} and \texttt{Compute/FlagDensity.lean} is a
compilation target with deliberate optimizations (edge-count pruning before
permutation enumeration; type-aware permutation that fixes labeled vertices
and enumerates only the remainder).  Adequacy theorems certify that the
compilation preserves the specification's meaning.  The generalization is:
when reflection on the natural specification is computable in principle but
stuck in practice, the right move is not to make the specification more
computable but to separate spec from implementation and connect them once.

\paragraph{Lesson 3 (Names can be elaboration-time metadata when type-class
search is too slow).}
Every flag and flag-algebra constant in the development is named according
to the scheme \lean{Flag\_n\_k\_m\_i} and \lean{FlagAlgebra\_n\_k\_m\_i},
where $n$ is the vertex count, $k$ and $m$ identify the type graph, and $i$
is a canonical index.  The convention is not organizational.  Tactics parse
trailing integers from a constant's \lean{Name}, which is a syntactic object
the elaborator can inspect without type-checking or attribute resolution.
\lean{ac\_sort\_pipeline} computes sort keys entirely in the fast
elaboration path.

The recorded counterexample is in \texttt{ErdosPentagon/Lemmas.lean}: a
commented-out call reads \texttt{-- sort\_at -- takes more than 10 minutes}.
The generic \lean{sort\_at}, which uses \lean{simp} without index guidance,
times out on expressions of around 25 flag terms.  The index-guided
\lean{ac\_sort} completes the same step in under a second by exploiting the
name encoding to bypass combinatorial search.  We do not claim that names
beat user attributes or a custom AST in general; on this development they
did, and the design axis is worth naming so that future formalizations
facing the same large-expression performance wall know to consider it.

\paragraph{Lesson 4 (External computation is a candidate generator, not a
proof).}
Two pieces of unverified external machinery feed the pentagon proof: the
JSON density tables produced by an enumeration script, and the rational PSD
certificates produced by an SDP solver.  Neither is trusted.  The density
tables drive the elaboration macros \lean{load\_flag\_pair\_density\_theorems}
and \lean{MulLoader}: for each table entry the macros generate the
statement Lean should prove (e.g., \lean{flagDensity2 F0 F1 G = 3/10}) and
prove it by unfolding the concrete flag definitions, rewriting with the
adequacy theorem, and discharging the resulting decidable proposition.

A wrong table entry does not corrupt the proof; it fails as a theorem at the
evaluator step, because after adequacy rewriting the goal is a decidable
equality between two rationals computed independently.  The SDP certificate
is handled analogously: the matrix \lean{P} and its claimed
\lean{LDL$^\top$} decomposition are written explicitly in
\lean{ErdosPentagon.MatrixDef}, and the equality \lean{P = LP * Matrix.diagonal dP * LP.transpose}
is verified by Lean rather than asserted.  The transferable framing is: the
external tool proposes values; the elaboration layer generates the
corresponding theorem statements; the verified evaluator accepts only those
it can reproduce.  JSON tables and solver output are test vectors against
the spec, not certificates.

\paragraph{Lesson 5 (When a constraint is already expressible in the
textbook framework, the formalization choice is the timing of commitment).}
\emph{This lesson is meta-level; its generality is conjectural.}
Razborov's framework specifies which graphs are forbidden via a universal
first-order theory chosen at the outset.  Our design defers the same
specification to a parameter on the semantic order: the forbidden-subgraph
assumption appears as a predicate we instantiate only when proving a
specific theorem.  Both designs can express the same mathematical
statements; what differs is which intermediate lemmas are reusable across
choices of constraint.

We do not claim our deferred design is preferable to Razborov's.  We do
report that, in our formalization, deferring the specification has the
consequence that random-extension theorems became the definitional content
of the forbidden-subgraph relation rather than analytic ingredients inside
a fixed theory.  The general lesson, if it generalizes, is that when a
constraint is already expressible in the textbook framework, the
formalization design question is when in the development to commit to a
specific instance of it, and that this question deserves explicit attention
rather than being absorbed into the choice of framework.

\paragraph{Lesson 6 (Honest description of the trust profile).}
This is a reporting lesson, not a methodological one.  Our formalization
uses two evaluators with different trusted-computing-base profiles:
\lean{native\_decide}, which compiles the goal to native code and evaluates
it outside the kernel; and \lean{decide +kernel}, which evaluates inside the
kernel.  They are not interchangeable.

The roughly 2,800 density-table lemmas are discharged by
\lean{native\_decide}.  This was not a deliberate trust choice:
\lean{decide +kernel} does not let our scripts type-check on the
proof-relevant encoding the specification uses, even after the optimizations
of Section~\ref{sec:reflection}.  The three SDP matrix equalities (sizes
$8 \times 8$, $6 \times 6$, $5 \times 5$ over $\mathbb{Q}$) are discharged
by \lean{decide +kernel}.  Here the choice is deliberate: the matrices are
small enough that kernel evaluation finishes, and the upper bound proof
would collapse if a non-PSD matrix were erroneously accepted, so we pay the
kernel-evaluation cost to keep PSD verification inside the kernel.

We report this asymmetry rather than presenting it as a design pattern.  A
reader who depends on the artifact may want to know that the load-bearing
PSD step is kernel-checked and that the bulk of the density bookkeeping is
checked by compiled code; the profile is information they can act on.  A
reader who finds the profile insufficient---for example, who would want
\lean{decide +kernel} on the density tables as well---can read this as a
record of where the current development sits, not as a claim that the
current placement is principled.  The trust profile of a large formalization
is not always chosen, but it is always reportable.


\section{Related Work}
\label{sec:related}

\paragraph{Formalizations of combinatorics.}
Several combinatorial theorems have been formalized in Lean and related
systems.  Dillies and Mehta~\cite{dillies2022szemeredi} formalized
Szemer\'edi's Regularity Lemma in Lean~4 using a graph-theoretic presentation;
our work faces similar challenges with large Mathlib dependencies but at a
different point in the proof-engineering spectrum (our main difficulty is
data-heavy SDP verification rather than epsilon-delta regularity arguments).
Mehta~\cite{mehta2022kruskal} formalized the Kruskal-Katona theorem.
Subercaseaux et al.~\cite{subercaseaux2024hexagon} formally verified the
empty hexagon number using a SAT-based pipeline.  Our work differs in that
the central method (flag algebras) relies on an algebraic quotient
construction and measure-theoretic limits, rather than purely combinatorial
enumeration or SAT solving.  Gowers et al.~\cite{gowers2024formalizing}
recently announced a Lean~4 formalization of the polynomial Freiman-Ruzsa
conjecture proof; their experience with large-scale Lean formalization of
recent combinatorics results is closely related to ours.

\paragraph{Proof by reflection in proof assistants.}
Proof by reflection is a classical technique in the Coq community (see
Chlipala~\cite{chlipala2013cpdt} for a survey), and has been used in Lean
for verified computation (e.g., in Mathlib's \lean{decide} and
\lean{norm\_num} tactics).  Cohen et al.~\cite{cohen2013refinements} develop
a systematic methodology for building efficiently computable representations
alongside their abstract counterparts in Coq; our \lean{Sym2Graph} /
\lean{SimpleGraph} pair follows the same philosophy adapted to Lean~4's
type-class and instance-search mechanism.  Our use of elaboration-time
theorem generation from JSON data (Section~\ref{sec:reflection}) is related to
the approach of Amin and Rompf~\cite{amin2017collapsing}, who generate
verified code from high-level specifications; we generate verified density
lemmas from externally computed tables.

\paragraph{Computer-assisted flag algebra arguments.}
The Flagmatic software~\cite{vaughan2013flagmatic} automates the computation
of flag algebra certificates as a Python tool, but produces results that
are validated only informally.  To our knowledge, our work provides the
first formally verified end-to-end pipeline from flag algebra certificates
to combinatorial theorems.  Razborov~\cite{razborov2013flag} gives a
comprehensive survey of flag algebra applications; the $\sim$100 results
catalogued there all involve the same recurring concerns identified in
Section~\ref{sec:intro}: abstract structure, extensional finite counting,
data-heavy computation, and algebraic bookkeeping.  This common shape
suggests our infrastructure is broadly applicable, although
Section~\ref{sec:unsure} flags the scaling question.

\paragraph{SDP certificate verification.}
The verification of SDP certificates has been studied in the context of
sum-of-squares proofs; Harrison~\cite{harrison2007verifying} verifies
positivity certificates for polynomials in HOL Light.  Our LDL$^\top$
approach is structurally similar but operates over $\mathbb{Q}$ (exact
arithmetic) rather than $\mathbb{R}$, which is essential for
kernel-checkable verification.  Parrilo and
Sturmfels~\cite{parrilo2003minimizing} connect algebraic certificates
(sum-of-squares) to SDP duality in the polynomial setting; our setting is
analogous but over the discrete flag algebra rather than the polynomial
ring.

\paragraph{Tactic metaprogramming.}
Custom tactics for domain-specific proof automation have been developed in
many contexts~\cite{ebner2017structured,sozeau2008coq}.  Our tactics exploit
the naming convention as a form of \emph{reflected metadata}: the name of a
constant encodes information that the tactic exploits without having to
reduce the mathematical content of the expression.  This is related to, but
distinct from, approaches based on type-class inference or user-facing
attribute annotations, and is most closely analogous to the ``term-mode
reflection'' pattern described by Malecha and
Morrisett~\cite{malecha2014towards}.

\paragraph{Graph limits and Turán densities.}
The mathematical foundation of flag algebras (convergent graph sequences,
positive homomorphisms, the graphon limit) is developed in depth by
Lov\'asz~\cite{lovasz2012large}.  To our knowledge, no formalization of
graph limit theory exists in any proof assistant; our \lean{FlagSequence.lean}
and \lean{RandomHom.lean} provide a partial formalization sufficient for
the flag algebra method, but a full formalization of the graphon theory
remains an open challenge.


\section{Conclusion}
\label{sec:conclusion}

We have reported on a Lean~4 formalization of Razborov's flag algebra method
for simple graphs.  The development carries a specification of the
mathematical theory---flags as quotients modulo isomorphism, the quotient
algebra of density-expansion relations, positive homomorphisms and the
semantic non-negativity cone, and the predicate-based forbidden-subgraph framework---and a
bridge of reflection, tactic, and API layers that turns specifications
written in this language into proofs Lean's kernel can check.  The case
studies prove the flag-algebra form of Mantel's theorem, the Erd\H{o}s
pentagon theorem ($\tdensity{C_5}{K_3} = 24/625$), two Goodman-style
inequalities, and an inequality bounding the $C_4$-density of induced-triangle-free
graphs.

We do not claim to have found
\emph{the} right way to formalize flag algebras.  We have described the
choices we made for the mathematical specification, the engineering
obstacles those choices exposed, the bridge layer we built around them, the
recurring proof patterns we observed, the questions our work surfaced but
did not resolve, and the lessons we believe are transferable beyond flag
algebras.  In places where an outcome was forced by feasibility rather than
chosen for principled reasons---the use of \lean{native\_decide} for the
density tables being the most prominent example---we have said so.

We expect the architecture to transfer to further graph flag algebra
arguments because the same concerns recur: quotient-level semantics, exact
finite counting, externally generated but internally checked certificates,
and large-scale algebraic bookkeeping.  Whether it transfers to richer
combinatorial structures (hypergraphs, directed graphs, oriented graphs) is
a question we have not tested.  A full formalization of graphon theory,
connecting flag algebra limits to the Lov\'asz theory~\cite{lovasz2012large},
would provide a richer mathematical foundation for future extensions.
Automating the discovery of SDP certificates from within Lean (rather than
importing them from external solvers) remains a longer-term goal.  Most
immediately, replacing \lean{native\_decide} with \lean{decide +kernel}
throughout, or with a formally verified external checker, would eliminate
the remaining dependency on the native compiler---though, as
Section~\ref{sec:lessons} notes, this is a feasibility constraint of the
current encoding, not a logical limitation of the method.


% =========================================================================
\bibliographystyle{ACM-Reference-Format}
\begin{thebibliography}{99}

\bibitem{razborov2007flag}
A.~A.~Razborov.
\newblock Flag algebras.
\newblock \textit{Journal of Symbolic Logic}, 72(4):1239--1282, 2007.

\bibitem{razborov2008}
A.~A.~Razborov.
\newblock On the minimal density of triangles in graphs.
\newblock \textit{Combinatorics, Probability and Computing},
  17(4):603--618, 2008.

\bibitem{razborov2013flag}
A.~A.~Razborov.
\newblock Flag algebras: An interim report.
\newblock In \textit{The Mathematics of Paul Erd\H{o}s II}, pages 207--232.
  Springer, 2013.

\bibitem{grzesik2012}
A.~Grzesik.
\newblock On the maximum number of five-cycles in a triangle-free graph.
\newblock \textit{Journal of Combinatorial Theory, Series B},
  102(5):1061--1066, 2012.

\bibitem{hatami2012}
H.~Hatami, J.~Hladk\'{y}, D.~Kr\'{a}l, S.~Norine, and A.~Razborov.
\newblock On the number of pentagons in triangle-free graphs.
\newblock \textit{Journal of Combinatorial Theory, Series A},
  120(3):722--732, 2013.

\bibitem{dillies2022szemeredi}
Y.~Dillies and B.~Mehta.
\newblock Formalising Szemer\'{e}di's regularity lemma in Lean.
\newblock In \textit{Proc.\ ITP 2022}, LIPIcs~237, 2022.

\bibitem{mehta2022kruskal}
B.~Mehta.
\newblock Formalising the Kruskal-Katona theorem in Lean.
\newblock In \textit{Proc.\ ITP 2022}, 2022.

\bibitem{subercaseaux2024hexagon}
B.~Subercaseaux, M.~J.~H.~Heule, J.~Mackey, J.~Meadows, R.~Tao, and
  C.~Wu.
\newblock Formal verification of the empty hexagon number.
\newblock In \textit{Proc.\ ITP 2024}, LIPIcs~309, 2024.

\bibitem{gowers2024formalizing}
W.~T.~Gowers, D.~Green, F.~Manners, and T.~Tao.
\newblock On a conjecture of Marton.
\newblock \textit{arXiv:2311.05762}, 2023.
\newblock (Lean formalization by T.~Bloom et al., 2024.)

\bibitem{vaughan2013flagmatic}
E.~R.~Vaughan.
\newblock Flagmatic 2.0: A user-friendly implementation of flag algebra
  arguments.
\newblock Preprint, 2013. Available at \texttt{https://github.com/jsloan/flagmatic}.

\bibitem{chlipala2013cpdt}
A.~Chlipala.
\newblock \textit{Certified Programming with Dependent Types}.
\newblock MIT Press, 2013.

\bibitem{cohen2013refinements}
C.~Cohen, M.~Denes, and A.~M\"{o}rtberg.
\newblock Refinements for free!
\newblock In \textit{Proc.\ CPP 2013}, LNCS~8307, pages 147--162, 2013.

\bibitem{amin2017collapsing}
N.~Amin and T.~Rompf.
\newblock Type soundness proofs with definitional interpreters.
\newblock In \textit{Proc.\ POPL 2017}, pages 666--679, 2017.

\bibitem{harrison2007verifying}
J.~Harrison.
\newblock Verifying nonlinear real formulas via sums of squares.
\newblock In \textit{Proc.\ TPHOLs 2007}, LNCS~4732, pages 102--118, 2007.

\bibitem{parrilo2003minimizing}
P.~A.~Parrilo and B.~Sturmfels.
\newblock Minimizing polynomial functions.
\newblock In \textit{Algorithmic and Quantitative Real Algebraic Geometry},
  DIMACS Series in Discrete Math., pages 83--99, 2003.

\bibitem{ebner2017structured}
G.~Ebner, S.~Ullrich, J.~Roesch, J.~Avigad, and L.~de~Moura.
\newblock A metaprogramming framework for formal verification.
\newblock \textit{Proc.\ ACM Program.\ Lang.}, 1(ICFP):34:1--34:29, 2017.

\bibitem{sozeau2008coq}
M.~Sozeau and N.~Oury.
\newblock First-class type classes.
\newblock In \textit{Proc.\ TPHOLs 2008}, LNCS~5170, 2008.

\bibitem{malecha2014towards}
G.~Malecha and G.~Morrisett.
\newblock Towards foundational verification of cyber-physical systems.
\newblock In \textit{Proc.\ SoSyM 2014}, 2014.

\bibitem{billingsley1999convergence}
P.~Billingsley.
\newblock \textit{Convergence of Probability Measures}, 2nd ed.
\newblock Wiley, 1999.

\bibitem{lovasz2012large}
L.~Lov\'{a}sz.
\newblock \textit{Large Networks and Graph Limits}.
\newblock American Mathematical Society, 2012.

\bibitem{lovasz2006limits}
L.~Lov\'{a}sz and B.~Szegedy.
\newblock Limits of dense graph sequences.
\newblock \textit{Journal of Combinatorial Theory, Series~B}, 96(6):933--957,
  2006.

\end{thebibliography}

\end{document}
